% BEGIN EMBEDDED VERIFICATION MATERIALS
\begin{filecontents*}[overwrite]{verify_harman_certificates.py}
"""Verify the Harman-region proof using rational arithmetic only.

No numerical optimizer, symbolic algebra system, or nonstandard dependency
is imported. The floating-point discovery program is not imported.
"""
import json
from fractions import Fraction as F
from pathlib import Path

N=8
names=['m','h','a','b','c','f','z','t']

def row(**kw):
    ans=[F(0)]*N
    for key,value in kw.items():ans[names.index(key)]=F(value)
    return ans

def unpack(item):
    a,b,label=item
    assert len(a)==N
    return [F(x) for x in a],F(b),label

def check_provenance(data):
    summary=data['summary'];region=summary['region']
    lo,hi=F(summary['theta_lower']),F(summary['theta_upper'])
    assert F(1,25)<=lo<=hi<=F(1,24)
    assert region in ['H1','H2']
    allowed={
        'theta_lower':(row(t=-1),-lo),
        'theta_upper':(row(t=1),hi),
        'm_half':(row(m=1),F(1,2)),
        'h_le_m':(row(h=1,m=-1),F(0)),
        'sum_lengths':(row(m=1,h=1),F(1)),
        'a_upper':(row(a=2,m=-1),F(0)),
        'b_upper':(row(b=2,h=-1),F(0)),
        'c_upper':(row(c=2,m=1,h=1),F(1)),
        'f_upper':(row(f=1,t=1),F(1)),
        'z_m_half':(row(z=1,m=1),F(1,2)),
        'z_h_le_m':(row(z=1,h=1,m=-1),F(0)),
        'z_h_low':(row(z=1,h=-1,t=2),F(0)),
    }
    if region=='H1':
        allowed.update({
            'z_prime':(row(z=1,a=-1),F(0)),
            'z_h1_left':(row(z=1,h=2,m=-1,t=-2),F(0)),
            'z_h1_right':(row(z=1,m=1,t=-4),F(0)),
        })
        # The H1 (m,h) section is the triangle with these three vertices.
        vertices=[(2*t,2*t,t) for t in (lo,hi)]+[(4*t,2*t,t) for t in (lo,hi)]+[(4*t,3*t,t) for t in (lo,hi)]
    else:
        allowed.update({
            'z_prime':(row(z=1,b=-1),F(0)),
            'z_h_high':(row(z=1,h=1,t=-4),F(0)),
            'z_band_high':(row(z=1,m=4,h=3,t=-2),F(2)),
            'z_band_low':(row(z=1,m=-4,h=-3,t=-2),F(-2)),
        })
        # The H2 (m,h) section is the following parallelogram.
        vertices=[]
        for t in (lo,hi):
            vertices.extend([(F(1,2)-2*t,2*t,t),(F(1,2)-t,2*t,t),(F(1,2)-F(7,2)*t,4*t,t),(F(1,2)-F(5,2)*t,4*t,t)])

    def monomial(p):
        assert len(p)==3 and all(isinstance(x,int) and x>=0 for x in p)
        assert sum(p)>0
        lm,lh,l0=p[0]-p[2],p[1]-p[2],p[2]
        diag=(row(a=-2*p[0],b=-2*p[1],c=-2*p[2],f=1),F(0))
        mean=(row(m=lm,h=lh,a=-2*p[0],b=-2*p[1],c=-2*p[2],f=1,t=1),F(1-l0))
        halasz=(row(m=2*lm,h=2*lh,a=-6*p[0],b=-6*p[1],c=-6*p[2],f=1,t=1),F(1-2*l0))
        lengths=[(l0+lm*m+lh*h,t) for m,h,t in vertices]
        return diag,mean,halasz,lengths

    base=[unpack(x) for x in data['base']]
    seen_regions=set()
    for a,b,label in base:
        if isinstance(label,str):
            assert label in allowed and (a,b)==allowed[label],label
            seen_regions.add(label)
        elif label[0]=='nonnegative':
            j=label[1];assert isinstance(j,int) and 0<=j<N
            rr=[F(0)]*N;rr[j]=-1
            assert (a,b)==(rr,F(0))
        else:
            kind,p=label;diag,mean,halasz,lengths=monomial(p)
            if kind=='mean':
                assert ((a,b)==diag and min(L+t for L,t in lengths)>=1) or ((a,b)==mean and max(L+t for L,t in lengths)<=1),label
            elif kind=='halasz':
                assert (a,b)==halasz and max(2*L+t for L,t in lengths)<=1,label
            else:raise AssertionError(label)
    assert seen_regions==set(allowed)
    assert {label[1] for a,b,label in base if isinstance(label,list) and label[0]=='nonnegative'}==set(range(N))

    dis=[]
    for pair in data['disjunctions']:
        assert len(pair)==2
        aa=unpack(pair[0]);bb=unpack(pair[1])
        ka,p=aa[2];kb,pb=bb[2]
        assert p==pb
        diag,mean,halasz,lengths=monomial(p)
        if ka=='mean_diag':
            assert kb=='mean_off' and (aa[0],aa[1])==diag and (bb[0],bb[1])==mean
        else:
            assert ka=='halasz_diag' and kb=='halasz_off' and (aa[0],aa[1])==diag and (bb[0],bb[1])==halasz
        dis.append([aa,bb])
    target=[F(x) for x in data['target']]
    assert target==row(a=-2,b=-2,c=-2,f=1,z=F(summary['gain']))
    assert F(summary['gain'])>0
    return base,dis,target

def check_tree(paths):
    assert len(paths)==len(set(paths))
    def rec(ps):
        if () in ps:
            assert len(ps)==1
            return 1,0
        first={p[0][0] for p in ps};assert len(first)==1
        sides={p[0][1] for p in ps};assert sides=={0,1}
        leaves=0;depth=0
        for side in [0,1]:
            ls,ds=rec([p[1:] for p in ps if p[0][1]==side])
            leaves+=ls;depth=max(depth,ds+1)
        return leaves,depth
    return rec(paths)

def verify(path):
    data=json.loads(Path(path).read_text())
    base,dis,target=check_provenance(data)
    paths=[];max_weight=F(0);min_bound=F(0)
    for leaf in data['leaves']:
        p=tuple((int(i),int(side)) for i,side in leaf['path'])
        assert len({i for i,side in p})==len(p)
        assert all(0<=i<len(dis) and side in [0,1] for i,side in p)
        paths.append(p)
        cons=base+[dis[i][side] for i,side in p]
        indices=[i for i,v in leaf['dual']['weights']]
        assert len(indices)==len(set(indices))
        sums=[F(0)]*N;bound=F(0);weight=F(0)
        for i,v in leaf['dual']['weights']:
            assert isinstance(i,int) and 0<=i<len(cons)
            v=F(v);assert v>=0
            a,b,_=cons[i]
            sums=[x+v*y for x,y in zip(sums,a)]
            bound+=v*b;weight+=v
        assert sums==target
        assert bound==F(leaf['dual']['bound'])<=0
        max_weight=max(max_weight,weight);min_bound=min(min_bound,bound)
    count,depth=check_tree(paths)
    assert not data['summary']['bad'] and data['summary']['stack']==0
    assert count==data['summary']['leaves']
    assert 2*count-1==data['summary']['nodes']
    assert depth==data['summary']['maxdepth']
    print(json.dumps({'certificate':str(path),'verified':True,'region':data['summary']['region'],'theta_range':[data['summary']['theta_lower'],data['summary']['theta_upper']],'gain':data['summary']['gain'],'leaves':count,'max_depth':depth,'max_dual_weight_sum':str(max_weight),'smallest_leaf_bound':str(min_bound)}))

if __name__=='__main__':
    import sys
    paths=sys.argv[1:] or [str(Path(__file__).parent/f'harman_{r}_uniform_certificate.json') for r in ['H1','H2']]
    for path in paths:verify(path)
\end{filecontents*}
\begin{filecontents*}[overwrite]{harman_H1_uniform_certificate.json}
{"summary": {"region": "H1", "theta_lower": "1/25", "theta_upper": "1/24", "gain": "1/100", "nodes": 139, "leaves": 70, "stack": 0, "maxdepth": 17, "seconds": 0.619877815246582, "bad": []}, "base": [[["-1", "0", "0", "0", "0", "0", "0", "0"], "0", ["nonnegative", 0]], [["0", "-1", "0", "0", "0", "0", "0", "0"], "0", ["nonnegative", 1]], [["0", "0", "-1", "0", "0", "0", "0", "0"], "0", ["nonnegative", 2]], [["0", "0", "0", "-1", "0", "0", "0", "0"], "0", ["nonnegative", 3]], [["0", "0", "0", "0", "-1", "0", "0", "0"], "0", ["nonnegative", 4]], [["0", "0", "0", "0", "0", "-1", "0", "0"], "0", ["nonnegative", 5]], [["0", "0", "0", "0", "0", "0", "-1", "0"], "0", ["nonnegative", 6]], [["0", "0", "0", "0", "0", "0", "0", "-1"], "0", ["nonnegative", 7]], [["0", "0", "0", "0", "0", "0", "0", "-1"], "-1/25", "theta_lower"], [["0", "0", "0", "0", "0", "0", "0", "1"], "1/24", "theta_upper"], [["1", "0", "0", "0", "0", "0", "0", "0"], "1/2", "m_half"], [["-1", "1", "0", "0", "0", "0", "0", "0"], "0", "h_le_m"], [["1", "1", "0", "0", "0", "0", "0", "0"], "1", "sum_lengths"], [["-1", "0", "2", "0", "0", "0", "0", "0"], "0", "a_upper"], [["0", "-1", "0", "2", "0", "0", "0", "0"], "0", "b_upper"], [["1", "1", "0", "0", "2", "0", "0", "0"], "1", "c_upper"], [["0", "0", "0", "0", "0", "1", "0", "1"], "1", "f_upper"], [["1", "0", "0", "0", "0", "0", "1", "0"], "1/2", "z_m_half"], [["-1", "1", "0", "0", "0", "0", "1", "0"], "0", "z_h_le_m"], [["0", "-1", "0", "0", "0", "0", "1", "2"], "0", "z_h_low"], [["0", "0", "-1", "0", "0", "0", "1", "0"], "0", "z_prime"], [["-1", "2", "0", "0", "0", "0", "1", "-2"], "0", "z_h1_left"], [["1", "0", "0", "0", "0", "0", "1", "-4"], "0", "z_h1_right"], [["-1", "-1", "0", "0", "-2", "1", "0", "1"], "0", ["mean", [0, 0, 1]]], [["0", "0", "0", "0", "-4", "1", "0", "0"], "0", ["mean", [0, 0, 2]]], [["0", "1", "0", "-2", "0", "1", "0", "1"], "1", ["mean", [0, 1, 0]]], [["0", "2", "0", "-6", "0", "1", "0", "1"], "1", ["halasz", [0, 1, 0]]], [["-1", "0", "0", "-2", "-2", "1", "0", "1"], "0", ["mean", [0, 1, 1]]], [["0", "2", "0", "-4", "0", "1", "0", "1"], "1", ["mean", [0, 2, 0]]], [["0", "4", "0", "-12", "0", "1", "0", "1"], "1", ["halasz", [0, 2, 0]]], [["0", "3", "0", "-6", "0", "1", "0", "1"], "1", ["mean", [0, 3, 0]]], [["0", "6", "0", "-18", "0", "1", "0", "1"], "1", ["halasz", [0, 3, 0]]], [["0", "0", "0", "-6", "-2", "1", "0", "0"], "0", ["mean", [0, 3, 1]]], [["0", "4", "0", "-8", "0", "1", "0", "1"], "1", ["mean", [0, 4, 0]]], [["0", "0", "0", "-8", "-2", "1", "0", "0"], "0", ["mean", [0, 4, 1]]], [["0", "5", "0", "-10", "0", "1", "0", "1"], "1", ["mean", [0, 5, 0]]], [["0", "0", "0", "-10", "-2", "1", "0", "0"], "0", ["mean", [0, 5, 1]]], [["0", "6", "0", "-12", "0", "1", "0", "1"], "1", ["mean", [0, 6, 0]]], [["0", "0", "0", "-12", "-2", "1", "0", "0"], "0", ["mean", [0, 6, 1]]], [["0", "7", "0", "-14", "0", "1", "0", "1"], "1", ["mean", [0, 7, 0]]], [["0", "0", "0", "-14", "-2", "1", "0", "0"], "0", ["mean", [0, 7, 1]]], [["0", "0", "0", "-16", "-2", "1", "0", "0"], "0", ["mean", [0, 8, 1]]], [["0", "0", "0", "-18", "-2", "1", "0", "0"], "0", ["mean", [0, 9, 1]]], [["0", "0", "0", "-24", "0", "1", "0", "0"], "0", ["mean", [0, 12, 0]]], [["1", "0", "-2", "0", "0", "1", "0", "1"], "1", ["mean", [1, 0, 0]]], [["2", "0", "-6", "0", "0", "1", "0", "1"], "1", ["halasz", [1, 0, 0]]], [["0", "-1", "-2", "0", "-2", "1", "0", "1"], "0", ["mean", [1, 0, 1]]], [["1", "1", "-2", "-2", "0", "1", "0", "1"], "1", ["mean", [1, 1, 0]]], [["2", "2", "-6", "-6", "0", "1", "0", "1"], "1", ["halasz", [1, 1, 0]]], [["0", "0", "-2", "-2", "-2", "1", "0", "0"], "0", ["mean", [1, 1, 1]]], [["1", "2", "-2", "-4", "0", "1", "0", "1"], "1", ["mean", [1, 2, 0]]], [["2", "4", "-6", "-12", "0", "1", "0", "1"], "1", ["halasz", [1, 2, 0]]], [["0", "0", "-2", "-4", "-2", "1", "0", "0"], "0", ["mean", [1, 2, 1]]], [["1", "3", "-2", "-6", "0", "1", "0", "1"], "1", ["mean", [1, 3, 0]]], [["0", "0", "-2", "-6", "-2", "1", "0", "0"], "0", ["mean", [1, 3, 1]]], [["1", "4", "-2", "-8", "0", "1", "0", "1"], "1", ["mean", [1, 4, 0]]], [["0", "0", "-2", "-8", "-2", "1", "0", "0"], "0", ["mean", [1, 4, 1]]], [["1", "5", "-2", "-10", "0", "1", "0", "1"], "1", ["mean", [1, 5, 0]]], [["0", "0", "-2", "-10", "-2", "1", "0", "0"], "0", ["mean", [1, 5, 1]]], [["1", "6", "-2", "-12", "0", "1", "0", "1"], "1", ["mean", [1, 6, 0]]], [["0", "0", "-2", "-12", "-2", "1", "0", "0"], "0", ["mean", [1, 6, 1]]], [["0", "0", "-2", "-14", "-2", "1", "0", "0"], "0", ["mean", [1, 7, 1]]], [["0", "0", "-2", "-22", "0", "1", "0", "0"], "0", ["mean", [1, 11, 0]]], [["0", "0", "-2", "-24", "0", "1", "0", "0"], "0", ["mean", [1, 12, 0]]], [["2", "0", "-4", "0", "0", "1", "0", "1"], "1", ["mean", [2, 0, 0]]], [["4", "0", "-12", "0", "0", "1", "0", "1"], "1", ["halasz", [2, 0, 0]]], [["0", "0", "-4", "0", "-2", "1", "0", "0"], "0", ["mean", [2, 0, 1]]], [["2", "1", "-4", "-2", "0", "1", "0", "1"], "1", ["mean", [2, 1, 0]]], [["4", "2", "-12", "-6", "0", "1", "0", "1"], "1", ["halasz", [2, 1, 0]]], [["0", "0", "-4", "-2", "-2", "1", "0", "0"], "0", ["mean", [2, 1, 1]]], [["2", "2", "-4", "-4", "0", "1", "0", "1"], "1", ["mean", [2, 2, 0]]], [["0", "0", "-4", "-4", "-2", "1", "0", "0"], "0", ["mean", [2, 2, 1]]], [["2", "3", "-4", "-6", "0", "1", "0", "1"], "1", ["mean", [2, 3, 0]]], [["0", "0", "-4", "-6", "-2", "1", "0", "0"], "0", ["mean", [2, 3, 1]]], [["2", "4", "-4", "-8", "0", "1", "0", "1"], "1", ["mean", [2, 4, 0]]], [["0", "0", "-4", "-8", "-2", "1", "0", "0"], "0", ["mean", [2, 4, 1]]], [["2", "5", "-4", "-10", "0", "1", "0", "1"], "1", ["mean", [2, 5, 0]]], [["0", "0", "-4", "-10", "-2", "1", "0", "0"], "0", ["mean", [2, 5, 1]]], [["0", "0", "-4", "-12", "-2", "1", "0", "0"], "0", ["mean", [2, 6, 1]]], [["0", "0", "-4", "-20", "0", "1", "0", "0"], "0", ["mean", [2, 10, 0]]], [["0", "0", "-4", "-22", "0", "1", "0", "0"], "0", ["mean", [2, 11, 0]]], [["0", "0", "-4", "-24", "0", "1", "0", "0"], "0", ["mean", [2, 12, 0]]], [["3", "0", "-6", "0", "0", "1", "0", "1"], "1", ["mean", [3, 0, 0]]], [["0", "0", "-6", "0", "-2", "1", "0", "0"], "0", ["mean", [3, 0, 1]]], [["3", "1", "-6", "-2", "0", "1", "0", "1"], "1", ["mean", [3, 1, 0]]], [["0", "0", "-6", "-2", "-2", "1", "0", "0"], "0", ["mean", [3, 1, 1]]], [["3", "2", "-6", "-4", "0", "1", "0", "1"], "1", ["mean", [3, 2, 0]]], [["0", "0", "-6", "-4", "-2", "1", "0", "0"], "0", ["mean", [3, 2, 1]]], [["3", "3", "-6", "-6", "0", "1", "0", "1"], "1", ["mean", [3, 3, 0]]], [["0", "0", "-6", "-6", "-2", "1", "0", "0"], "0", ["mean", [3, 3, 1]]], [["0", "0", "-6", "-8", "-2", "1", "0", "0"], "0", ["mean", [3, 4, 1]]], [["0", "0", "-6", "-10", "-2", "1", "0", "0"], "0", ["mean", [3, 5, 1]]], [["0", "0", "-6", "-18", "0", "1", "0", "0"], "0", ["mean", [3, 9, 0]]], [["0", "0", "-6", "-20", "0", "1", "0", "0"], "0", ["mean", [3, 10, 0]]], [["0", "0", "-6", "-22", "0", "1", "0", "0"], "0", ["mean", [3, 11, 0]]], [["0", "0", "-6", "-24", "0", "1", "0", "0"], "0", ["mean", [3, 12, 0]]], [["4", "0", "-8", "0", "0", "1", "0", "1"], "1", ["mean", [4, 0, 0]]], [["0", "0", "-8", "0", "-2", "1", "0", "0"], "0", ["mean", [4, 0, 1]]], [["4", "1", "-8", "-2", "0", "1", "0", "1"], "1", ["mean", [4, 1, 0]]], [["0", "0", "-8", "-2", "-2", "1", "0", "0"], "0", ["mean", [4, 1, 1]]], [["4", "2", "-8", "-4", "0", "1", "0", "1"], "1", ["mean", [4, 2, 0]]], [["0", "0", "-8", "-4", "-2", "1", "0", "0"], "0", ["mean", [4, 2, 1]]], [["0", "0", "-8", "-6", "-2", "1", "0", "0"], "0", ["mean", [4, 3, 1]]], [["0", "0", "-8", "-8", "-2", "1", "0", "0"], "0", ["mean", [4, 4, 1]]], [["0", "0", "-8", "-16", "0", "1", "0", "0"], "0", ["mean", [4, 8, 0]]], [["0", "0", "-8", "-18", "0", "1", "0", "0"], "0", ["mean", [4, 9, 0]]], [["0", "0", "-8", "-20", "0", "1", "0", "0"], "0", ["mean", [4, 10, 0]]], [["0", "0", "-8", "-22", "0", "1", "0", "0"], "0", ["mean", [4, 11, 0]]], [["0", "0", "-8", "-24", "0", "1", "0", "0"], "0", ["mean", [4, 12, 0]]], [["5", "0", "-10", "0", "0", "1", "0", "1"], "1", ["mean", [5, 0, 0]]], [["0", "0", "-10", "0", "-2", "1", "0", "0"], "0", ["mean", [5, 0, 1]]], [["5", "1", "-10", "-2", "0", "1", "0", "1"], "1", ["mean", [5, 1, 0]]], [["0", "0", "-10", "-2", "-2", "1", "0", "0"], "0", ["mean", [5, 1, 1]]], [["0", "0", "-10", "-4", "-2", "1", "0", "0"], "0", ["mean", [5, 2, 1]]], [["0", "0", "-10", "-6", "-2", "1", "0", "0"], "0", ["mean", [5, 3, 1]]], [["0", "0", "-10", "-14", "0", "1", "0", "0"], "0", ["mean", [5, 7, 0]]], [["0", "0", "-10", "-16", "0", "1", "0", "0"], "0", ["mean", [5, 8, 0]]], [["0", "0", "-10", "-18", "0", "1", "0", "0"], "0", ["mean", [5, 9, 0]]], [["0", "0", "-10", "-20", "0", "1", "0", "0"], "0", ["mean", [5, 10, 0]]], [["0", "0", "-10", "-22", "0", "1", "0", "0"], "0", ["mean", [5, 11, 0]]], [["0", "0", "-10", "-24", "0", "1", "0", "0"], "0", ["mean", [5, 12, 0]]], [["0", "0", "-12", "0", "-2", "1", "0", "0"], "0", ["mean", [6, 0, 1]]], [["0", "0", "-12", "-2", "-2", "1", "0", "0"], "0", ["mean", [6, 1, 1]]], [["0", "0", "-12", "-4", "-2", "1", "0", "0"], "0", ["mean", [6, 2, 1]]], [["0", "0", "-12", "-12", "0", "1", "0", "0"], "0", ["mean", [6, 6, 0]]], [["0", "0", "-12", "-14", "0", "1", "0", "0"], "0", ["mean", [6, 7, 0]]], [["0", "0", "-12", "-16", "0", "1", "0", "0"], "0", ["mean", [6, 8, 0]]], [["0", "0", "-12", "-18", "0", "1", "0", "0"], "0", ["mean", [6, 9, 0]]], [["0", "0", "-12", "-20", "0", "1", "0", "0"], "0", ["mean", [6, 10, 0]]], [["0", "0", "-12", "-22", "0", "1", "0", "0"], "0", ["mean", [6, 11, 0]]], [["0", "0", "-12", "-24", "0", "1", "0", "0"], "0", ["mean", [6, 12, 0]]], [["0", "0", "-14", "0", "-2", "1", "0", "0"], "0", ["mean", [7, 0, 1]]], [["0", "0", "-14", "-2", "-2", "1", "0", "0"], "0", ["mean", [7, 1, 1]]], [["0", "0", "-14", "-10", "0", "1", "0", "0"], "0", ["mean", [7, 5, 0]]], [["0", "0", "-14", "-12", "0", "1", "0", "0"], "0", ["mean", [7, 6, 0]]], [["0", "0", "-14", "-14", "0", "1", "0", "0"], "0", ["mean", [7, 7, 0]]], [["0", "0", "-14", "-16", "0", "1", "0", "0"], "0", ["mean", [7, 8, 0]]], [["0", "0", "-14", "-18", "0", "1", "0", "0"], "0", ["mean", [7, 9, 0]]], [["0", "0", "-14", "-20", "0", "1", "0", "0"], "0", ["mean", [7, 10, 0]]], [["0", "0", "-14", "-22", "0", "1", "0", "0"], "0", ["mean", [7, 11, 0]]], [["0", "0", "-16", "0", "-2", "1", "0", "0"], "0", ["mean", [8, 0, 1]]], [["0", "0", "-16", "-8", "0", "1", "0", "0"], "0", ["mean", [8, 4, 0]]], [["0", "0", "-16", "-10", "0", "1", "0", "0"], "0", ["mean", [8, 5, 0]]], [["0", "0", "-16", "-12", "0", "1", "0", "0"], "0", ["mean", [8, 6, 0]]], [["0", "0", "-16", "-14", "0", "1", "0", "0"], "0", ["mean", [8, 7, 0]]], [["0", "0", "-16", "-16", "0", "1", "0", "0"], "0", ["mean", [8, 8, 0]]], [["0", "0", "-16", "-18", "0", "1", "0", "0"], "0", ["mean", [8, 9, 0]]], [["0", "0", "-16", "-20", "0", "1", "0", "0"], "0", ["mean", [8, 10, 0]]], [["0", "0", "-18", "-6", "0", "1", "0", "0"], "0", ["mean", [9, 3, 0]]], [["0", "0", "-18", "-8", "0", "1", "0", "0"], "0", ["mean", [9, 4, 0]]], [["0", "0", "-18", "-10", "0", "1", "0", "0"], "0", ["mean", [9, 5, 0]]], [["0", "0", "-18", "-12", "0", "1", "0", "0"], "0", ["mean", [9, 6, 0]]], [["0", "0", "-18", "-14", "0", "1", "0", "0"], "0", ["mean", [9, 7, 0]]], [["0", "0", "-18", "-16", "0", "1", "0", "0"], "0", ["mean", [9, 8, 0]]], [["0", "0", "-18", "-18", "0", "1", "0", "0"], "0", ["mean", [9, 9, 0]]], [["0", "0", "-20", "-4", "0", "1", "0", "0"], "0", ["mean", [10, 2, 0]]], [["0", "0", "-20", "-6", "0", "1", "0", "0"], "0", ["mean", [10, 3, 0]]], [["0", "0", "-20", "-8", "0", "1", "0", "0"], "0", ["mean", [10, 4, 0]]], [["0", "0", "-20", "-10", "0", "1", "0", "0"], "0", ["mean", [10, 5, 0]]], [["0", "0", "-20", "-12", "0", "1", "0", "0"], "0", ["mean", [10, 6, 0]]], [["0", "0", "-20", "-14", "0", "1", "0", "0"], "0", ["mean", [10, 7, 0]]], [["0", "0", "-20", "-16", "0", "1", "0", "0"], "0", ["mean", [10, 8, 0]]], [["0", "0", "-22", "-2", "0", "1", "0", "0"], "0", ["mean", [11, 1, 0]]], [["0", "0", "-22", "-4", "0", "1", "0", "0"], "0", ["mean", [11, 2, 0]]], [["0", "0", "-22", "-6", "0", "1", "0", "0"], "0", ["mean", [11, 3, 0]]], [["0", "0", "-22", "-8", "0", "1", "0", "0"], "0", ["mean", [11, 4, 0]]], [["0", "0", "-22", "-10", "0", "1", "0", "0"], "0", ["mean", [11, 5, 0]]], [["0", "0", "-22", "-12", "0", "1", "0", "0"], "0", ["mean", [11, 6, 0]]], [["0", "0", "-22", "-14", "0", "1", "0", "0"], "0", ["mean", [11, 7, 0]]], [["0", "0", "-24", "0", "0", "1", "0", "0"], "0", ["mean", [12, 0, 0]]], [["0", "0", "-24", "-2", "0", "1", "0", "0"], "0", ["mean", [12, 1, 0]]], [["0", "0", "-24", "-4", "0", "1", "0", "0"], "0", ["mean", [12, 2, 0]]], [["0", "0", "-24", "-6", "0", "1", "0", "0"], "0", ["mean", [12, 3, 0]]], [["0", "0", "-24", "-8", "0", "1", "0", "0"], "0", ["mean", [12, 4, 0]]], [["0", "0", "-24", "-10", "0", "1", "0", "0"], "0", ["mean", [12, 5, 0]]], [["0", "0", "-24", "-12", "0", "1", "0", "0"], "0", ["mean", [12, 6, 0]]]], "disjunctions": [[[["0", "0", "0", "0", "-2", "1", "0", "0"], "0", ["halasz_diag", [0, 0, 1]]], [["-2", "-2", "0", "0", "-6", "1", "0", "1"], "-1", ["halasz_off", [0, 0, 1]]]], [[["0", "0", "0", "0", "-4", "1", "0", "0"], "0", ["halasz_diag", [0, 0, 2]]], [["-4", "-4", "0", "0", "-12", "1", "0", "1"], "-3", ["halasz_off", [0, 0, 2]]]], [[["0", "0", "0", "-2", "-2", "1", "0", "0"], "0", ["halasz_diag", [0, 1, 1]]], [["-2", "0", "0", "-6", "-6", "1", "0", "1"], "-1", ["halasz_off", [0, 1, 1]]]], [[["0", "0", "0", "-4", "-2", "1", "0", "0"], "0", ["mean_diag", [0, 2, 1]]], [["-1", "1", "0", "-4", "-2", "1", "0", "1"], "0", ["mean_off", [0, 2, 1]]]], [[["0", "0", "0", "-4", "-2", "1", "0", "0"], "0", ["halasz_diag", [0, 2, 1]]], [["-2", "2", "0", "-12", "-6", "1", "0", "1"], "-1", ["halasz_off", [0, 2, 1]]]], [[["0", "0", "0", "-6", "-2", "1", "0", "0"], "0", ["halasz_diag", [0, 3, 1]]], [["-2", "4", "0", "-18", "-6", "1", "0", "1"], "-1", ["halasz_off", [0, 3, 1]]]], [[["0", "0", "0", "-8", "0", "1", "0", "0"], "0", ["halasz_diag", [0, 4, 0]]], [["0", "8", "0", "-24", "0", "1", "0", "1"], "1", ["halasz_off", [0, 4, 0]]]], [[["0", "0", "0", "-8", "-2", "1", "0", "0"], "0", ["halasz_diag", [0, 4, 1]]], [["-2", "6", "0", "-24", "-6", "1", "0", "1"], "-1", ["halasz_off", [0, 4, 1]]]], [[["0", "0", "0", "-10", "0", "1", "0", "0"], "0", ["halasz_diag", [0, 5, 0]]], [["0", "10", "0", "-30", "0", "1", "0", "1"], "1", ["halasz_off", [0, 5, 0]]]], [[["0", "0", "0", "-10", "-2", "1", "0", "0"], "0", ["halasz_diag", [0, 5, 1]]], [["-2", "8", "0", "-30", "-6", "1", "0", "1"], "-1", ["halasz_off", [0, 5, 1]]]], [[["0", "0", "0", "-12", "0", "1", "0", "0"], "0", ["halasz_diag", [0, 6, 0]]], [["0", "12", "0", "-36", "0", "1", "0", "1"], "1", ["halasz_off", [0, 6, 0]]]], [[["0", "0", "0", "-12", "-2", "1", "0", "0"], "0", ["halasz_diag", [0, 6, 1]]], [["-2", "10", "0", "-36", "-6", "1", "0", "1"], "-1", ["halasz_off", [0, 6, 1]]]], [[["0", "0", "0", "-14", "0", "1", "0", "0"], "0", ["halasz_diag", [0, 7, 0]]], [["0", "14", "0", "-42", "0", "1", "0", "1"], "1", ["halasz_off", [0, 7, 0]]]], [[["0", "0", "0", "-14", "-2", "1", "0", "0"], "0", ["halasz_diag", [0, 7, 1]]], [["-2", "12", "0", "-42", "-6", "1", "0", "1"], "-1", ["halasz_off", [0, 7, 1]]]], [[["0", "0", "0", "-16", "0", "1", "0", "0"], "0", ["mean_diag", [0, 8, 0]]], [["0", "8", "0", "-16", "0", "1", "0", "1"], "1", ["mean_off", [0, 8, 0]]]], [[["0", "0", "0", "-16", "0", "1", "0", "0"], "0", ["halasz_diag", [0, 8, 0]]], [["0", "16", "0", "-48", "0", "1", "0", "1"], "1", ["halasz_off", [0, 8, 0]]]], [[["0", "0", "0", "-16", "-2", "1", "0", "0"], "0", ["halasz_diag", [0, 8, 1]]], [["-2", "14", "0", "-48", "-6", "1", "0", "1"], "-1", ["halasz_off", [0, 8, 1]]]], [[["0", "0", "0", "-18", "0", "1", "0", "0"], "0", ["mean_diag", [0, 9, 0]]], [["0", "9", "0", "-18", "0", "1", "0", "1"], "1", ["mean_off", [0, 9, 0]]]], [[["0", "0", "0", "-18", "0", "1", "0", "0"], "0", ["halasz_diag", [0, 9, 0]]], [["0", "18", "0", "-54", "0", "1", "0", "1"], "1", ["halasz_off", [0, 9, 0]]]], [[["0", "0", "0", "-18", "-2", "1", "0", "0"], "0", ["halasz_diag", [0, 9, 1]]], [["-2", "16", "0", "-54", "-6", "1", "0", "1"], "-1", ["halasz_off", [0, 9, 1]]]], [[["0", "0", "0", "-20", "0", "1", "0", "0"], "0", ["mean_diag", [0, 10, 0]]], [["0", "10", "0", "-20", "0", "1", "0", "1"], "1", ["mean_off", [0, 10, 0]]]], [[["0", "0", "0", "-20", "0", "1", "0", "0"], "0", ["halasz_diag", [0, 10, 0]]], [["0", "20", "0", "-60", "0", "1", "0", "1"], "1", ["halasz_off", [0, 10, 0]]]], [[["0", "0", "0", "-22", "0", "1", "0", "0"], "0", ["mean_diag", [0, 11, 0]]], [["0", "11", "0", "-22", "0", "1", "0", "1"], "1", ["mean_off", [0, 11, 0]]]], [[["0", "0", "0", "-22", "0", "1", "0", "0"], "0", ["halasz_diag", [0, 11, 0]]], [["0", "22", "0", "-66", "0", "1", "0", "1"], "1", ["halasz_off", [0, 11, 0]]]], [[["0", "0", "0", "-24", "0", "1", "0", "0"], "0", ["halasz_diag", [0, 12, 0]]], [["0", "24", "0", "-72", "0", "1", "0", "1"], "1", ["halasz_off", [0, 12, 0]]]], [[["0", "0", "-2", "0", "-2", "1", "0", "0"], "0", ["halasz_diag", [1, 0, 1]]], [["0", "-2", "-6", "0", "-6", "1", "0", "1"], "-1", ["halasz_off", [1, 0, 1]]]], [[["0", "0", "-2", "-2", "-2", "1", "0", "0"], "0", ["halasz_diag", [1, 1, 1]]], [["0", "0", "-6", "-6", "-6", "1", "0", "1"], "-1", ["halasz_off", [1, 1, 1]]]], [[["0", "0", "-2", "-4", "-2", "1", "0", "0"], "0", ["halasz_diag", [1, 2, 1]]], [["0", "2", "-6", "-12", "-6", "1", "0", "1"], "-1", ["halasz_off", [1, 2, 1]]]], [[["0", "0", "-2", "-6", "0", "1", "0", "0"], "0", ["halasz_diag", [1, 3, 0]]], [["2", "6", "-6", "-18", "0", "1", "0", "1"], "1", ["halasz_off", [1, 3, 0]]]], [[["0", "0", "-2", "-6", "-2", "1", "0", "0"], "0", ["halasz_diag", [1, 3, 1]]], [["0", "4", "-6", "-18", "-6", "1", "0", "1"], "-1", ["halasz_off", [1, 3, 1]]]], [[["0", "0", "-2", "-8", "0", "1", "0", "0"], "0", ["halasz_diag", [1, 4, 0]]], [["2", "8", "-6", "-24", "0", "1", "0", "1"], "1", ["halasz_off", [1, 4, 0]]]], [[["0", "0", "-2", "-8", "-2", "1", "0", "0"], "0", ["halasz_diag", [1, 4, 1]]], [["0", "6", "-6", "-24", "-6", "1", "0", "1"], "-1", ["halasz_off", [1, 4, 1]]]], [[["0", "0", "-2", "-10", "0", "1", "0", "0"], "0", ["halasz_diag", [1, 5, 0]]], [["2", "10", "-6", "-30", "0", "1", "0", "1"], "1", ["halasz_off", [1, 5, 0]]]], [[["0", "0", "-2", "-10", "-2", "1", "0", "0"], "0", ["halasz_diag", [1, 5, 1]]], [["0", "8", "-6", "-30", "-6", "1", "0", "1"], "-1", ["halasz_off", [1, 5, 1]]]], [[["0", "0", "-2", "-12", "0", "1", "0", "0"], "0", ["halasz_diag", [1, 6, 0]]], [["2", "12", "-6", "-36", "0", "1", "0", "1"], "1", ["halasz_off", [1, 6, 0]]]], [[["0", "0", "-2", "-12", "-2", "1", "0", "0"], "0", ["halasz_diag", [1, 6, 1]]], [["0", "10", "-6", "-36", "-6", "1", "0", "1"], "-1", ["halasz_off", [1, 6, 1]]]], [[["0", "0", "-2", "-14", "0", "1", "0", "0"], "0", ["mean_diag", [1, 7, 0]]], [["1", "7", "-2", "-14", "0", "1", "0", "1"], "1", ["mean_off", [1, 7, 0]]]], [[["0", "0", "-2", "-14", "0", "1", "0", "0"], "0", ["halasz_diag", [1, 7, 0]]], [["2", "14", "-6", "-42", "0", "1", "0", "1"], "1", ["halasz_off", [1, 7, 0]]]], [[["0", "0", "-2", "-14", "-2", "1", "0", "0"], "0", ["halasz_diag", [1, 7, 1]]], [["0", "12", "-6", "-42", "-6", "1", "0", "1"], "-1", ["halasz_off", [1, 7, 1]]]], [[["0", "0", "-2", "-16", "0", "1", "0", "0"], "0", ["mean_diag", [1, 8, 0]]], [["1", "8", "-2", "-16", "0", "1", "0", "1"], "1", ["mean_off", [1, 8, 0]]]], [[["0", "0", "-2", "-16", "0", "1", "0", "0"], "0", ["halasz_diag", [1, 8, 0]]], [["2", "16", "-6", "-48", "0", "1", "0", "1"], "1", ["halasz_off", [1, 8, 0]]]], [[["0", "0", "-2", "-18", "0", "1", "0", "0"], "0", ["mean_diag", [1, 9, 0]]], [["1", "9", "-2", "-18", "0", "1", "0", "1"], "1", ["mean_off", [1, 9, 0]]]], [[["0", "0", "-2", "-18", "0", "1", "0", "0"], "0", ["halasz_diag", [1, 9, 0]]], [["2", "18", "-6", "-54", "0", "1", "0", "1"], "1", ["halasz_off", [1, 9, 0]]]], [[["0", "0", "-2", "-20", "0", "1", "0", "0"], "0", ["mean_diag", [1, 10, 0]]], [["1", "10", "-2", "-20", "0", "1", "0", "1"], "1", ["mean_off", [1, 10, 0]]]], [[["0", "0", "-2", "-20", "0", "1", "0", "0"], "0", ["halasz_diag", [1, 10, 0]]], [["2", "20", "-6", "-60", "0", "1", "0", "1"], "1", ["halasz_off", [1, 10, 0]]]], [[["0", "0", "-2", "-22", "0", "1", "0", "0"], "0", ["halasz_diag", [1, 11, 0]]], [["2", "22", "-6", "-66", "0", "1", "0", "1"], "1", ["halasz_off", [1, 11, 0]]]], [[["0", "0", "-2", "-24", "0", "1", "0", "0"], "0", ["halasz_diag", [1, 12, 0]]], [["2", "24", "-6", "-72", "0", "1", "0", "1"], "1", ["halasz_off", [1, 12, 0]]]], [[["0", "0", "-4", "0", "-2", "1", "0", "0"], "0", ["halasz_diag", [2, 0, 1]]], [["2", "-2", "-12", "0", "-6", "1", "0", "1"], "-1", ["halasz_off", [2, 0, 1]]]], [[["0", "0", "-4", "-2", "-2", "1", "0", "0"], "0", ["halasz_diag", [2, 1, 1]]], [["2", "0", "-12", "-6", "-6", "1", "0", "1"], "-1", ["halasz_off", [2, 1, 1]]]], [[["0", "0", "-4", "-4", "0", "1", "0", "0"], "0", ["halasz_diag", [2, 2, 0]]], [["4", "4", "-12", "-12", "0", "1", "0", "1"], "1", ["halasz_off", [2, 2, 0]]]], [[["0", "0", "-4", "-4", "-2", "1", "0", "0"], "0", ["halasz_diag", [2, 2, 1]]], [["2", "2", "-12", "-12", "-6", "1", "0", "1"], "-1", ["halasz_off", [2, 2, 1]]]], [[["0", "0", "-4", "-6", "0", "1", "0", "0"], "0", ["halasz_diag", [2, 3, 0]]], [["4", "6", "-12", "-18", "0", "1", "0", "1"], "1", ["halasz_off", [2, 3, 0]]]], [[["0", "0", "-4", "-6", "-2", "1", "0", "0"], "0", ["halasz_diag", [2, 3, 1]]], [["2", "4", "-12", "-18", "-6", "1", "0", "1"], "-1", ["halasz_off", [2, 3, 1]]]], [[["0", "0", "-4", "-8", "0", "1", "0", "0"], "0", ["halasz_diag", [2, 4, 0]]], [["4", "8", "-12", "-24", "0", "1", "0", "1"], "1", ["halasz_off", [2, 4, 0]]]], [[["0", "0", "-4", "-8", "-2", "1", "0", "0"], "0", ["halasz_diag", [2, 4, 1]]], [["2", "6", "-12", "-24", "-6", "1", "0", "1"], "-1", ["halasz_off", [2, 4, 1]]]], [[["0", "0", "-4", "-10", "0", "1", "0", "0"], "0", ["halasz_diag", [2, 5, 0]]], [["4", "10", "-12", "-30", "0", "1", "0", "1"], "1", ["halasz_off", [2, 5, 0]]]], [[["0", "0", "-4", "-10", "-2", "1", "0", "0"], "0", ["halasz_diag", [2, 5, 1]]], [["2", "8", "-12", "-30", "-6", "1", "0", "1"], "-1", ["halasz_off", [2, 5, 1]]]], [[["0", "0", "-4", "-12", "0", "1", "0", "0"], "0", ["mean_diag", [2, 6, 0]]], [["2", "6", "-4", "-12", "0", "1", "0", "1"], "1", ["mean_off", [2, 6, 0]]]], [[["0", "0", "-4", "-12", "0", "1", "0", "0"], "0", ["halasz_diag", [2, 6, 0]]], [["4", "12", "-12", "-36", "0", "1", "0", "1"], "1", ["halasz_off", [2, 6, 0]]]], [[["0", "0", "-4", "-12", "-2", "1", "0", "0"], "0", ["halasz_diag", [2, 6, 1]]], [["2", "10", "-12", "-36", "-6", "1", "0", "1"], "-1", ["halasz_off", [2, 6, 1]]]], [[["0", "0", "-4", "-14", "0", "1", "0", "0"], "0", ["mean_diag", [2, 7, 0]]], [["2", "7", "-4", "-14", "0", "1", "0", "1"], "1", ["mean_off", [2, 7, 0]]]], [[["0", "0", "-4", "-14", "0", "1", "0", "0"], "0", ["halasz_diag", [2, 7, 0]]], [["4", "14", "-12", "-42", "0", "1", "0", "1"], "1", ["halasz_off", [2, 7, 0]]]], [[["0", "0", "-4", "-16", "0", "1", "0", "0"], "0", ["mean_diag", [2, 8, 0]]], [["2", "8", "-4", "-16", "0", "1", "0", "1"], "1", ["mean_off", [2, 8, 0]]]], [[["0", "0", "-4", "-16", "0", "1", "0", "0"], "0", ["halasz_diag", [2, 8, 0]]], [["4", "16", "-12", "-48", "0", "1", "0", "1"], "1", ["halasz_off", [2, 8, 0]]]], [[["0", "0", "-4", "-18", "0", "1", "0", "0"], "0", ["mean_diag", [2, 9, 0]]], [["2", "9", "-4", "-18", "0", "1", "0", "1"], "1", ["mean_off", [2, 9, 0]]]], [[["0", "0", "-4", "-18", "0", "1", "0", "0"], "0", ["halasz_diag", [2, 9, 0]]], [["4", "18", "-12", "-54", "0", "1", "0", "1"], "1", ["halasz_off", [2, 9, 0]]]], [[["0", "0", "-4", "-20", "0", "1", "0", "0"], "0", ["halasz_diag", [2, 10, 0]]], [["4", "20", "-12", "-60", "0", "1", "0", "1"], "1", ["halasz_off", [2, 10, 0]]]], [[["0", "0", "-4", "-22", "0", "1", "0", "0"], "0", ["halasz_diag", [2, 11, 0]]], [["4", "22", "-12", "-66", "0", "1", "0", "1"], "1", ["halasz_off", [2, 11, 0]]]], [[["0", "0", "-4", "-24", "0", "1", "0", "0"], "0", ["halasz_diag", [2, 12, 0]]], [["4", "24", "-12", "-72", "0", "1", "0", "1"], "1", ["halasz_off", [2, 12, 0]]]], [[["0", "0", "-6", "0", "0", "1", "0", "0"], "0", ["halasz_diag", [3, 0, 0]]], [["6", "0", "-18", "0", "0", "1", "0", "1"], "1", ["halasz_off", [3, 0, 0]]]], [[["0", "0", "-6", "0", "-2", "1", "0", "0"], "0", ["halasz_diag", [3, 0, 1]]], [["4", "-2", "-18", "0", "-6", "1", "0", "1"], "-1", ["halasz_off", [3, 0, 1]]]], [[["0", "0", "-6", "-2", "0", "1", "0", "0"], "0", ["halasz_diag", [3, 1, 0]]], [["6", "2", "-18", "-6", "0", "1", "0", "1"], "1", ["halasz_off", [3, 1, 0]]]], [[["0", "0", "-6", "-2", "-2", "1", "0", "0"], "0", ["halasz_diag", [3, 1, 1]]], [["4", "0", "-18", "-6", "-6", "1", "0", "1"], "-1", ["halasz_off", [3, 1, 1]]]], [[["0", "0", "-6", "-4", "0", "1", "0", "0"], "0", ["halasz_diag", [3, 2, 0]]], [["6", "4", "-18", "-12", "0", "1", "0", "1"], "1", ["halasz_off", [3, 2, 0]]]], [[["0", "0", "-6", "-4", "-2", "1", "0", "0"], "0", ["halasz_diag", [3, 2, 1]]], [["4", "2", "-18", "-12", "-6", "1", "0", "1"], "-1", ["halasz_off", [3, 2, 1]]]], [[["0", "0", "-6", "-6", "0", "1", "0", "0"], "0", ["halasz_diag", [3, 3, 0]]], [["6", "6", "-18", "-18", "0", "1", "0", "1"], "1", ["halasz_off", [3, 3, 0]]]], [[["0", "0", "-6", "-6", "-2", "1", "0", "0"], "0", ["halasz_diag", [3, 3, 1]]], [["4", "4", "-18", "-18", "-6", "1", "0", "1"], "-1", ["halasz_off", [3, 3, 1]]]], [[["0", "0", "-6", "-8", "0", "1", "0", "0"], "0", ["mean_diag", [3, 4, 0]]], [["3", "4", "-6", "-8", "0", "1", "0", "1"], "1", ["mean_off", [3, 4, 0]]]], [[["0", "0", "-6", "-8", "0", "1", "0", "0"], "0", ["halasz_diag", [3, 4, 0]]], [["6", "8", "-18", "-24", "0", "1", "0", "1"], "1", ["halasz_off", [3, 4, 0]]]], [[["0", "0", "-6", "-8", "-2", "1", "0", "0"], "0", ["halasz_diag", [3, 4, 1]]], [["4", "6", "-18", "-24", "-6", "1", "0", "1"], "-1", ["halasz_off", [3, 4, 1]]]], [[["0", "0", "-6", "-10", "0", "1", "0", "0"], "0", ["mean_diag", [3, 5, 0]]], [["3", "5", "-6", "-10", "0", "1", "0", "1"], "1", ["mean_off", [3, 5, 0]]]], [[["0", "0", "-6", "-10", "0", "1", "0", "0"], "0", ["halasz_diag", [3, 5, 0]]], [["6", "10", "-18", "-30", "0", "1", "0", "1"], "1", ["halasz_off", [3, 5, 0]]]], [[["0", "0", "-6", "-10", "-2", "1", "0", "0"], "0", ["halasz_diag", [3, 5, 1]]], [["4", "8", "-18", "-30", "-6", "1", "0", "1"], "-1", ["halasz_off", [3, 5, 1]]]], [[["0", "0", "-6", "-12", "0", "1", "0", "0"], "0", ["mean_diag", [3, 6, 0]]], [["3", "6", "-6", "-12", "0", "1", "0", "1"], "1", ["mean_off", [3, 6, 0]]]], [[["0", "0", "-6", "-12", "0", "1", "0", "0"], "0", ["halasz_diag", [3, 6, 0]]], [["6", "12", "-18", "-36", "0", "1", "0", "1"], "1", ["halasz_off", [3, 6, 0]]]], [[["0", "0", "-6", "-14", "0", "1", "0", "0"], "0", ["mean_diag", [3, 7, 0]]], [["3", "7", "-6", "-14", "0", "1", "0", "1"], "1", ["mean_off", [3, 7, 0]]]], [[["0", "0", "-6", "-14", "0", "1", "0", "0"], "0", ["halasz_diag", [3, 7, 0]]], [["6", "14", "-18", "-42", "0", "1", "0", "1"], "1", ["halasz_off", [3, 7, 0]]]], [[["0", "0", "-6", "-16", "0", "1", "0", "0"], "0", ["mean_diag", [3, 8, 0]]], [["3", "8", "-6", "-16", "0", "1", "0", "1"], "1", ["mean_off", [3, 8, 0]]]], [[["0", "0", "-6", "-16", "0", "1", "0", "0"], "0", ["halasz_diag", [3, 8, 0]]], [["6", "16", "-18", "-48", "0", "1", "0", "1"], "1", ["halasz_off", [3, 8, 0]]]], [[["0", "0", "-6", "-18", "0", "1", "0", "0"], "0", ["halasz_diag", [3, 9, 0]]], [["6", "18", "-18", "-54", "0", "1", "0", "1"], "1", ["halasz_off", [3, 9, 0]]]], [[["0", "0", "-6", "-20", "0", "1", "0", "0"], "0", ["halasz_diag", [3, 10, 0]]], [["6", "20", "-18", "-60", "0", "1", "0", "1"], "1", ["halasz_off", [3, 10, 0]]]], [[["0", "0", "-6", "-22", "0", "1", "0", "0"], "0", ["halasz_diag", [3, 11, 0]]], [["6", "22", "-18", "-66", "0", "1", "0", "1"], "1", ["halasz_off", [3, 11, 0]]]], [[["0", "0", "-6", "-24", "0", "1", "0", "0"], "0", ["halasz_diag", [3, 12, 0]]], [["6", "24", "-18", "-72", "0", "1", "0", "1"], "1", ["halasz_off", [3, 12, 0]]]], [[["0", "0", "-8", "0", "0", "1", "0", "0"], "0", ["halasz_diag", [4, 0, 0]]], [["8", "0", "-24", "0", "0", "1", "0", "1"], "1", ["halasz_off", [4, 0, 0]]]], [[["0", "0", "-8", "0", "-2", "1", "0", "0"], "0", ["halasz_diag", [4, 0, 1]]], [["6", "-2", "-24", "0", "-6", "1", "0", "1"], "-1", ["halasz_off", [4, 0, 1]]]], [[["0", "0", "-8", "-2", "0", "1", "0", "0"], "0", ["halasz_diag", [4, 1, 0]]], [["8", "2", "-24", "-6", "0", "1", "0", "1"], "1", ["halasz_off", [4, 1, 0]]]], [[["0", "0", "-8", "-2", "-2", "1", "0", "0"], "0", ["halasz_diag", [4, 1, 1]]], [["6", "0", "-24", "-6", "-6", "1", "0", "1"], "-1", ["halasz_off", [4, 1, 1]]]], [[["0", "0", "-8", "-4", "0", "1", "0", "0"], "0", ["halasz_diag", [4, 2, 0]]], [["8", "4", "-24", "-12", "0", "1", "0", "1"], "1", ["halasz_off", [4, 2, 0]]]], [[["0", "0", "-8", "-4", "-2", "1", "0", "0"], "0", ["halasz_diag", [4, 2, 1]]], [["6", "2", "-24", "-12", "-6", "1", "0", "1"], "-1", ["halasz_off", [4, 2, 1]]]], [[["0", "0", "-8", "-6", "0", "1", "0", "0"], "0", ["mean_diag", [4, 3, 0]]], [["4", "3", "-8", "-6", "0", "1", "0", "1"], "1", ["mean_off", [4, 3, 0]]]], [[["0", "0", "-8", "-6", "0", "1", "0", "0"], "0", ["halasz_diag", [4, 3, 0]]], [["8", "6", "-24", "-18", "0", "1", "0", "1"], "1", ["halasz_off", [4, 3, 0]]]], [[["0", "0", "-8", "-6", "-2", "1", "0", "0"], "0", ["halasz_diag", [4, 3, 1]]], [["6", "4", "-24", "-18", "-6", "1", "0", "1"], "-1", ["halasz_off", [4, 3, 1]]]], [[["0", "0", "-8", "-8", "0", "1", "0", "0"], "0", ["mean_diag", [4, 4, 0]]], [["4", "4", "-8", "-8", "0", "1", "0", "1"], "1", ["mean_off", [4, 4, 0]]]], [[["0", "0", "-8", "-8", "0", "1", "0", "0"], "0", ["halasz_diag", [4, 4, 0]]], [["8", "8", "-24", "-24", "0", "1", "0", "1"], "1", ["halasz_off", [4, 4, 0]]]], [[["0", "0", "-8", "-8", "-2", "1", "0", "0"], "0", ["halasz_diag", [4, 4, 1]]], [["6", "6", "-24", "-24", "-6", "1", "0", "1"], "-1", ["halasz_off", [4, 4, 1]]]], [[["0", "0", "-8", "-10", "0", "1", "0", "0"], "0", ["mean_diag", [4, 5, 0]]], [["4", "5", "-8", "-10", "0", "1", "0", "1"], "1", ["mean_off", [4, 5, 0]]]], [[["0", "0", "-8", "-10", "0", "1", "0", "0"], "0", ["halasz_diag", [4, 5, 0]]], [["8", "10", "-24", "-30", "0", "1", "0", "1"], "1", ["halasz_off", [4, 5, 0]]]], [[["0", "0", "-8", "-12", "0", "1", "0", "0"], "0", ["mean_diag", [4, 6, 0]]], [["4", "6", "-8", "-12", "0", "1", "0", "1"], "1", ["mean_off", [4, 6, 0]]]], [[["0", "0", "-8", "-12", "0", "1", "0", "0"], "0", ["halasz_diag", [4, 6, 0]]], [["8", "12", "-24", "-36", "0", "1", "0", "1"], "1", ["halasz_off", [4, 6, 0]]]], [[["0", "0", "-8", "-14", "0", "1", "0", "0"], "0", ["mean_diag", [4, 7, 0]]], [["4", "7", "-8", "-14", "0", "1", "0", "1"], "1", ["mean_off", [4, 7, 0]]]], [[["0", "0", "-8", "-14", "0", "1", "0", "0"], "0", ["halasz_diag", [4, 7, 0]]], [["8", "14", "-24", "-42", "0", "1", "0", "1"], "1", ["halasz_off", [4, 7, 0]]]], [[["0", "0", "-8", "-16", "0", "1", "0", "0"], "0", ["halasz_diag", [4, 8, 0]]], [["8", "16", "-24", "-48", "0", "1", "0", "1"], "1", ["halasz_off", [4, 8, 0]]]], [[["0", "0", "-8", "-18", "0", "1", "0", "0"], "0", ["halasz_diag", [4, 9, 0]]], [["8", "18", "-24", "-54", "0", "1", "0", "1"], "1", ["halasz_off", [4, 9, 0]]]], [[["0", "0", "-8", "-20", "0", "1", "0", "0"], "0", ["halasz_diag", [4, 10, 0]]], [["8", "20", "-24", "-60", "0", "1", "0", "1"], "1", ["halasz_off", [4, 10, 0]]]], [[["0", "0", "-8", "-22", "0", "1", "0", "0"], "0", ["halasz_diag", [4, 11, 0]]], [["8", "22", "-24", "-66", "0", "1", "0", "1"], "1", ["halasz_off", [4, 11, 0]]]], [[["0", "0", "-8", "-24", "0", "1", "0", "0"], "0", ["halasz_diag", [4, 12, 0]]], [["8", "24", "-24", "-72", "0", "1", "0", "1"], "1", ["halasz_off", [4, 12, 0]]]], [[["0", "0", "-10", "0", "0", "1", "0", "0"], "0", ["halasz_diag", [5, 0, 0]]], [["10", "0", "-30", "0", "0", "1", "0", "1"], "1", ["halasz_off", [5, 0, 0]]]], [[["0", "0", "-10", "0", "-2", "1", "0", "0"], "0", ["halasz_diag", [5, 0, 1]]], [["8", "-2", "-30", "0", "-6", "1", "0", "1"], "-1", ["halasz_off", [5, 0, 1]]]], [[["0", "0", "-10", "-2", "0", "1", "0", "0"], "0", ["halasz_diag", [5, 1, 0]]], [["10", "2", "-30", "-6", "0", "1", "0", "1"], "1", ["halasz_off", [5, 1, 0]]]], [[["0", "0", "-10", "-2", "-2", "1", "0", "0"], "0", ["halasz_diag", [5, 1, 1]]], [["8", "0", "-30", "-6", "-6", "1", "0", "1"], "-1", ["halasz_off", [5, 1, 1]]]], [[["0", "0", "-10", "-4", "0", "1", "0", "0"], "0", ["mean_diag", [5, 2, 0]]], [["5", "2", "-10", "-4", "0", "1", "0", "1"], "1", ["mean_off", [5, 2, 0]]]], [[["0", "0", "-10", "-4", "0", "1", "0", "0"], "0", ["halasz_diag", [5, 2, 0]]], [["10", "4", "-30", "-12", "0", "1", "0", "1"], "1", ["halasz_off", [5, 2, 0]]]], [[["0", "0", "-10", "-4", "-2", "1", "0", "0"], "0", ["halasz_diag", [5, 2, 1]]], [["8", "2", "-30", "-12", "-6", "1", "0", "1"], "-1", ["halasz_off", [5, 2, 1]]]], [[["0", "0", "-10", "-6", "0", "1", "0", "0"], "0", ["mean_diag", [5, 3, 0]]], [["5", "3", "-10", "-6", "0", "1", "0", "1"], "1", ["mean_off", [5, 3, 0]]]], [[["0", "0", "-10", "-6", "0", "1", "0", "0"], "0", ["halasz_diag", [5, 3, 0]]], [["10", "6", "-30", "-18", "0", "1", "0", "1"], "1", ["halasz_off", [5, 3, 0]]]], [[["0", "0", "-10", "-6", "-2", "1", "0", "0"], "0", ["halasz_diag", [5, 3, 1]]], [["8", "4", "-30", "-18", "-6", "1", "0", "1"], "-1", ["halasz_off", [5, 3, 1]]]], [[["0", "0", "-10", "-8", "0", "1", "0", "0"], "0", ["mean_diag", [5, 4, 0]]], [["5", "4", "-10", "-8", "0", "1", "0", "1"], "1", ["mean_off", [5, 4, 0]]]], [[["0", "0", "-10", "-8", "0", "1", "0", "0"], "0", ["halasz_diag", [5, 4, 0]]], [["10", "8", "-30", "-24", "0", "1", "0", "1"], "1", ["halasz_off", [5, 4, 0]]]], [[["0", "0", "-10", "-10", "0", "1", "0", "0"], "0", ["mean_diag", [5, 5, 0]]], [["5", "5", "-10", "-10", "0", "1", "0", "1"], "1", ["mean_off", [5, 5, 0]]]], [[["0", "0", "-10", "-10", "0", "1", "0", "0"], "0", ["halasz_diag", [5, 5, 0]]], [["10", "10", "-30", "-30", "0", "1", "0", "1"], "1", ["halasz_off", [5, 5, 0]]]], [[["0", "0", "-10", "-12", "0", "1", "0", "0"], "0", ["mean_diag", [5, 6, 0]]], [["5", "6", "-10", "-12", "0", "1", "0", "1"], "1", ["mean_off", [5, 6, 0]]]], [[["0", "0", "-10", "-12", "0", "1", "0", "0"], "0", ["halasz_diag", [5, 6, 0]]], [["10", "12", "-30", "-36", "0", "1", "0", "1"], "1", ["halasz_off", [5, 6, 0]]]], [[["0", "0", "-10", "-14", "0", "1", "0", "0"], "0", ["halasz_diag", [5, 7, 0]]], [["10", "14", "-30", "-42", "0", "1", "0", "1"], "1", ["halasz_off", [5, 7, 0]]]], [[["0", "0", "-10", "-16", "0", "1", "0", "0"], "0", ["halasz_diag", [5, 8, 0]]], [["10", "16", "-30", "-48", "0", "1", "0", "1"], "1", ["halasz_off", [5, 8, 0]]]], [[["0", "0", "-10", "-18", "0", "1", "0", "0"], "0", ["halasz_diag", [5, 9, 0]]], [["10", "18", "-30", "-54", "0", "1", "0", "1"], "1", ["halasz_off", [5, 9, 0]]]], [[["0", "0", "-10", "-20", "0", "1", "0", "0"], "0", ["halasz_diag", [5, 10, 0]]], [["10", "20", "-30", "-60", "0", "1", "0", "1"], "1", ["halasz_off", [5, 10, 0]]]], [[["0", "0", "-10", "-22", "0", "1", "0", "0"], "0", ["halasz_diag", [5, 11, 0]]], [["10", "22", "-30", "-66", "0", "1", "0", "1"], "1", ["halasz_off", [5, 11, 0]]]], [[["0", "0", "-10", "-24", "0", "1", "0", "0"], "0", ["halasz_diag", [5, 12, 0]]], [["10", "24", "-30", "-72", "0", "1", "0", "1"], "1", ["halasz_off", [5, 12, 0]]]], [[["0", "0", "-12", "0", "0", "1", "0", "0"], "0", ["mean_diag", [6, 0, 0]]], [["6", "0", "-12", "0", "0", "1", "0", "1"], "1", ["mean_off", [6, 0, 0]]]], [[["0", "0", "-12", "0", "0", "1", "0", "0"], "0", ["halasz_diag", [6, 0, 0]]], [["12", "0", "-36", "0", "0", "1", "0", "1"], "1", ["halasz_off", [6, 0, 0]]]], [[["0", "0", "-12", "0", "-2", "1", "0", "0"], "0", ["halasz_diag", [6, 0, 1]]], [["10", "-2", "-36", "0", "-6", "1", "0", "1"], "-1", ["halasz_off", [6, 0, 1]]]], [[["0", "0", "-12", "-2", "0", "1", "0", "0"], "0", ["mean_diag", [6, 1, 0]]], [["6", "1", "-12", "-2", "0", "1", "0", "1"], "1", ["mean_off", [6, 1, 0]]]], [[["0", "0", "-12", "-2", "0", "1", "0", "0"], "0", ["halasz_diag", [6, 1, 0]]], [["12", "2", "-36", "-6", "0", "1", "0", "1"], "1", ["halasz_off", [6, 1, 0]]]], [[["0", "0", "-12", "-2", "-2", "1", "0", "0"], "0", ["halasz_diag", [6, 1, 1]]], [["10", "0", "-36", "-6", "-6", "1", "0", "1"], "-1", ["halasz_off", [6, 1, 1]]]], [[["0", "0", "-12", "-4", "0", "1", "0", "0"], "0", ["mean_diag", [6, 2, 0]]], [["6", "2", "-12", "-4", "0", "1", "0", "1"], "1", ["mean_off", [6, 2, 0]]]], [[["0", "0", "-12", "-4", "0", "1", "0", "0"], "0", ["halasz_diag", [6, 2, 0]]], [["12", "4", "-36", "-12", "0", "1", "0", "1"], "1", ["halasz_off", [6, 2, 0]]]], [[["0", "0", "-12", "-4", "-2", "1", "0", "0"], "0", ["halasz_diag", [6, 2, 1]]], [["10", "2", "-36", "-12", "-6", "1", "0", "1"], "-1", ["halasz_off", [6, 2, 1]]]], [[["0", "0", "-12", "-6", "0", "1", "0", "0"], "0", ["mean_diag", [6, 3, 0]]], [["6", "3", "-12", "-6", "0", "1", "0", "1"], "1", ["mean_off", [6, 3, 0]]]], [[["0", "0", "-12", "-6", "0", "1", "0", "0"], "0", ["halasz_diag", [6, 3, 0]]], [["12", "6", "-36", "-18", "0", "1", "0", "1"], "1", ["halasz_off", [6, 3, 0]]]], [[["0", "0", "-12", "-8", "0", "1", "0", "0"], "0", ["mean_diag", [6, 4, 0]]], [["6", "4", "-12", "-8", "0", "1", "0", "1"], "1", ["mean_off", [6, 4, 0]]]], [[["0", "0", "-12", "-8", "0", "1", "0", "0"], "0", ["halasz_diag", [6, 4, 0]]], [["12", "8", "-36", "-24", "0", "1", "0", "1"], "1", ["halasz_off", [6, 4, 0]]]], [[["0", "0", "-12", "-10", "0", "1", "0", "0"], "0", ["mean_diag", [6, 5, 0]]], [["6", "5", "-12", "-10", "0", "1", "0", "1"], "1", ["mean_off", [6, 5, 0]]]], [[["0", "0", "-12", "-10", "0", "1", "0", "0"], "0", ["halasz_diag", [6, 5, 0]]], [["12", "10", "-36", "-30", "0", "1", "0", "1"], "1", ["halasz_off", [6, 5, 0]]]], [[["0", "0", "-12", "-12", "0", "1", "0", "0"], "0", ["halasz_diag", [6, 6, 0]]], [["12", "12", "-36", "-36", "0", "1", "0", "1"], "1", ["halasz_off", [6, 6, 0]]]], [[["0", "0", "-12", "-14", "0", "1", "0", "0"], "0", ["halasz_diag", [6, 7, 0]]], [["12", "14", "-36", "-42", "0", "1", "0", "1"], "1", ["halasz_off", [6, 7, 0]]]], [[["0", "0", "-12", "-16", "0", "1", "0", "0"], "0", ["halasz_diag", [6, 8, 0]]], [["12", "16", "-36", "-48", "0", "1", "0", "1"], "1", ["halasz_off", [6, 8, 0]]]], [[["0", "0", "-12", "-18", "0", "1", "0", "0"], "0", ["halasz_diag", [6, 9, 0]]], [["12", "18", "-36", "-54", "0", "1", "0", "1"], "1", ["halasz_off", [6, 9, 0]]]], [[["0", "0", "-12", "-20", "0", "1", "0", "0"], "0", ["halasz_diag", [6, 10, 0]]], [["12", "20", "-36", "-60", "0", "1", "0", "1"], "1", ["halasz_off", [6, 10, 0]]]], [[["0", "0", "-12", "-22", "0", "1", "0", "0"], "0", ["halasz_diag", [6, 11, 0]]], [["12", "22", "-36", "-66", "0", "1", "0", "1"], "1", ["halasz_off", [6, 11, 0]]]], [[["0", "0", "-12", "-24", "0", "1", "0", "0"], "0", ["halasz_diag", [6, 12, 0]]], [["12", "24", "-36", "-72", "0", "1", "0", "1"], "1", ["halasz_off", [6, 12, 0]]]], [[["0", "0", "-14", "0", "0", "1", "0", "0"], "0", ["mean_diag", [7, 0, 0]]], [["7", "0", "-14", "0", "0", "1", "0", "1"], "1", ["mean_off", [7, 0, 0]]]], [[["0", "0", "-14", "0", "0", "1", "0", "0"], "0", ["halasz_diag", [7, 0, 0]]], [["14", "0", "-42", "0", "0", "1", "0", "1"], "1", ["halasz_off", [7, 0, 0]]]], [[["0", "0", "-14", "0", "-2", "1", "0", "0"], "0", ["halasz_diag", [7, 0, 1]]], [["12", "-2", "-42", "0", "-6", "1", "0", "1"], "-1", ["halasz_off", [7, 0, 1]]]], [[["0", "0", "-14", "-2", "0", "1", "0", "0"], "0", ["mean_diag", [7, 1, 0]]], [["7", "1", "-14", "-2", "0", "1", "0", "1"], "1", ["mean_off", [7, 1, 0]]]], [[["0", "0", "-14", "-2", "0", "1", "0", "0"], "0", ["halasz_diag", [7, 1, 0]]], [["14", "2", "-42", "-6", "0", "1", "0", "1"], "1", ["halasz_off", [7, 1, 0]]]], [[["0", "0", "-14", "-2", "-2", "1", "0", "0"], "0", ["halasz_diag", [7, 1, 1]]], [["12", "0", "-42", "-6", "-6", "1", "0", "1"], "-1", ["halasz_off", [7, 1, 1]]]], [[["0", "0", "-14", "-4", "0", "1", "0", "0"], "0", ["mean_diag", [7, 2, 0]]], [["7", "2", "-14", "-4", "0", "1", "0", "1"], "1", ["mean_off", [7, 2, 0]]]], [[["0", "0", "-14", "-4", "0", "1", "0", "0"], "0", ["halasz_diag", [7, 2, 0]]], [["14", "4", "-42", "-12", "0", "1", "0", "1"], "1", ["halasz_off", [7, 2, 0]]]], [[["0", "0", "-14", "-6", "0", "1", "0", "0"], "0", ["mean_diag", [7, 3, 0]]], [["7", "3", "-14", "-6", "0", "1", "0", "1"], "1", ["mean_off", [7, 3, 0]]]], [[["0", "0", "-14", "-6", "0", "1", "0", "0"], "0", ["halasz_diag", [7, 3, 0]]], [["14", "6", "-42", "-18", "0", "1", "0", "1"], "1", ["halasz_off", [7, 3, 0]]]], [[["0", "0", "-14", "-8", "0", "1", "0", "0"], "0", ["mean_diag", [7, 4, 0]]], [["7", "4", "-14", "-8", "0", "1", "0", "1"], "1", ["mean_off", [7, 4, 0]]]], [[["0", "0", "-14", "-8", "0", "1", "0", "0"], "0", ["halasz_diag", [7, 4, 0]]], [["14", "8", "-42", "-24", "0", "1", "0", "1"], "1", ["halasz_off", [7, 4, 0]]]], [[["0", "0", "-14", "-10", "0", "1", "0", "0"], "0", ["halasz_diag", [7, 5, 0]]], [["14", "10", "-42", "-30", "0", "1", "0", "1"], "1", ["halasz_off", [7, 5, 0]]]], [[["0", "0", "-14", "-12", "0", "1", "0", "0"], "0", ["halasz_diag", [7, 6, 0]]], [["14", "12", "-42", "-36", "0", "1", "0", "1"], "1", ["halasz_off", [7, 6, 0]]]], [[["0", "0", "-14", "-14", "0", "1", "0", "0"], "0", ["halasz_diag", [7, 7, 0]]], [["14", "14", "-42", "-42", "0", "1", "0", "1"], "1", ["halasz_off", [7, 7, 0]]]], [[["0", "0", "-14", "-16", "0", "1", "0", "0"], "0", ["halasz_diag", [7, 8, 0]]], [["14", "16", "-42", "-48", "0", "1", "0", "1"], "1", ["halasz_off", [7, 8, 0]]]], [[["0", "0", "-14", "-18", "0", "1", "0", "0"], "0", ["halasz_diag", [7, 9, 0]]], [["14", "18", "-42", "-54", "0", "1", "0", "1"], "1", ["halasz_off", [7, 9, 0]]]], [[["0", "0", "-14", "-20", "0", "1", "0", "0"], "0", ["halasz_diag", [7, 10, 0]]], [["14", "20", "-42", "-60", "0", "1", "0", "1"], "1", ["halasz_off", [7, 10, 0]]]], [[["0", "0", "-14", "-22", "0", "1", "0", "0"], "0", ["halasz_diag", [7, 11, 0]]], [["14", "22", "-42", "-66", "0", "1", "0", "1"], "1", ["halasz_off", [7, 11, 0]]]], [[["0", "0", "-16", "0", "0", "1", "0", "0"], "0", ["mean_diag", [8, 0, 0]]], [["8", "0", "-16", "0", "0", "1", "0", "1"], "1", ["mean_off", [8, 0, 0]]]], [[["0", "0", "-16", "0", "0", "1", "0", "0"], "0", ["halasz_diag", [8, 0, 0]]], [["16", "0", "-48", "0", "0", "1", "0", "1"], "1", ["halasz_off", [8, 0, 0]]]], [[["0", "0", "-16", "0", "-2", "1", "0", "0"], "0", ["halasz_diag", [8, 0, 1]]], [["14", "-2", "-48", "0", "-6", "1", "0", "1"], "-1", ["halasz_off", [8, 0, 1]]]], [[["0", "0", "-16", "-2", "0", "1", "0", "0"], "0", ["mean_diag", [8, 1, 0]]], [["8", "1", "-16", "-2", "0", "1", "0", "1"], "1", ["mean_off", [8, 1, 0]]]], [[["0", "0", "-16", "-2", "0", "1", "0", "0"], "0", ["halasz_diag", [8, 1, 0]]], [["16", "2", "-48", "-6", "0", "1", "0", "1"], "1", ["halasz_off", [8, 1, 0]]]], [[["0", "0", "-16", "-4", "0", "1", "0", "0"], "0", ["mean_diag", [8, 2, 0]]], [["8", "2", "-16", "-4", "0", "1", "0", "1"], "1", ["mean_off", [8, 2, 0]]]], [[["0", "0", "-16", "-4", "0", "1", "0", "0"], "0", ["halasz_diag", [8, 2, 0]]], [["16", "4", "-48", "-12", "0", "1", "0", "1"], "1", ["halasz_off", [8, 2, 0]]]], [[["0", "0", "-16", "-6", "0", "1", "0", "0"], "0", ["mean_diag", [8, 3, 0]]], [["8", "3", "-16", "-6", "0", "1", "0", "1"], "1", ["mean_off", [8, 3, 0]]]], [[["0", "0", "-16", "-6", "0", "1", "0", "0"], "0", ["halasz_diag", [8, 3, 0]]], [["16", "6", "-48", "-18", "0", "1", "0", "1"], "1", ["halasz_off", [8, 3, 0]]]], [[["0", "0", "-16", "-8", "0", "1", "0", "0"], "0", ["halasz_diag", [8, 4, 0]]], [["16", "8", "-48", "-24", "0", "1", "0", "1"], "1", ["halasz_off", [8, 4, 0]]]], [[["0", "0", "-16", "-10", "0", "1", "0", "0"], "0", ["halasz_diag", [8, 5, 0]]], [["16", "10", "-48", "-30", "0", "1", "0", "1"], "1", ["halasz_off", [8, 5, 0]]]], [[["0", "0", "-16", "-12", "0", "1", "0", "0"], "0", ["halasz_diag", [8, 6, 0]]], [["16", "12", "-48", "-36", "0", "1", "0", "1"], "1", ["halasz_off", [8, 6, 0]]]], [[["0", "0", "-16", "-14", "0", "1", "0", "0"], "0", ["halasz_diag", [8, 7, 0]]], [["16", "14", "-48", "-42", "0", "1", "0", "1"], "1", ["halasz_off", [8, 7, 0]]]], [[["0", "0", "-16", "-16", "0", "1", "0", "0"], "0", ["halasz_diag", [8, 8, 0]]], [["16", "16", "-48", "-48", "0", "1", "0", "1"], "1", ["halasz_off", [8, 8, 0]]]], [[["0", "0", "-16", "-18", "0", "1", "0", "0"], "0", ["halasz_diag", [8, 9, 0]]], [["16", "18", "-48", "-54", "0", "1", "0", "1"], "1", ["halasz_off", [8, 9, 0]]]], [[["0", "0", "-16", "-20", "0", "1", "0", "0"], "0", ["halasz_diag", [8, 10, 0]]], [["16", "20", "-48", "-60", "0", "1", "0", "1"], "1", ["halasz_off", [8, 10, 0]]]], [[["0", "0", "-18", "0", "0", "1", "0", "0"], "0", ["mean_diag", [9, 0, 0]]], [["9", "0", "-18", "0", "0", "1", "0", "1"], "1", ["mean_off", [9, 0, 0]]]], [[["0", "0", "-18", "0", "0", "1", "0", "0"], "0", ["halasz_diag", [9, 0, 0]]], [["18", "0", "-54", "0", "0", "1", "0", "1"], "1", ["halasz_off", [9, 0, 0]]]], [[["0", "0", "-18", "-2", "0", "1", "0", "0"], "0", ["mean_diag", [9, 1, 0]]], [["9", "1", "-18", "-2", "0", "1", "0", "1"], "1", ["mean_off", [9, 1, 0]]]], [[["0", "0", "-18", "-2", "0", "1", "0", "0"], "0", ["halasz_diag", [9, 1, 0]]], [["18", "2", "-54", "-6", "0", "1", "0", "1"], "1", ["halasz_off", [9, 1, 0]]]], [[["0", "0", "-18", "-4", "0", "1", "0", "0"], "0", ["mean_diag", [9, 2, 0]]], [["9", "2", "-18", "-4", "0", "1", "0", "1"], "1", ["mean_off", [9, 2, 0]]]], [[["0", "0", "-18", "-4", "0", "1", "0", "0"], "0", ["halasz_diag", [9, 2, 0]]], [["18", "4", "-54", "-12", "0", "1", "0", "1"], "1", ["halasz_off", [9, 2, 0]]]], [[["0", "0", "-18", "-6", "0", "1", "0", "0"], "0", ["halasz_diag", [9, 3, 0]]], [["18", "6", "-54", "-18", "0", "1", "0", "1"], "1", ["halasz_off", [9, 3, 0]]]], [[["0", "0", "-18", "-8", "0", "1", "0", "0"], "0", ["halasz_diag", [9, 4, 0]]], [["18", "8", "-54", "-24", "0", "1", "0", "1"], "1", ["halasz_off", [9, 4, 0]]]], [[["0", "0", "-18", "-10", "0", "1", "0", "0"], "0", ["halasz_diag", [9, 5, 0]]], [["18", "10", "-54", "-30", "0", "1", "0", "1"], "1", ["halasz_off", [9, 5, 0]]]], [[["0", "0", "-18", "-12", "0", "1", "0", "0"], "0", ["halasz_diag", [9, 6, 0]]], [["18", "12", "-54", "-36", "0", "1", "0", "1"], "1", ["halasz_off", [9, 6, 0]]]], [[["0", "0", "-18", "-14", "0", "1", "0", "0"], "0", ["halasz_diag", [9, 7, 0]]], [["18", "14", "-54", "-42", "0", "1", "0", "1"], "1", ["halasz_off", [9, 7, 0]]]], [[["0", "0", "-18", "-16", "0", "1", "0", "0"], "0", ["halasz_diag", [9, 8, 0]]], [["18", "16", "-54", "-48", "0", "1", "0", "1"], "1", ["halasz_off", [9, 8, 0]]]], [[["0", "0", "-18", "-18", "0", "1", "0", "0"], "0", ["halasz_diag", [9, 9, 0]]], [["18", "18", "-54", "-54", "0", "1", "0", "1"], "1", ["halasz_off", [9, 9, 0]]]], [[["0", "0", "-20", "0", "0", "1", "0", "0"], "0", ["mean_diag", [10, 0, 0]]], [["10", "0", "-20", "0", "0", "1", "0", "1"], "1", ["mean_off", [10, 0, 0]]]], [[["0", "0", "-20", "0", "0", "1", "0", "0"], "0", ["halasz_diag", [10, 0, 0]]], [["20", "0", "-60", "0", "0", "1", "0", "1"], "1", ["halasz_off", [10, 0, 0]]]], [[["0", "0", "-20", "-2", "0", "1", "0", "0"], "0", ["mean_diag", [10, 1, 0]]], [["10", "1", "-20", "-2", "0", "1", "0", "1"], "1", ["mean_off", [10, 1, 0]]]], [[["0", "0", "-20", "-2", "0", "1", "0", "0"], "0", ["halasz_diag", [10, 1, 0]]], [["20", "2", "-60", "-6", "0", "1", "0", "1"], "1", ["halasz_off", [10, 1, 0]]]], [[["0", "0", "-20", "-4", "0", "1", "0", "0"], "0", ["halasz_diag", [10, 2, 0]]], [["20", "4", "-60", "-12", "0", "1", "0", "1"], "1", ["halasz_off", [10, 2, 0]]]], [[["0", "0", "-20", "-6", "0", "1", "0", "0"], "0", ["halasz_diag", [10, 3, 0]]], [["20", "6", "-60", "-18", "0", "1", "0", "1"], "1", ["halasz_off", [10, 3, 0]]]], [[["0", "0", "-20", "-8", "0", "1", "0", "0"], "0", ["halasz_diag", [10, 4, 0]]], [["20", "8", "-60", "-24", "0", "1", "0", "1"], "1", ["halasz_off", [10, 4, 0]]]], [[["0", "0", "-20", "-10", "0", "1", "0", "0"], "0", ["halasz_diag", [10, 5, 0]]], [["20", "10", "-60", "-30", "0", "1", "0", "1"], "1", ["halasz_off", [10, 5, 0]]]], [[["0", "0", "-20", "-12", "0", "1", "0", "0"], "0", ["halasz_diag", [10, 6, 0]]], [["20", "12", "-60", "-36", "0", "1", "0", "1"], "1", ["halasz_off", [10, 6, 0]]]], [[["0", "0", "-20", "-14", "0", "1", "0", "0"], "0", ["halasz_diag", [10, 7, 0]]], [["20", "14", "-60", "-42", "0", "1", "0", "1"], "1", ["halasz_off", [10, 7, 0]]]], [[["0", "0", "-20", "-16", "0", "1", "0", "0"], "0", ["halasz_diag", [10, 8, 0]]], [["20", "16", "-60", "-48", "0", "1", "0", "1"], "1", ["halasz_off", [10, 8, 0]]]], [[["0", "0", "-22", "0", "0", "1", "0", "0"], "0", ["mean_diag", [11, 0, 0]]], [["11", "0", "-22", "0", "0", "1", "0", "1"], "1", ["mean_off", [11, 0, 0]]]], [[["0", "0", "-22", "0", "0", "1", "0", "0"], "0", ["halasz_diag", [11, 0, 0]]], [["22", "0", "-66", "0", "0", "1", "0", "1"], "1", ["halasz_off", [11, 0, 0]]]], [[["0", "0", "-22", "-2", "0", "1", "0", "0"], "0", ["halasz_diag", [11, 1, 0]]], [["22", "2", "-66", "-6", "0", "1", "0", "1"], "1", ["halasz_off", [11, 1, 0]]]], [[["0", "0", "-22", "-4", "0", "1", "0", "0"], "0", ["halasz_diag", [11, 2, 0]]], [["22", "4", "-66", "-12", "0", "1", "0", "1"], "1", ["halasz_off", [11, 2, 0]]]], [[["0", "0", "-22", "-6", "0", "1", "0", "0"], "0", ["halasz_diag", [11, 3, 0]]], [["22", "6", "-66", "-18", "0", "1", "0", "1"], "1", ["halasz_off", [11, 3, 0]]]], [[["0", "0", "-22", "-8", "0", "1", "0", "0"], "0", ["halasz_diag", [11, 4, 0]]], [["22", "8", "-66", "-24", "0", "1", "0", "1"], "1", ["halasz_off", [11, 4, 0]]]], [[["0", "0", "-22", "-10", "0", "1", "0", "0"], "0", ["halasz_diag", [11, 5, 0]]], [["22", "10", "-66", "-30", "0", "1", "0", "1"], "1", ["halasz_off", [11, 5, 0]]]], [[["0", "0", "-22", "-12", "0", "1", "0", "0"], "0", ["halasz_diag", [11, 6, 0]]], [["22", "12", "-66", "-36", "0", "1", "0", "1"], "1", ["halasz_off", [11, 6, 0]]]], [[["0", "0", "-22", "-14", "0", "1", "0", "0"], "0", ["halasz_diag", [11, 7, 0]]], [["22", "14", "-66", "-42", "0", "1", "0", "1"], "1", ["halasz_off", [11, 7, 0]]]], [[["0", "0", "-24", "0", "0", "1", "0", "0"], "0", ["halasz_diag", [12, 0, 0]]], [["24", "0", "-72", "0", "0", "1", "0", "1"], "1", ["halasz_off", [12, 0, 0]]]], [[["0", "0", "-24", "-2", "0", "1", "0", "0"], "0", ["halasz_diag", [12, 1, 0]]], [["24", "2", "-72", "-6", "0", "1", "0", "1"], "1", ["halasz_off", [12, 1, 0]]]], [[["0", "0", "-24", "-4", "0", "1", "0", "0"], "0", ["halasz_diag", [12, 2, 0]]], [["24", "4", "-72", "-12", "0", "1", "0", "1"], "1", ["halasz_off", [12, 2, 0]]]], [[["0", "0", "-24", "-6", "0", "1", "0", "0"], "0", ["halasz_diag", [12, 3, 0]]], [["24", "6", "-72", "-18", "0", "1", "0", "1"], "1", ["halasz_off", [12, 3, 0]]]], [[["0", "0", "-24", "-8", "0", "1", "0", "0"], "0", ["halasz_diag", [12, 4, 0]]], [["24", "8", "-72", "-24", "0", "1", "0", "1"], "1", ["halasz_off", [12, 4, 0]]]], [[["0", "0", "-24", "-10", "0", "1", "0", "0"], "0", ["halasz_diag", [12, 5, 0]]], [["24", "10", "-72", "-30", "0", "1", "0", "1"], "1", ["halasz_off", [12, 5, 0]]]], [[["0", "0", "-24", "-12", "0", "1", "0", "0"], "0", ["halasz_diag", [12, 6, 0]]], [["24", "12", "-72", "-36", "0", "1", "0", "1"], "1", ["halasz_off", [12, 6, 0]]]]], "target": ["0", "0", "-2", "-2", "-2", "1", "1/100", "0"], "leaves": [{"path": [[0, 1], [53, 1], [116, 1], [15, 1], [118, 1], [2, 1], [58, 1], [139, 1], [25, 1], [18, 1], [121, 1], [4, 1], [3, 1]], "dual": {"bound": "-1/1750", "weights": [[6, "27/175"], [8, "41/120"], [22, "23/140"], [43, "1/840"], [46, "13/28"], [183, "23/840"], [184, "1/70"], [188, "69/140"]]}}, {"path": [[0, 1], [53, 1], [116, 1], [15, 1], [118, 1], [2, 1], [58, 1], [139, 1], [25, 1], [18, 1], [121, 1], [4, 1], [3, 0], [180, 1], [196, 1], [209, 1], [220, 1]], "dual": {"bound": "-1/2445", "weights": [[6, "12661/48900"], [8, "95/1956"], [21, "263/978"], [43, "5/489"], [46, "183/326"], [187, "23/1956"], [188, "263/652"], [192, "13/978"]]}}, {"path": [[0, 1], [53, 1], [116, 1], [15, 1], [118, 1], [2, 1], [58, 1], [139, 1], [25, 1], [18, 1], [121, 1], [4, 1], [3, 0], [180, 1], [196, 1], [209, 1], [220, 0]], "dual": {"bound": "-1/3825", "weights": [[6, "3497/15300"], [8, "29/612"], [21, "73/306"], [46, "1/2"], [187, "7/612"], [188, "95/204"], [191, "2/153"], [192, "1/102"]]}}, {"path": [[0, 1], [53, 1], [116, 1], [15, 1], [118, 1], [2, 1], [58, 1], [139, 1], [25, 1], [18, 1], [121, 1], [4, 1], [3, 0], [180, 1], [196, 1], [209, 0]], "dual": {"bound": "-1/9425", "weights": [[6, "8573/37700"], [8, "79/1508"], [21, "179/754"], [46, "1/2"], [187, "19/1508"], [188, "697/1508"], [190, "11/754"], [191, "4/377"]]}}, {"path": [[0, 1], [53, 1], [116, 1], [15, 1], [118, 1], [2, 1], [58, 1], [139, 1], [25, 1], [18, 1], [121, 1], [4, 1], [3, 0], [180, 1], [196, 0]], "dual": {"bound": "0", "weights": [[21, "1/100"], [46, "3/160"], [49, "767/800"], [184, "1/1600"], [188, "33/1600"], [189, "1/1600"], [190, "1/1600"]]}}, {"path": [[0, 1], [53, 1], [116, 1], [15, 1], [118, 1], [2, 1], [58, 1], [139, 1], [25, 1], [18, 1], [121, 1], [4, 1], [3, 0], [180, 0], [161, 1]], "dual": {"bound": "0", "weights": [[21, "1/100"], [46, "13/700"], [49, "671/700"], [184, "1/1400"], [188, "29/1400"], [189, "1/1400"], [190, "1/1400"]]}}, {"path": [[0, 1], [53, 1], [116, 1], [15, 1], [118, 1], [2, 1], [58, 1], [139, 1], [25, 1], [18, 1], [121, 1], [4, 1], [3, 0], [180, 0], [161, 0]], "dual": {"bound": "0", "weights": [[21, "1/100"], [46, "11/600"], [49, "23/24"], [183, "1/1200"], [184, "1/1200"], [188, "1/48"], [190, "1/1200"]]}}, {"path": [[0, 1], [53, 1], [116, 1], [15, 1], [118, 1], [2, 1], [58, 1], [139, 1], [25, 1], [18, 1], [121, 1], [4, 0], [180, 1], [196, 1], [209, 1], [220, 1]], "dual": {"bound": "0", "weights": [[21, "1/100"], [43, "1/2200"], [46, "21/1100"], [49, "1061/1100"], [184, "1/2200"], [187, "3/200"], [191, "1/2200"]]}}, {"path": [[0, 1], [53, 1], [116, 1], [15, 1], [118, 1], [2, 1], [58, 1], [139, 1], [25, 1], [18, 1], [121, 1], [4, 0], [180, 1], [196, 1], [209, 1], [220, 0]], "dual": {"bound": "0", "weights": [[6, "959/4100"], [21, "10/41"], [46, "19/41"], [184, "1/82"], [187, "1/2"], [190, "1/82"], [191, "1/82"]]}}, {"path": [[0, 1], [53, 1], [116, 1], [15, 1], [118, 1], [2, 1], [58, 1], [139, 1], [25, 1], [18, 1], [121, 1], [4, 0], [180, 1], [196, 1], [209, 0]], "dual": {"bound": "0", "weights": [[21, "1/100"], [46, "17/900"], [49, "863/900"], [184, "1/1800"], [187, "37/1800"], [189, "1/1800"], [190, "1/1800"]]}}, {"path": [[0, 1], [53, 1], [116, 1], [15, 1], [118, 1], [2, 1], [58, 1], [139, 1], [25, 1], [18, 1], [121, 1], [4, 0], [180, 1], [196, 0]], "dual": {"bound": "0", "weights": [[21, "1/100"], [46, "3/160"], [49, "767/800"], [184, "1/1600"], [187, "33/1600"], [188, "1/1600"], [189, "1/1600"]]}}, {"path": [[0, 1], [53, 1], [116, 1], [15, 1], [118, 1], [2, 1], [58, 1], [139, 1], [25, 1], [18, 1], [121, 1], [4, 0], [180, 0], [161, 1]], "dual": {"bound": "0", "weights": [[21, "1/100"], [46, "13/700"], [49, "671/700"], [184, "1/1400"], [187, "29/1400"], [188, "1/1400"], [189, "1/1400"]]}}, {"path": [[0, 1], [53, 1], [116, 1], [15, 1], [118, 1], [2, 1], [58, 1], [139, 1], [25, 1], [18, 1], [121, 1], [4, 0], [180, 0], [161, 0]], "dual": {"bound": "0", "weights": [[21, "1/100"], [46, "11/600"], [49, "23/24"], [183, "1/1200"], [184, "1/1200"], [187, "1/48"], [189, "1/1200"]]}}, {"path": [[0, 1], [53, 1], [116, 1], [15, 1], [118, 1], [2, 1], [58, 1], [139, 1], [25, 1], [18, 1], [121, 0], [4, 1], [3, 1]], "dual": {"bound": "-1/1500", "weights": [[6, "109/600"], [8, "9/40"], [22, "23/120"], [46, "11/24"], [183, "1/40"], [184, "1/60"], [186, "1/120"], [188, "59/120"]]}}, {"path": [[0, 1], [53, 1], [116, 1], [15, 1], [118, 1], [2, 1], [58, 1], [139, 1], [25, 1], [18, 1], [121, 0], [4, 1], [3, 0]], "dual": {"bound": "0", "weights": [[21, "1/100"], [46, "11/600"], [49, "24/25"], [183, "1/1200"], [184, "1/1200"], [186, "1/1200"], [188, "23/1200"]]}}, {"path": [[0, 1], [53, 1], [116, 1], [15, 1], [118, 1], [2, 1], [58, 1], [139, 1], [25, 1], [18, 1], [121, 0], [4, 0]], "dual": {"bound": "0", "weights": [[21, "1/100"], [46, "11/600"], [49, "24/25"], [183, "1/1200"], [184, "1/1200"], [186, "1/1200"], [187, "23/1200"]]}}, {"path": [[0, 1], [53, 1], [116, 1], [15, 1], [118, 1], [2, 1], [58, 1], [139, 1], [25, 1], [18, 0], [3, 1]], "dual": {"bound": "-4/5175", "weights": [[6, "3193/20700"], [8, "212/621"], [22, "34/207"], [46, "32/69"], [183, "17/621"], [184, "1/69"], [185, "1/621"], [186, "34/69"]]}}, {"path": [[0, 1], [53, 1], [116, 1], [15, 1], [118, 1], [2, 1], [58, 1], [139, 1], [25, 1], [18, 0], [3, 0], [161, 1], [180, 1]], "dual": {"bound": "0", "weights": [[21, "1/100"], [46, "3/160"], [49, "1543/1600"], [184, "1/1600"], [185, "1/1600"], [186, "3/200"], [188, "1/1600"]]}}, {"path": [[0, 1], [53, 1], [116, 1], [15, 1], [118, 1], [2, 1], [58, 1], [139, 1], [25, 1], [18, 0], [3, 0], [161, 1], [180, 0]], "dual": {"bound": "0", "weights": [[21, "1/100"], [46, "13/700"], [49, "671/700"], [184, "1/1400"], [186, "29/1400"], [187, "1/1400"], [188, "1/1400"]]}}, {"path": [[0, 1], [53, 1], [116, 1], [15, 1], [118, 1], [2, 1], [58, 1], [139, 1], [25, 1], [18, 0], [3, 0], [161, 0]], "dual": {"bound": "0", "weights": [[21, "1/100"], [46, "11/600"], [49, "23/24"], [183, "1/1200"], [184, "1/1200"], [186, "1/48"], [187, "1/1200"]]}}, {"path": [[0, 1], [53, 1], [116, 1], [15, 1], [118, 1], [2, 1], [58, 1], [139, 1], [25, 0]], "dual": {"bound": "-1319/37500", "weights": [[8, "401/6000"], [19, "1/100"], [43, "433/6000"], [181, "119/3000"], [182, "1/1200"], [183, "19/3000"], [184, "881/1000"]]}}, {"path": [[0, 1], [53, 1], [116, 1], [15, 1], [118, 1], [2, 1], [58, 1], [139, 0], [18, 1], [4, 1], [3, 1], [25, 1]], "dual": {"bound": "0", "weights": [[9, "6/325"], [19, "1/1300"], [22, "3/325"], [46, "9/1300"], [49, "319/325"], [183, "1/650"], [186, "3/325"], [187, "1/1300"]]}}, {"path": [[0, 1], [53, 1], [116, 1], [15, 1], [118, 1], [2, 1], [58, 1], [139, 0], [18, 1], [4, 1], [3, 1], [25, 0]], "dual": {"bound": "-26/705", "weights": [[6, "353/4700"], [8, "10/47"], [19, "4/47"], [43, "3/47"], [180, "1/141"], [185, "5/141"], [187, "42/47"]]}}, {"path": [[0, 1], [53, 1], [116, 1], [15, 1], [118, 1], [2, 1], [58, 1], [139, 0], [18, 1], [4, 1], [3, 0], [25, 1]], "dual": {"bound": "-1/553", "weights": [[6, "1721/7900"], [21, "241/1106"], [22, "11/1106"], [46, "216/553"], [111, "23/553"], [183, "25/553"], [186, "265/553"], [187, "24/553"]]}}, {"path": [[0, 1], [53, 1], [116, 1], [15, 1], [118, 1], [2, 1], [58, 1], [139, 0], [18, 1], [4, 1], [3, 0], [25, 0]], "dual": {"bound": "-26/705", "weights": [[6, "353/4700"], [8, "10/47"], [19, "4/47"], [180, "1/141"], [183, "3/47"], [185, "5/141"], [186, "18/47"], [187, "24/47"]]}}, {"path": [[0, 1], [53, 1], [116, 1], [15, 1], [118, 1], [2, 1], [58, 1], [139, 0], [18, 1], [4, 0], [25, 1]], "dual": {"bound": "-1/553", "weights": [[6, "1721/7900"], [21, "241/1106"], [22, "11/1106"], [46, "216/553"], [111, "23/553"], [183, "25/553"], [185, "265/553"], [186, "24/553"]]}}, {"path": [[0, 1], [53, 1], [116, 1], [15, 1], [118, 1], [2, 1], [58, 1], [139, 0], [18, 1], [4, 0], [25, 0]], "dual": {"bound": "-11/315", "weights": [[6, "37/6300"], [8, "5/63"], [19, "1/63"], [180, "1/126"], [181, "5/126"], [183, "1/14"], [185, "3/7"], [186, "19/42"]]}}, {"path": [[0, 1], [53, 1], [116, 1], [15, 1], [118, 1], [2, 1], [58, 1], [139, 0], [18, 0], [4, 1], [3, 1], [25, 1]], "dual": {"bound": "-1/1633", "weights": [[6, "37367/163300"], [22, "390/1633"], [46, "672/1633"], [111, "1/23"], [183, "65/1633"], [184, "8/1633"], [186, "745/1633"], [187, "72/1633"]]}}, {"path": [[0, 1], [53, 1], [116, 1], [15, 1], [118, 1], [2, 1], [58, 1], [139, 0], [18, 0], [4, 1], [3, 1], [25, 0]], "dual": {"bound": "-13/285", "weights": [[6, "181/1900"], [8, "5/19"], [19, "2/19"], [180, "1/114"], [184, "3/38"], [185, "5/114"], [187, "33/38"]]}}, {"path": [[0, 1], [53, 1], [116, 1], [15, 1], [118, 1], [2, 1], [58, 1], [139, 0], [18, 0], [4, 1], [3, 0], [25, 1]], "dual": {"bound": "-1/553", "weights": [[6, "1721/7900"], [21, "241/1106"], [22, "11/1106"], [46, "216/553"], [111, "23/553"], [183, "25/553"], [186, "265/553"], [187, "24/553"]]}}, {"path": [[0, 1], [53, 1], [116, 1], [15, 1], [118, 1], [2, 1], [58, 1], [139, 0], [18, 0], [4, 1], [3, 0], [25, 0]], "dual": {"bound": "-13/285", "weights": [[6, "181/1900"], [8, "5/19"], [19, "2/19"], [180, "1/114"], [184, "3/38"], [185, "5/114"], [187, "33/38"]]}}, {"path": [[0, 1], [53, 1], [116, 1], [15, 1], [118, 1], [2, 1], [58, 1], [139, 0], [18, 0], [4, 0], [25, 1]], "dual": {"bound": "-1/553", "weights": [[6, "1721/7900"], [21, "241/1106"], [22, "11/1106"], [46, "216/553"], [111, "23/553"], [183, "25/553"], [185, "265/553"], [186, "24/553"]]}}, {"path": [[0, 1], [53, 1], [116, 1], [15, 1], [118, 1], [2, 1], [58, 1], [139, 0], [18, 0], [4, 0], [25, 0]], "dual": {"bound": "-2/45", "weights": [[6, "101/9900"], [8, "10/99"], [19, "2/99"], [180, "1/99"], [181, "5/99"], [184, "1/11"], [186, "28/33"]]}}, {"path": [[0, 1], [53, 1], [116, 1], [15, 1], [118, 1], [2, 1], [58, 0], [139, 1], [25, 1], [4, 1], [3, 1]], "dual": {"bound": "-4/5075", "weights": [[6, "3197/20300"], [8, "19/58"], [22, "34/203"], [46, "94/203"], [182, "1/406"], [183, "11/406"], [184, "3/203"], [186, "100/203"]]}}, {"path": [[0, 1], [53, 1], [116, 1], [15, 1], [118, 1], [2, 1], [58, 0], [139, 1], [25, 1], [4, 1], [3, 0], [161, 1]], "dual": {"bound": "0", "weights": [[21, "1/100"], [46, "13/700"], [49, "337/350"], [182, "1/1400"], [184, "1/1400"], [186, "23/1400"], [187, "1/1400"]]}}, {"path": [[0, 1], [53, 1], [116, 1], [15, 1], [118, 1], [2, 1], [58, 0], [139, 1], [25, 1], [4, 1], [3, 0], [161, 0]], "dual": {"bound": "0", "weights": [[21, "1/100"], [46, "11/600"], [49, "23/24"], [183, "1/1200"], [184, "1/1200"], [186, "1/48"], [187, "1/1200"]]}}, {"path": [[0, 1], [53, 1], [116, 1], [15, 1], [118, 1], [2, 1], [58, 0], [139, 1], [25, 1], [4, 0], [161, 1]], "dual": {"bound": "0", "weights": [[21, "1/100"], [46, "13/700"], [49, "337/350"], [182, "1/1400"], [184, "1/1400"], [185, "23/1400"], [186, "1/1400"]]}}, {"path": [[0, 1], [53, 1], [116, 1], [15, 1], [118, 1], [2, 1], [58, 0], [139, 1], [25, 1], [4, 0], [161, 0]], "dual": {"bound": "0", "weights": [[21, "1/100"], [46, "11/600"], [49, "23/24"], [183, "1/1200"], [184, "1/1200"], [185, "1/48"], [186, "1/1200"]]}}, {"path": [[0, 1], [53, 1], [116, 1], [15, 1], [118, 1], [2, 1], [58, 0], [139, 1], [25, 0]], "dual": {"bound": "-17/432", "weights": [[6, "91/900"], [9, "7/18"], [22, "1/9"], [32, "1/18"], [181, "1/18"], [182, "1/9"], [184, "7/9"]]}}, {"path": [[0, 1], [53, 1], [116, 1], [15, 1], [118, 1], [2, 1], [58, 0], [139, 0], [4, 1], [3, 1], [25, 1]], "dual": {"bound": "-1/1633", "weights": [[6, "37367/163300"], [22, "390/1633"], [46, "672/1633"], [111, "1/23"], [182, "12/1633"], [183, "61/1633"], [185, "745/1633"], [186, "72/1633"]]}}, {"path": [[0, 1], [53, 1], [116, 1], [15, 1], [118, 1], [2, 1], [58, 0], [139, 0], [4, 1], [3, 1], [25, 0]], "dual": {"bound": "-1/24", "weights": [[6, "379/2100"], [9, "1/7"], [19, "2/21"], [22, "2/21"], [43, "1/42"], [182, "1/14"], [184, "1/21"], [186, "6/7"]]}}, {"path": [[0, 1], [53, 1], [116, 1], [15, 1], [118, 1], [2, 1], [58, 0], [139, 0], [4, 1], [3, 0], [25, 1]], "dual": {"bound": "-1/553", "weights": [[6, "1721/7900"], [21, "241/1106"], [22, "11/1106"], [46, "216/553"], [111, "23/553"], [183, "25/553"], [185, "265/553"], [186, "24/553"]]}}, {"path": [[0, 1], [53, 1], [116, 1], [15, 1], [118, 1], [2, 1], [58, 0], [139, 0], [4, 1], [3, 0], [25, 0]], "dual": {"bound": "-7/160", "weights": [[6, "19/100"], [9, "3/20"], [19, "1/10"], [22, "1/10"], [182, "1/10"], [184, "1/20"], [185, "1/20"], [186, "4/5"]]}}, {"path": [[0, 1], [53, 1], [116, 1], [15, 1], [118, 1], [2, 1], [58, 0], [139, 0], [4, 0], [25, 1]], "dual": {"bound": "-1/553", "weights": [[6, "1721/7900"], [21, "241/1106"], [22, "11/1106"], [46, "216/553"], [111, "23/553"], [183, "25/553"], [184, "265/553"], [185, "24/553"]]}}, {"path": [[0, 1], [53, 1], [116, 1], [15, 1], [118, 1], [2, 1], [58, 0], [139, 0], [4, 0], [25, 0]], "dual": {"bound": "-1/24", "weights": [[6, "183/1700"], [9, "7/17"], [22, "2/17"], [181, "1/17"], [182, "2/17"], [184, "1/17"], [185, "13/17"]]}}, {"path": [[0, 1], [53, 1], [116, 1], [15, 1], [118, 1], [2, 0]], "dual": {"bound": "-4691/146250", "weights": [[8, "899/14400"], [19, "1/150"], [21, "1/300"], [169, "541/7488"], [176, "1997/46800"], [179, "791/187200"], [180, "83/9360"], [181, "13603/15600"]]}}, {"path": [[0, 1], [53, 1], [116, 1], [15, 1], [118, 0], [2, 1], [18, 1], [3, 1], [25, 1]], "dual": {"bound": "-1/528", "weights": [[6, "503/2200"], [21, "1/1056"], [22, "251/1056"], [46, "9/22"], [111, "23/528"], [180, "25/528"], [183, "5/11"], [184, "1/22"]]}}, {"path": [[0, 1], [53, 1], [116, 1], [15, 1], [118, 0], [2, 1], [18, 1], [3, 1], [25, 0]], "dual": {"bound": "-50/1433", "weights": [[6, "95701/429900"], [19, "210/1433"], [22, "370/4299"], [43, "224/4299"], [180, "37/1433"], [181, "185/4299"], [182, "35/4299"], [184, "1248/1433"]]}}, {"path": [[0, 1], [53, 1], [116, 1], [15, 1], [118, 0], [2, 1], [18, 1], [3, 0], [25, 1]], "dual": {"bound": "-1/528", "weights": [[6, "503/2200"], [21, "241/1056"], [22, "1/96"], [46, "9/22"], [111, "23/528"], [180, "25/528"], [183, "5/11"], [184, "1/22"]]}}, {"path": [[0, 1], [53, 1], [116, 1], [15, 1], [118, 0], [2, 1], [18, 1], [3, 0], [25, 0]], "dual": {"bound": "-17/456", "weights": [[6, "3/950"], [22, "1/76"], [178, "7/912"], [180, "25/304"], [181, "41/912"], [183, "119/304"], [184, "9/19"]]}}, {"path": [[0, 1], [53, 1], [116, 1], [15, 1], [118, 0], [2, 1], [18, 0], [3, 1], [25, 1]], "dual": {"bound": "-1/528", "weights": [[6, "503/2200"], [21, "1/1056"], [22, "251/1056"], [46, "9/22"], [111, "23/528"], [180, "25/528"], [183, "5/11"], [184, "1/22"]]}}, {"path": [[0, 1], [53, 1], [116, 1], [15, 1], [118, 0], [2, 1], [18, 0], [3, 1], [25, 0]], "dual": {"bound": "-85/1923", "weights": [[6, "359/64100"], [22, "10/641"], [178, "35/3846"], [180, "3/641"], [181, "205/3846"], [182, "119/1282"], [184, "1077/1282"]]}}, {"path": [[0, 1], [53, 1], [116, 1], [15, 1], [118, 0], [2, 1], [18, 0], [3, 0], [25, 1]], "dual": {"bound": "-1/528", "weights": [[6, "503/2200"], [21, "241/1056"], [22, "1/96"], [46, "9/22"], [111, "23/528"], [180, "25/528"], [183, "5/11"], [184, "1/22"]]}}, {"path": [[0, 1], [53, 1], [116, 1], [15, 1], [118, 0], [2, 1], [18, 0], [3, 0], [25, 0]], "dual": {"bound": "-85/1923", "weights": [[6, "359/64100"], [22, "10/641"], [178, "35/3846"], [180, "3/641"], [181, "205/3846"], [182, "119/1282"], [184, "1077/1282"]]}}, {"path": [[0, 1], [53, 1], [116, 1], [15, 1], [118, 0], [2, 0]], "dual": {"bound": "-4144/73125", "weights": [[6, "419/11700"], [8, "44/2925"], [21, "134/2925"], [176, "238/2925"], [178, "61/2925"], [179, "1/225"], [180, "134/975"], [181, "737/975"]]}}, {"path": [[0, 1], [53, 1], [116, 1], [15, 0], [139, 1], [25, 1], [3, 1]], "dual": {"bound": "-3/3425", "weights": [[6, "2113/13700"], [8, "187/548"], [22, "45/274"], [46, "127/274"], [179, "1/548"], [180, "15/548"], [181, "2/137"], [182, "135/274"]]}}, {"path": [[0, 1], [53, 1], [116, 1], [15, 0], [139, 1], [25, 1], [3, 0], [161, 1]], "dual": {"bound": "0", "weights": [[21, "1/100"], [46, "13/700"], [49, "27/28"], [179, "1/1400"], [181, "1/1400"], [182, "3/200"], [183, "1/1400"]]}}, {"path": [[0, 1], [53, 1], [116, 1], [15, 0], [139, 1], [25, 1], [3, 0], [161, 0]], "dual": {"bound": "0", "weights": [[21, "1/100"], [46, "11/600"], [49, "23/24"], [180, "1/1200"], [181, "1/1200"], [182, "1/48"], [183, "1/1200"]]}}, {"path": [[0, 1], [53, 1], [116, 1], [15, 0], [139, 1], [25, 0]], "dual": {"bound": "-1349/30625", "weights": [[8, "1009/19600"], [21, "1/200"], [22, "1/200"], [176, "121/1960"], [177, "139/9800"], [179, "2033/19600"], [180, "109/19600"], [181, "1597/1960"]]}}, {"path": [[0, 1], [53, 1], [116, 1], [15, 0], [139, 0], [25, 1], [3, 1]], "dual": {"bound": "-1/1081", "weights": [[6, "24719/108100"], [22, "258/1081"], [46, "444/1081"], [111, "1/23"], [179, "6/1081"], [180, "43/1081"], [181, "48/1081"], [182, "493/1081"]]}}, {"path": [[0, 1], [53, 1], [116, 1], [15, 0], [139, 0], [25, 1], [3, 0]], "dual": {"bound": "-1/553", "weights": [[6, "1721/7900"], [21, "241/1106"], [22, "11/1106"], [46, "216/553"], [111, "23/553"], [180, "25/553"], [181, "24/553"], [182, "265/553"]]}}, {"path": [[0, 1], [53, 1], [116, 1], [15, 0], [139, 0], [25, 0]], "dual": {"bound": "-4391/101250", "weights": [[8, "4/75"], [21, "1/200"], [22, "1/200"], [176, "1009/16200"], [177, "29/2025"], [178, "109/16200"], [179, "559/5400"], [181, "4391/5400"]]}}, {"path": [[0, 1], [53, 1], [116, 0], [2, 1], [3, 1]], "dual": {"bound": "-1/136", "weights": [[6, "883/1700"], [9, "21/17"], [22, "9/17"], [46, "7/17"], [178, "2/17"], [179, "1/17"], [180, "7/17"]]}}, {"path": [[0, 1], [53, 1], [116, 0], [2, 1], [3, 0]], "dual": {"bound": "-1/112", "weights": [[6, "393/700"], [9, "3/2"], [21, "3/14"], [22, "5/14"], [46, "2/7"], [176, "1/14"], [178, "1/7"], [180, "1/2"]]}}, {"path": [[0, 1], [53, 1], [116, 0], [2, 0]], "dual": {"bound": "-61/984", "weights": [[6, "359/4100"], [9, "11/41"], [22, "4/41"], [176, "4/41"], [177, "1/41"], [178, "7/41"], [179, "29/41"]]}}, {"path": [[0, 1], [53, 0], [3, 1], [25, 1]], "dual": {"bound": "-1/384", "weights": [[6, "271/400"], [9, "23/16"], [21, "7/32"], [22, "15/32"], [46, "9/16"], [177, "1/8"], [178, "1/4"], [179, "1/16"]]}}, {"path": [[0, 1], [53, 0], [3, 1], [25, 0]], "dual": {"bound": "-11/240", "weights": [[6, "39/100"], [9, "13/10"], [21, "1/20"], [22, "7/20"], [176, "1/10"], [177, "1/5"], [178, "1/10"], [179, "3/5"]]}}, {"path": [[0, 1], [53, 0], [3, 0], [25, 1]], "dual": {"bound": "-1/384", "weights": [[6, "271/400"], [9, "23/16"], [21, "11/32"], [22, "11/32"], [46, "9/16"], [177, "1/8"], [178, "1/4"], [179, "1/16"]]}}, {"path": [[0, 1], [53, 0], [3, 0], [25, 0]], "dual": {"bound": "-11/240", "weights": [[6, "39/100"], [9, "13/10"], [21, "1/10"], [22, "3/10"], [176, "1/10"], [177, "1/5"], [178, "1/10"], [179, "3/5"]]}}, {"path": [[0, 0]], "dual": {"bound": "0", "weights": [[3, "201/200"], [20, "1/100"], [69, "199/400"], [176, "201/400"]]}}]}
\end{filecontents*}
\begin{filecontents*}[overwrite]{harman_H2_uniform_certificate.json}
{"summary": {"region": "H2", "theta_lower": "1/25", "theta_upper": "1/24", "gain": "1/100", "nodes": 71, "leaves": 36, "stack": 0, "maxdepth": 9, "seconds": 0.2483537197113037, "bad": []}, "base": [[["-1", "0", "0", "0", "0", "0", "0", "0"], "0", ["nonnegative", 0]], [["0", "-1", "0", "0", "0", "0", "0", "0"], "0", ["nonnegative", 1]], [["0", "0", "-1", "0", "0", "0", "0", "0"], "0", ["nonnegative", 2]], [["0", "0", "0", "-1", "0", "0", "0", "0"], "0", ["nonnegative", 3]], [["0", "0", "0", "0", "-1", "0", "0", "0"], "0", ["nonnegative", 4]], [["0", "0", "0", "0", "0", "-1", "0", "0"], "0", ["nonnegative", 5]], [["0", "0", "0", "0", "0", "0", "-1", "0"], "0", ["nonnegative", 6]], [["0", "0", "0", "0", "0", "0", "0", "-1"], "0", ["nonnegative", 7]], [["0", "0", "0", "0", "0", "0", "0", "-1"], "-1/25", "theta_lower"], [["0", "0", "0", "0", "0", "0", "0", "1"], "1/24", "theta_upper"], [["1", "0", "0", "0", "0", "0", "0", "0"], "1/2", "m_half"], [["-1", "1", "0", "0", "0", "0", "0", "0"], "0", "h_le_m"], [["1", "1", "0", "0", "0", "0", "0", "0"], "1", "sum_lengths"], [["-1", "0", "2", "0", "0", "0", "0", "0"], "0", "a_upper"], [["0", "-1", "0", "2", "0", "0", "0", "0"], "0", "b_upper"], [["1", "1", "0", "0", "2", "0", "0", "0"], "1", "c_upper"], [["0", "0", "0", "0", "0", "1", "0", "1"], "1", "f_upper"], [["1", "0", "0", "0", "0", "0", "1", "0"], "1/2", "z_m_half"], [["-1", "1", "0", "0", "0", "0", "1", "0"], "0", "z_h_le_m"], [["0", "-1", "0", "0", "0", "0", "1", "2"], "0", "z_h_low"], [["0", "0", "0", "-1", "0", "0", "1", "0"], "0", "z_prime"], [["0", "1", "0", "0", "0", "0", "1", "-4"], "0", "z_h_high"], [["4", "3", "0", "0", "0", "0", "1", "-2"], "2", "z_band_high"], [["-4", "-3", "0", "0", "0", "0", "1", "-2"], "-2", "z_band_low"], [["-1", "-1", "0", "0", "-2", "1", "0", "1"], "0", ["mean", [0, 0, 1]]], [["0", "0", "0", "0", "-6", "1", "0", "0"], "0", ["mean", [0, 0, 3]]], [["0", "1", "0", "-2", "0", "1", "0", "1"], "1", ["mean", [0, 1, 0]]], [["0", "2", "0", "-6", "0", "1", "0", "1"], "1", ["halasz", [0, 1, 0]]], [["-1", "0", "0", "-2", "-2", "1", "0", "1"], "0", ["mean", [0, 1, 1]]], [["0", "0", "0", "-2", "-4", "1", "0", "0"], "0", ["mean", [0, 1, 2]]], [["0", "0", "0", "-2", "-6", "1", "0", "0"], "0", ["mean", [0, 1, 3]]], [["0", "2", "0", "-4", "0", "1", "0", "1"], "1", ["mean", [0, 2, 0]]], [["0", "4", "0", "-12", "0", "1", "0", "1"], "1", ["halasz", [0, 2, 0]]], [["-1", "1", "0", "-4", "-2", "1", "0", "1"], "0", ["mean", [0, 2, 1]]], [["0", "0", "0", "-4", "-4", "1", "0", "0"], "0", ["mean", [0, 2, 2]]], [["0", "3", "0", "-6", "0", "1", "0", "1"], "1", ["mean", [0, 3, 0]]], [["0", "0", "0", "-6", "-4", "1", "0", "0"], "0", ["mean", [0, 3, 2]]], [["0", "4", "0", "-8", "0", "1", "0", "1"], "1", ["mean", [0, 4, 0]]], [["0", "0", "0", "-8", "-4", "1", "0", "0"], "0", ["mean", [0, 4, 2]]], [["0", "5", "0", "-10", "0", "1", "0", "1"], "1", ["mean", [0, 5, 0]]], [["0", "0", "0", "-10", "-4", "1", "0", "0"], "0", ["mean", [0, 5, 2]]], [["0", "0", "0", "-12", "-4", "1", "0", "0"], "0", ["mean", [0, 6, 2]]], [["0", "0", "0", "-14", "-2", "1", "0", "0"], "0", ["mean", [0, 7, 1]]], [["0", "0", "0", "-14", "-4", "1", "0", "0"], "0", ["mean", [0, 7, 2]]], [["0", "0", "0", "-16", "-2", "1", "0", "0"], "0", ["mean", [0, 8, 1]]], [["0", "0", "0", "-18", "-2", "1", "0", "0"], "0", ["mean", [0, 9, 1]]], [["0", "0", "0", "-20", "-2", "1", "0", "0"], "0", ["mean", [0, 10, 1]]], [["0", "0", "0", "-22", "-2", "1", "0", "0"], "0", ["mean", [0, 11, 1]]], [["0", "0", "0", "-24", "0", "1", "0", "0"], "0", ["mean", [0, 12, 0]]], [["0", "0", "0", "-24", "-2", "1", "0", "0"], "0", ["mean", [0, 12, 1]]], [["1", "0", "-2", "0", "0", "1", "0", "1"], "1", ["mean", [1, 0, 0]]], [["2", "0", "-6", "0", "0", "1", "0", "1"], "1", ["halasz", [1, 0, 0]]], [["0", "-1", "-2", "0", "-2", "1", "0", "1"], "0", ["mean", [1, 0, 1]]], [["0", "0", "-2", "0", "-4", "1", "0", "0"], "0", ["mean", [1, 0, 2]]], [["1", "1", "-2", "-2", "0", "1", "0", "1"], "1", ["mean", [1, 1, 0]]], [["0", "0", "-2", "-2", "-2", "1", "0", "0"], "0", ["mean", [1, 1, 1]]], [["0", "0", "-2", "-2", "-4", "1", "0", "0"], "0", ["mean", [1, 1, 2]]], [["1", "2", "-2", "-4", "0", "1", "0", "1"], "1", ["mean", [1, 2, 0]]], [["0", "0", "-2", "-4", "-2", "1", "0", "0"], "0", ["mean", [1, 2, 1]]], [["1", "3", "-2", "-6", "0", "1", "0", "1"], "1", ["mean", [1, 3, 0]]], [["0", "0", "-2", "-6", "-2", "1", "0", "0"], "0", ["mean", [1, 3, 1]]], [["0", "0", "-2", "-8", "-2", "1", "0", "0"], "0", ["mean", [1, 4, 1]]], [["0", "0", "-2", "-10", "-2", "1", "0", "0"], "0", ["mean", [1, 5, 1]]], [["0", "0", "-2", "-12", "-2", "1", "0", "0"], "0", ["mean", [1, 6, 1]]], [["0", "0", "-2", "-14", "0", "1", "0", "0"], "0", ["mean", [1, 7, 0]]], [["0", "0", "-2", "-14", "-2", "1", "0", "0"], "0", ["mean", [1, 7, 1]]], [["0", "0", "-2", "-16", "0", "1", "0", "0"], "0", ["mean", [1, 8, 0]]], [["0", "0", "-2", "-18", "0", "1", "0", "0"], "0", ["mean", [1, 9, 0]]], [["0", "0", "-2", "-20", "0", "1", "0", "0"], "0", ["mean", [1, 10, 0]]], [["0", "0", "-2", "-22", "0", "1", "0", "0"], "0", ["mean", [1, 11, 0]]], [["0", "0", "-2", "-24", "0", "1", "0", "0"], "0", ["mean", [1, 12, 0]]], [["2", "0", "-4", "0", "0", "1", "0", "1"], "1", ["mean", [2, 0, 0]]], [["0", "0", "-4", "0", "-2", "1", "0", "0"], "0", ["mean", [2, 0, 1]]], [["0", "0", "-4", "-2", "-2", "1", "0", "0"], "0", ["mean", [2, 1, 1]]], [["0", "0", "-4", "-4", "0", "1", "0", "0"], "0", ["mean", [2, 2, 0]]], [["0", "0", "-4", "-4", "-2", "1", "0", "0"], "0", ["mean", [2, 2, 1]]], [["0", "0", "-4", "-6", "0", "1", "0", "0"], "0", ["mean", [2, 3, 0]]], [["0", "0", "-4", "-8", "0", "1", "0", "0"], "0", ["mean", [2, 4, 0]]], [["0", "0", "-4", "-10", "0", "1", "0", "0"], "0", ["mean", [2, 5, 0]]], [["0", "0", "-4", "-12", "0", "1", "0", "0"], "0", ["mean", [2, 6, 0]]], [["0", "0", "-4", "-14", "0", "1", "0", "0"], "0", ["mean", [2, 7, 0]]], [["0", "0", "-4", "-16", "0", "1", "0", "0"], "0", ["mean", [2, 8, 0]]], [["0", "0", "-6", "0", "0", "1", "0", "0"], "0", ["mean", [3, 0, 0]]], [["0", "0", "-6", "-2", "0", "1", "0", "0"], "0", ["mean", [3, 1, 0]]], [["0", "0", "-6", "-4", "0", "1", "0", "0"], "0", ["mean", [3, 2, 0]]], [["0", "0", "-6", "-6", "0", "1", "0", "0"], "0", ["mean", [3, 3, 0]]]], "disjunctions": [[[["0", "0", "0", "0", "-2", "1", "0", "0"], "0", ["halasz_diag", [0, 0, 1]]], [["-2", "-2", "0", "0", "-6", "1", "0", "1"], "-1", ["halasz_off", [0, 0, 1]]]], [[["0", "0", "0", "0", "-4", "1", "0", "0"], "0", ["mean_diag", [0, 0, 2]]], [["-2", "-2", "0", "0", "-4", "1", "0", "1"], "-1", ["mean_off", [0, 0, 2]]]], [[["0", "0", "0", "0", "-4", "1", "0", "0"], "0", ["halasz_diag", [0, 0, 2]]], [["-4", "-4", "0", "0", "-12", "1", "0", "1"], "-3", ["halasz_off", [0, 0, 2]]]], [[["0", "0", "0", "0", "-6", "1", "0", "0"], "0", ["halasz_diag", [0, 0, 3]]], [["-6", "-6", "0", "0", "-18", "1", "0", "1"], "-5", ["halasz_off", [0, 0, 3]]]], [[["0", "0", "0", "-2", "-2", "1", "0", "0"], "0", ["halasz_diag", [0, 1, 1]]], [["-2", "0", "0", "-6", "-6", "1", "0", "1"], "-1", ["halasz_off", [0, 1, 1]]]], [[["0", "0", "0", "-2", "-4", "1", "0", "0"], "0", ["halasz_diag", [0, 1, 2]]], [["-4", "-2", "0", "-6", "-12", "1", "0", "1"], "-3", ["halasz_off", [0, 1, 2]]]], [[["0", "0", "0", "-2", "-6", "1", "0", "0"], "0", ["halasz_diag", [0, 1, 3]]], [["-6", "-4", "0", "-6", "-18", "1", "0", "1"], "-5", ["halasz_off", [0, 1, 3]]]], [[["0", "0", "0", "-4", "-2", "1", "0", "0"], "0", ["halasz_diag", [0, 2, 1]]], [["-2", "2", "0", "-12", "-6", "1", "0", "1"], "-1", ["halasz_off", [0, 2, 1]]]], [[["0", "0", "0", "-4", "-4", "1", "0", "0"], "0", ["halasz_diag", [0, 2, 2]]], [["-4", "0", "0", "-12", "-12", "1", "0", "1"], "-3", ["halasz_off", [0, 2, 2]]]], [[["0", "0", "0", "-6", "0", "1", "0", "0"], "0", ["halasz_diag", [0, 3, 0]]], [["0", "6", "0", "-18", "0", "1", "0", "1"], "1", ["halasz_off", [0, 3, 0]]]], [[["0", "0", "0", "-6", "-2", "1", "0", "0"], "0", ["mean_diag", [0, 3, 1]]], [["-1", "2", "0", "-6", "-2", "1", "0", "1"], "0", ["mean_off", [0, 3, 1]]]], [[["0", "0", "0", "-6", "-2", "1", "0", "0"], "0", ["halasz_diag", [0, 3, 1]]], [["-2", "4", "0", "-18", "-6", "1", "0", "1"], "-1", ["halasz_off", [0, 3, 1]]]], [[["0", "0", "0", "-6", "-4", "1", "0", "0"], "0", ["halasz_diag", [0, 3, 2]]], [["-4", "2", "0", "-18", "-12", "1", "0", "1"], "-3", ["halasz_off", [0, 3, 2]]]], [[["0", "0", "0", "-8", "0", "1", "0", "0"], "0", ["halasz_diag", [0, 4, 0]]], [["0", "8", "0", "-24", "0", "1", "0", "1"], "1", ["halasz_off", [0, 4, 0]]]], [[["0", "0", "0", "-8", "-2", "1", "0", "0"], "0", ["mean_diag", [0, 4, 1]]], [["-1", "3", "0", "-8", "-2", "1", "0", "1"], "0", ["mean_off", [0, 4, 1]]]], [[["0", "0", "0", "-8", "-2", "1", "0", "0"], "0", ["halasz_diag", [0, 4, 1]]], [["-2", "6", "0", "-24", "-6", "1", "0", "1"], "-1", ["halasz_off", [0, 4, 1]]]], [[["0", "0", "0", "-8", "-4", "1", "0", "0"], "0", ["halasz_diag", [0, 4, 2]]], [["-4", "4", "0", "-24", "-12", "1", "0", "1"], "-3", ["halasz_off", [0, 4, 2]]]], [[["0", "0", "0", "-10", "0", "1", "0", "0"], "0", ["halasz_diag", [0, 5, 0]]], [["0", "10", "0", "-30", "0", "1", "0", "1"], "1", ["halasz_off", [0, 5, 0]]]], [[["0", "0", "0", "-10", "-2", "1", "0", "0"], "0", ["mean_diag", [0, 5, 1]]], [["-1", "4", "0", "-10", "-2", "1", "0", "1"], "0", ["mean_off", [0, 5, 1]]]], [[["0", "0", "0", "-10", "-2", "1", "0", "0"], "0", ["halasz_diag", [0, 5, 1]]], [["-2", "8", "0", "-30", "-6", "1", "0", "1"], "-1", ["halasz_off", [0, 5, 1]]]], [[["0", "0", "0", "-10", "-4", "1", "0", "0"], "0", ["halasz_diag", [0, 5, 2]]], [["-4", "6", "0", "-30", "-12", "1", "0", "1"], "-3", ["halasz_off", [0, 5, 2]]]], [[["0", "0", "0", "-12", "0", "1", "0", "0"], "0", ["mean_diag", [0, 6, 0]]], [["0", "6", "0", "-12", "0", "1", "0", "1"], "1", ["mean_off", [0, 6, 0]]]], [[["0", "0", "0", "-12", "0", "1", "0", "0"], "0", ["halasz_diag", [0, 6, 0]]], [["0", "12", "0", "-36", "0", "1", "0", "1"], "1", ["halasz_off", [0, 6, 0]]]], [[["0", "0", "0", "-12", "-2", "1", "0", "0"], "0", ["mean_diag", [0, 6, 1]]], [["-1", "5", "0", "-12", "-2", "1", "0", "1"], "0", ["mean_off", [0, 6, 1]]]], [[["0", "0", "0", "-12", "-2", "1", "0", "0"], "0", ["halasz_diag", [0, 6, 1]]], [["-2", "10", "0", "-36", "-6", "1", "0", "1"], "-1", ["halasz_off", [0, 6, 1]]]], [[["0", "0", "0", "-12", "-4", "1", "0", "0"], "0", ["halasz_diag", [0, 6, 2]]], [["-4", "8", "0", "-36", "-12", "1", "0", "1"], "-3", ["halasz_off", [0, 6, 2]]]], [[["0", "0", "0", "-14", "0", "1", "0", "0"], "0", ["mean_diag", [0, 7, 0]]], [["0", "7", "0", "-14", "0", "1", "0", "1"], "1", ["mean_off", [0, 7, 0]]]], [[["0", "0", "0", "-14", "0", "1", "0", "0"], "0", ["halasz_diag", [0, 7, 0]]], [["0", "14", "0", "-42", "0", "1", "0", "1"], "1", ["halasz_off", [0, 7, 0]]]], [[["0", "0", "0", "-14", "-2", "1", "0", "0"], "0", ["halasz_diag", [0, 7, 1]]], [["-2", "12", "0", "-42", "-6", "1", "0", "1"], "-1", ["halasz_off", [0, 7, 1]]]], [[["0", "0", "0", "-14", "-4", "1", "0", "0"], "0", ["halasz_diag", [0, 7, 2]]], [["-4", "10", "0", "-42", "-12", "1", "0", "1"], "-3", ["halasz_off", [0, 7, 2]]]], [[["0", "0", "0", "-16", "0", "1", "0", "0"], "0", ["mean_diag", [0, 8, 0]]], [["0", "8", "0", "-16", "0", "1", "0", "1"], "1", ["mean_off", [0, 8, 0]]]], [[["0", "0", "0", "-16", "0", "1", "0", "0"], "0", ["halasz_diag", [0, 8, 0]]], [["0", "16", "0", "-48", "0", "1", "0", "1"], "1", ["halasz_off", [0, 8, 0]]]], [[["0", "0", "0", "-16", "-2", "1", "0", "0"], "0", ["halasz_diag", [0, 8, 1]]], [["-2", "14", "0", "-48", "-6", "1", "0", "1"], "-1", ["halasz_off", [0, 8, 1]]]], [[["0", "0", "0", "-18", "0", "1", "0", "0"], "0", ["mean_diag", [0, 9, 0]]], [["0", "9", "0", "-18", "0", "1", "0", "1"], "1", ["mean_off", [0, 9, 0]]]], [[["0", "0", "0", "-18", "0", "1", "0", "0"], "0", ["halasz_diag", [0, 9, 0]]], [["0", "18", "0", "-54", "0", "1", "0", "1"], "1", ["halasz_off", [0, 9, 0]]]], [[["0", "0", "0", "-18", "-2", "1", "0", "0"], "0", ["halasz_diag", [0, 9, 1]]], [["-2", "16", "0", "-54", "-6", "1", "0", "1"], "-1", ["halasz_off", [0, 9, 1]]]], [[["0", "0", "0", "-20", "0", "1", "0", "0"], "0", ["mean_diag", [0, 10, 0]]], [["0", "10", "0", "-20", "0", "1", "0", "1"], "1", ["mean_off", [0, 10, 0]]]], [[["0", "0", "0", "-20", "0", "1", "0", "0"], "0", ["halasz_diag", [0, 10, 0]]], [["0", "20", "0", "-60", "0", "1", "0", "1"], "1", ["halasz_off", [0, 10, 0]]]], [[["0", "0", "0", "-20", "-2", "1", "0", "0"], "0", ["halasz_diag", [0, 10, 1]]], [["-2", "18", "0", "-60", "-6", "1", "0", "1"], "-1", ["halasz_off", [0, 10, 1]]]], [[["0", "0", "0", "-22", "0", "1", "0", "0"], "0", ["mean_diag", [0, 11, 0]]], [["0", "11", "0", "-22", "0", "1", "0", "1"], "1", ["mean_off", [0, 11, 0]]]], [[["0", "0", "0", "-22", "0", "1", "0", "0"], "0", ["halasz_diag", [0, 11, 0]]], [["0", "22", "0", "-66", "0", "1", "0", "1"], "1", ["halasz_off", [0, 11, 0]]]], [[["0", "0", "0", "-22", "-2", "1", "0", "0"], "0", ["halasz_diag", [0, 11, 1]]], [["-2", "20", "0", "-66", "-6", "1", "0", "1"], "-1", ["halasz_off", [0, 11, 1]]]], [[["0", "0", "0", "-24", "0", "1", "0", "0"], "0", ["halasz_diag", [0, 12, 0]]], [["0", "24", "0", "-72", "0", "1", "0", "1"], "1", ["halasz_off", [0, 12, 0]]]], [[["0", "0", "0", "-24", "-2", "1", "0", "0"], "0", ["halasz_diag", [0, 12, 1]]], [["-2", "22", "0", "-72", "-6", "1", "0", "1"], "-1", ["halasz_off", [0, 12, 1]]]], [[["0", "0", "-2", "0", "-2", "1", "0", "0"], "0", ["halasz_diag", [1, 0, 1]]], [["0", "-2", "-6", "0", "-6", "1", "0", "1"], "-1", ["halasz_off", [1, 0, 1]]]], [[["0", "0", "-2", "0", "-4", "1", "0", "0"], "0", ["halasz_diag", [1, 0, 2]]], [["-2", "-4", "-6", "0", "-12", "1", "0", "1"], "-3", ["halasz_off", [1, 0, 2]]]], [[["0", "0", "-2", "-2", "0", "1", "0", "0"], "0", ["halasz_diag", [1, 1, 0]]], [["2", "2", "-6", "-6", "0", "1", "0", "1"], "1", ["halasz_off", [1, 1, 0]]]], [[["0", "0", "-2", "-2", "-2", "1", "0", "0"], "0", ["halasz_diag", [1, 1, 1]]], [["0", "0", "-6", "-6", "-6", "1", "0", "1"], "-1", ["halasz_off", [1, 1, 1]]]], [[["0", "0", "-2", "-2", "-4", "1", "0", "0"], "0", ["halasz_diag", [1, 1, 2]]], [["-2", "-2", "-6", "-6", "-12", "1", "0", "1"], "-3", ["halasz_off", [1, 1, 2]]]], [[["0", "0", "-2", "-4", "0", "1", "0", "0"], "0", ["halasz_diag", [1, 2, 0]]], [["2", "4", "-6", "-12", "0", "1", "0", "1"], "1", ["halasz_off", [1, 2, 0]]]], [[["0", "0", "-2", "-4", "-2", "1", "0", "0"], "0", ["halasz_diag", [1, 2, 1]]], [["0", "2", "-6", "-12", "-6", "1", "0", "1"], "-1", ["halasz_off", [1, 2, 1]]]], [[["0", "0", "-2", "-6", "0", "1", "0", "0"], "0", ["halasz_diag", [1, 3, 0]]], [["2", "6", "-6", "-18", "0", "1", "0", "1"], "1", ["halasz_off", [1, 3, 0]]]], [[["0", "0", "-2", "-6", "-2", "1", "0", "0"], "0", ["halasz_diag", [1, 3, 1]]], [["0", "4", "-6", "-18", "-6", "1", "0", "1"], "-1", ["halasz_off", [1, 3, 1]]]], [[["0", "0", "-2", "-8", "0", "1", "0", "0"], "0", ["mean_diag", [1, 4, 0]]], [["1", "4", "-2", "-8", "0", "1", "0", "1"], "1", ["mean_off", [1, 4, 0]]]], [[["0", "0", "-2", "-8", "0", "1", "0", "0"], "0", ["halasz_diag", [1, 4, 0]]], [["2", "8", "-6", "-24", "0", "1", "0", "1"], "1", ["halasz_off", [1, 4, 0]]]], [[["0", "0", "-2", "-8", "-2", "1", "0", "0"], "0", ["halasz_diag", [1, 4, 1]]], [["0", "6", "-6", "-24", "-6", "1", "0", "1"], "-1", ["halasz_off", [1, 4, 1]]]], [[["0", "0", "-2", "-10", "0", "1", "0", "0"], "0", ["mean_diag", [1, 5, 0]]], [["1", "5", "-2", "-10", "0", "1", "0", "1"], "1", ["mean_off", [1, 5, 0]]]], [[["0", "0", "-2", "-10", "0", "1", "0", "0"], "0", ["halasz_diag", [1, 5, 0]]], [["2", "10", "-6", "-30", "0", "1", "0", "1"], "1", ["halasz_off", [1, 5, 0]]]], [[["0", "0", "-2", "-10", "-2", "1", "0", "0"], "0", ["halasz_diag", [1, 5, 1]]], [["0", "8", "-6", "-30", "-6", "1", "0", "1"], "-1", ["halasz_off", [1, 5, 1]]]], [[["0", "0", "-2", "-12", "0", "1", "0", "0"], "0", ["mean_diag", [1, 6, 0]]], [["1", "6", "-2", "-12", "0", "1", "0", "1"], "1", ["mean_off", [1, 6, 0]]]], [[["0", "0", "-2", "-12", "0", "1", "0", "0"], "0", ["halasz_diag", [1, 6, 0]]], [["2", "12", "-6", "-36", "0", "1", "0", "1"], "1", ["halasz_off", [1, 6, 0]]]], [[["0", "0", "-2", "-12", "-2", "1", "0", "0"], "0", ["halasz_diag", [1, 6, 1]]], [["0", "10", "-6", "-36", "-6", "1", "0", "1"], "-1", ["halasz_off", [1, 6, 1]]]], [[["0", "0", "-2", "-14", "0", "1", "0", "0"], "0", ["halasz_diag", [1, 7, 0]]], [["2", "14", "-6", "-42", "0", "1", "0", "1"], "1", ["halasz_off", [1, 7, 0]]]], [[["0", "0", "-2", "-14", "-2", "1", "0", "0"], "0", ["halasz_diag", [1, 7, 1]]], [["0", "12", "-6", "-42", "-6", "1", "0", "1"], "-1", ["halasz_off", [1, 7, 1]]]], [[["0", "0", "-2", "-16", "0", "1", "0", "0"], "0", ["halasz_diag", [1, 8, 0]]], [["2", "16", "-6", "-48", "0", "1", "0", "1"], "1", ["halasz_off", [1, 8, 0]]]], [[["0", "0", "-2", "-18", "0", "1", "0", "0"], "0", ["halasz_diag", [1, 9, 0]]], [["2", "18", "-6", "-54", "0", "1", "0", "1"], "1", ["halasz_off", [1, 9, 0]]]], [[["0", "0", "-2", "-20", "0", "1", "0", "0"], "0", ["halasz_diag", [1, 10, 0]]], [["2", "20", "-6", "-60", "0", "1", "0", "1"], "1", ["halasz_off", [1, 10, 0]]]], [[["0", "0", "-2", "-22", "0", "1", "0", "0"], "0", ["halasz_diag", [1, 11, 0]]], [["2", "22", "-6", "-66", "0", "1", "0", "1"], "1", ["halasz_off", [1, 11, 0]]]], [[["0", "0", "-2", "-24", "0", "1", "0", "0"], "0", ["halasz_diag", [1, 12, 0]]], [["2", "24", "-6", "-72", "0", "1", "0", "1"], "1", ["halasz_off", [1, 12, 0]]]], [[["0", "0", "-4", "0", "0", "1", "0", "0"], "0", ["halasz_diag", [2, 0, 0]]], [["4", "0", "-12", "0", "0", "1", "0", "1"], "1", ["halasz_off", [2, 0, 0]]]], [[["0", "0", "-4", "0", "-2", "1", "0", "0"], "0", ["halasz_diag", [2, 0, 1]]], [["2", "-2", "-12", "0", "-6", "1", "0", "1"], "-1", ["halasz_off", [2, 0, 1]]]], [[["0", "0", "-4", "-2", "0", "1", "0", "0"], "0", ["mean_diag", [2, 1, 0]]], [["2", "1", "-4", "-2", "0", "1", "0", "1"], "1", ["mean_off", [2, 1, 0]]]], [[["0", "0", "-4", "-2", "0", "1", "0", "0"], "0", ["halasz_diag", [2, 1, 0]]], [["4", "2", "-12", "-6", "0", "1", "0", "1"], "1", ["halasz_off", [2, 1, 0]]]], [[["0", "0", "-4", "-2", "-2", "1", "0", "0"], "0", ["halasz_diag", [2, 1, 1]]], [["2", "0", "-12", "-6", "-6", "1", "0", "1"], "-1", ["halasz_off", [2, 1, 1]]]], [[["0", "0", "-4", "-4", "0", "1", "0", "0"], "0", ["halasz_diag", [2, 2, 0]]], [["4", "4", "-12", "-12", "0", "1", "0", "1"], "1", ["halasz_off", [2, 2, 0]]]], [[["0", "0", "-4", "-4", "-2", "1", "0", "0"], "0", ["halasz_diag", [2, 2, 1]]], [["2", "2", "-12", "-12", "-6", "1", "0", "1"], "-1", ["halasz_off", [2, 2, 1]]]], [[["0", "0", "-4", "-6", "0", "1", "0", "0"], "0", ["halasz_diag", [2, 3, 0]]], [["4", "6", "-12", "-18", "0", "1", "0", "1"], "1", ["halasz_off", [2, 3, 0]]]], [[["0", "0", "-4", "-8", "0", "1", "0", "0"], "0", ["halasz_diag", [2, 4, 0]]], [["4", "8", "-12", "-24", "0", "1", "0", "1"], "1", ["halasz_off", [2, 4, 0]]]], [[["0", "0", "-4", "-10", "0", "1", "0", "0"], "0", ["halasz_diag", [2, 5, 0]]], [["4", "10", "-12", "-30", "0", "1", "0", "1"], "1", ["halasz_off", [2, 5, 0]]]], [[["0", "0", "-4", "-12", "0", "1", "0", "0"], "0", ["halasz_diag", [2, 6, 0]]], [["4", "12", "-12", "-36", "0", "1", "0", "1"], "1", ["halasz_off", [2, 6, 0]]]], [[["0", "0", "-4", "-14", "0", "1", "0", "0"], "0", ["halasz_diag", [2, 7, 0]]], [["4", "14", "-12", "-42", "0", "1", "0", "1"], "1", ["halasz_off", [2, 7, 0]]]], [[["0", "0", "-4", "-16", "0", "1", "0", "0"], "0", ["halasz_diag", [2, 8, 0]]], [["4", "16", "-12", "-48", "0", "1", "0", "1"], "1", ["halasz_off", [2, 8, 0]]]], [[["0", "0", "-6", "0", "0", "1", "0", "0"], "0", ["halasz_diag", [3, 0, 0]]], [["6", "0", "-18", "0", "0", "1", "0", "1"], "1", ["halasz_off", [3, 0, 0]]]], [[["0", "0", "-6", "-2", "0", "1", "0", "0"], "0", ["halasz_diag", [3, 1, 0]]], [["6", "2", "-18", "-6", "0", "1", "0", "1"], "1", ["halasz_off", [3, 1, 0]]]], [[["0", "0", "-6", "-4", "0", "1", "0", "0"], "0", ["halasz_diag", [3, 2, 0]]], [["6", "4", "-18", "-12", "0", "1", "0", "1"], "1", ["halasz_off", [3, 2, 0]]]], [[["0", "0", "-6", "-6", "0", "1", "0", "0"], "0", ["halasz_diag", [3, 3, 0]]], [["6", "6", "-18", "-18", "0", "1", "0", "1"], "1", ["halasz_off", [3, 3, 0]]]]], "target": ["0", "0", "-2", "-2", "-2", "1", "1/100", "0"], "leaves": [{"path": [[2, 1], [69, 1], [1, 1], [71, 1], [22, 1]], "dual": {"bound": "0", "weights": [[8, "1/48"], [21, "23/2400"], [23, "1/2400"], [55, "47/50"], [86, "1/1200"], [88, "11/400"], [89, "3/100"], [90, "1/600"]]}}, {"path": [[2, 1], [69, 1], [1, 1], [71, 1], [22, 0], [72, 1]], "dual": {"bound": "0", "weights": [[9, "12/575"], [21, "1/100"], [55, "563/575"], [86, "1/2300"], [88, "21/2300"], [89, "21/2300"], [90, "1/575"], [91, "1/2300"]]}}, {"path": [[2, 1], [69, 1], [1, 1], [71, 1], [22, 0], [72, 0]], "dual": {"bound": "-1/576", "weights": [[6, "269/600"], [9, "23/24"], [21, "11/48"], [22, "11/48"], [86, "1/24"], [88, "3/8"], [90, "1/12"], [91, "1/2"]]}}, {"path": [[2, 1], [69, 1], [1, 1], [71, 0], [22, 1], [31, 1], [37, 1], [40, 1]], "dual": {"bound": "0", "weights": [[22, "1/100"], [29, "211/440"], [48, "1/2200"], [86, "1/2200"], [88, "21/1100"], [89, "1/2"], [93, "1/2200"]]}}, {"path": [[2, 1], [69, 1], [1, 1], [71, 0], [22, 1], [31, 1], [37, 1], [40, 0]], "dual": {"bound": "0", "weights": [[22, "1/100"], [29, "959/2000"], [86, "1/2000"], [88, "19/1000"], [89, "1/2"], [92, "1/2000"], [93, "1/2000"]]}}, {"path": [[2, 1], [69, 1], [1, 1], [71, 0], [22, 1], [31, 1], [37, 0], [34, 1]], "dual": {"bound": "0", "weights": [[22, "1/100"], [29, "863/1800"], [86, "1/1800"], [88, "17/900"], [89, "1/2"], [92, "1/1800"], [93, "1/1800"]]}}, {"path": [[2, 1], [69, 1], [1, 1], [71, 0], [22, 1], [31, 1], [37, 0], [34, 0]], "dual": {"bound": "0", "weights": [[22, "1/100"], [29, "767/1600"], [86, "1/1600"], [88, "3/160"], [89, "1/2"], [91, "1/1600"], [93, "1/1600"]]}}, {"path": [[2, 1], [69, 1], [1, 1], [71, 0], [22, 1], [31, 0], [27, 1]], "dual": {"bound": "0", "weights": [[22, "1/100"], [29, "671/1400"], [86, "1/1400"], [88, "13/700"], [89, "1/2"], [91, "1/1400"], [92, "1/1400"]]}}, {"path": [[2, 1], [69, 1], [1, 1], [71, 0], [22, 1], [31, 0], [27, 0]], "dual": {"bound": "0", "weights": [[22, "1/100"], [29, "23/48"], [86, "1/1200"], [88, "11/600"], [89, "1/2"], [90, "1/1200"], [92, "1/1200"]]}}, {"path": [[2, 1], [69, 1], [1, 1], [71, 0], [22, 0]], "dual": {"bound": "-1/576", "weights": [[6, "269/600"], [9, "23/24"], [21, "11/48"], [22, "11/48"], [86, "1/24"], [88, "3/8"], [89, "1/2"], [90, "1/12"]]}}, {"path": [[2, 1], [69, 1], [1, 0], [71, 1], [22, 1], [31, 1], [37, 1], [40, 1]], "dual": {"bound": "-1/2745", "weights": [[6, "12601/54900"], [8, "95/2196"], [23, "263/1098"], [48, "5/549"], [86, "23/2196"], [88, "343/732"], [89, "1/2"], [93, "13/1098"]]}}, {"path": [[2, 1], [69, 1], [1, 0], [71, 1], [22, 1], [31, 1], [37, 1], [40, 0]], "dual": {"bound": "-1/3825", "weights": [[6, "3497/15300"], [8, "29/612"], [23, "73/306"], [86, "7/612"], [88, "95/204"], [89, "1/2"], [92, "2/153"], [93, "1/102"]]}}, {"path": [[2, 1], [69, 1], [1, 0], [71, 1], [22, 1], [31, 1], [37, 0], [34, 1]], "dual": {"bound": "-1/9425", "weights": [[6, "8573/37700"], [8, "79/1508"], [23, "179/754"], [86, "19/1508"], [88, "697/1508"], [89, "1/2"], [92, "4/377"], [93, "11/754"]]}}, {"path": [[2, 1], [69, 1], [1, 0], [71, 1], [22, 1], [31, 1], [37, 0], [34, 0], [57, 1]], "dual": {"bound": "-7/9225", "weights": [[6, "8581/36900"], [8, "1/492"], [23, "179/738"], [86, "1/1476"], [88, "245/492"], [89, "113/246"], [93, "11/738"], [94, "10/369"]]}}, {"path": [[2, 1], [69, 1], [1, 0], [71, 1], [22, 1], [31, 1], [37, 0], [34, 0], [57, 0]], "dual": {"bound": "-4/6975", "weights": [[6, "6121/27900"], [8, "11/186"], [23, "64/279"], [86, "4/279"], [88, "85/186"], [89, "136/279"], [91, "1/62"], [94, "7/279"]]}}, {"path": [[2, 1], [69, 1], [1, 0], [71, 1], [22, 1], [31, 0], [27, 1], [72, 1]], "dual": {"bound": "0", "weights": [[23, "1/100"], [55, "671/700"], [88, "29/1400"], [89, "13/700"], [91, "1/1400"], [92, "1/1400"], [93, "1/1400"]]}}, {"path": [[2, 1], [69, 1], [1, 0], [71, 1], [22, 1], [31, 0], [27, 1], [72, 0]], "dual": {"bound": "-7821/286250", "weights": [[8, "198/5725"], [20, "1/100"], [86, "693/45800"], [87, "693/45800"], [88, "20821/45800"], [91, "1287/22900"], [92, "99/22900"], [93, "20821/45800"]]}}, {"path": [[2, 1], [69, 1], [1, 0], [71, 1], [22, 1], [31, 0], [27, 0], [72, 1]], "dual": {"bound": "0", "weights": [[23, "1/100"], [55, "23/24"], [88, "1/48"], [89, "11/600"], [90, "1/1200"], [92, "1/1200"], [93, "1/1200"]]}}, {"path": [[2, 1], [69, 1], [1, 0], [71, 1], [22, 1], [31, 0], [27, 0], [72, 0]], "dual": {"bound": "-6549/215000", "weights": [[8, "21/34400"], [21, "1/100"], [86, "611/34400"], [87, "611/34400"], [88, "15367/34400"], [90, "7/1376"], [92, "2269/34400"], [93, "15367/34400"]]}}, {"path": [[2, 1], [69, 1], [1, 0], [71, 1], [22, 0], [72, 1]], "dual": {"bound": "-1/576", "weights": [[6, "269/600"], [9, "23/24"], [21, "11/48"], [23, "11/48"], [88, "1/2"], [89, "3/8"], [90, "1/12"], [91, "1/24"]]}}, {"path": [[2, 1], [69, 1], [1, 0], [71, 1], [22, 0], [72, 0]], "dual": {"bound": "-17/504", "weights": [[6, "179/2100"], [9, "1/3"], [21, "2/21"], [86, "1/42"], [87, "1/42"], [88, "3/7"], [90, "2/21"], [91, "3/7"]]}}, {"path": [[2, 1], [69, 1], [1, 0], [71, 0]], "dual": {"bound": "-55/3232", "weights": [[6, "99/10100"], [9, "3/404"], [22, "2/101"], [48, "4/101"], [59, "1/101"], [86, "9/404"], [88, "175/404"], [89, "50/101"]]}}, {"path": [[2, 1], [69, 0]], "dual": {"bound": "-23/1176", "weights": [[6, "19/1225"], [9, "1/49"], [22, "5/196"], [29, "83/196"], [39, "1/196"], [48, "9/196"], [86, "5/196"], [87, "1/2"]]}}, {"path": [[2, 0], [69, 1], [71, 1], [22, 1], [31, 1], [37, 1], [72, 1], [40, 1]], "dual": {"bound": "0", "weights": [[23, "1/100"], [48, "1/2200"], [74, "211/440"], [86, "1/2"], [88, "21/1100"], [92, "1/2200"], [93, "1/2200"]]}}, {"path": [[2, 0], [69, 1], [71, 1], [22, 1], [31, 1], [37, 1], [72, 1], [40, 0]], "dual": {"bound": "0", "weights": [[23, "1/100"], [74, "959/2000"], [86, "1/2"], [88, "19/1000"], [91, "1/2000"], [92, "1/2000"], [93, "1/2000"]]}}, {"path": [[2, 0], [69, 1], [71, 1], [22, 1], [31, 1], [37, 1], [72, 0]], "dual": {"bound": "-1/57", "weights": [[6, "97/300"], [8, "25/38"], [19, "37/114"], [23, "1/114"], [86, "1/2"], [87, "1/114"], [91, "1/57"], [92, "9/19"]]}}, {"path": [[2, 0], [69, 1], [71, 1], [22, 1], [31, 1], [37, 0], [72, 1], [34, 1]], "dual": {"bound": "0", "weights": [[23, "1/100"], [74, "863/1800"], [86, "1/2"], [88, "17/900"], [91, "1/1800"], [92, "1/1800"], [93, "1/1800"]]}}, {"path": [[2, 0], [69, 1], [71, 1], [22, 1], [31, 1], [37, 0], [72, 1], [34, 0]], "dual": {"bound": "0", "weights": [[6, "767/3300"], [23, "8/33"], [86, "1/2"], [88, "5/11"], [90, "1/66"], [92, "1/66"], [93, "1/66"]]}}, {"path": [[2, 0], [69, 1], [71, 1], [22, 1], [31, 1], [37, 0], [72, 0]], "dual": {"bound": "-299/15072", "weights": [[6, "243/15700"], [9, "13/628"], [23, "4/157"], [86, "1/2"], [87, "4/157"], [90, "3/628"], [91, "29/628"], [92, "133/314"]]}}, {"path": [[2, 0], [69, 1], [71, 1], [22, 1], [31, 0], [72, 1], [27, 1]], "dual": {"bound": "0", "weights": [[23, "1/100"], [55, "671/700"], [86, "29/1400"], [88, "13/700"], [90, "1/1400"], [91, "1/1400"], [92, "1/1400"]]}}, {"path": [[2, 0], [69, 1], [71, 1], [22, 1], [31, 0], [72, 1], [27, 0]], "dual": {"bound": "0", "weights": [[23, "1/100"], [55, "23/24"], [86, "1/48"], [88, "11/600"], [89, "1/1200"], [91, "1/1200"], [92, "1/1200"]]}}, {"path": [[2, 0], [69, 1], [71, 1], [22, 1], [31, 0], [72, 0]], "dual": {"bound": "-23/992", "weights": [[6, "69/3100"], [9, "3/124"], [23, "1/31"], [86, "1/2"], [87, "1/31"], [89, "1/124"], [90, "7/124"], [91, "25/62"]]}}, {"path": [[2, 0], [69, 1], [71, 1], [22, 0], [72, 1]], "dual": {"bound": "-1/576", "weights": [[6, "269/600"], [9, "23/24"], [21, "11/48"], [23, "11/48"], [86, "1/2"], [88, "3/8"], [89, "1/12"], [90, "1/24"]]}}, {"path": [[2, 0], [69, 1], [71, 1], [22, 0], [72, 0]], "dual": {"bound": "-11/432", "weights": [[6, "191/900"], [9, "13/18"], [21, "1/6"], [23, "1/18"], [86, "1/2"], [87, "1/18"], [89, "1/9"], [90, "1/3"]]}}, {"path": [[2, 0], [69, 1], [71, 0]], "dual": {"bound": "-23/2016", "weights": [[6, "83/4200"], [9, "1/84"], [23, "5/168"], [39, "1/56"], [48, "1/24"], [86, "1/2"], [87, "5/168"], [88, "23/56"]]}}, {"path": [[2, 0], [69, 0]], "dual": {"bound": "0", "weights": [[20, "1/100"], [29, "67/200"], [58, "33/100"], [87, "67/200"]]}}]}
\end{filecontents*}
\begin{filecontents*}[overwrite]{verify_threefactor_fraction.py}
"""Pure Fraction verifier of 56 fixed-moment Holder identities.
The table is static: no optimization or symbolic algebra is performed.
Branch order: DDD, DDO, DOD, DOO, ODD, ODO, OOD, OOO.
Each target bounds V^2 W^2 U^2 R. Coordinates: T,M,H,K,V,W,U.
"""
from fractions import Fraction as F
from itertools import product

TRIPLES = [((3, 0, 0), (0, 5, 0), (0, 0, 2)), ((3, 0, 0), (0, 4, 0), (0, 0, 2)), ((2, 0, 0), (0, 5, 0), (0, 0, 3)), ((2, 0, 0), (0, 4, 0), (0, 0, 3)), ((2, 1, 0), (0, 6, 0), (0, 0, 2)), ((2, 1, 0), (0, 7, 0), (0, 0, 2)), ((2, 0, 0), (0, 6, 0), (0, 1, 2))]

TARGETS = [
    [
        ('0', '9/10', '1', '1', '1/5', '0', '0'),
        ('7/15', '1', '1', '1/30', '0', '0', '0'),
        ('1/6', '1', '1/12', '1', '0', '0', '0'),
        ('2/3', '1', '0', '1/30', '0', '0', '0'),
        ('3/10', '1/20', '1', '1', '0', '0', '0'),
        ('4/5', '0', '1', '1/30', '0', '0', '0'),
        ('1/2', '0', '1/12', '1', '0', '0', '0'),
        ('1', '1/20', '0', '0', '0', '0', '0'),
    ],
    [
        ('0', '3/4', '1', '1', '1/2', '0', '0'),
        ('5/12', '1', '1', '1/12', '0', '0', '0'),
        ('1/6', '1', '1/6', '1', '0', '0', '0'),
        ('2/3', '1', '1/6', '0', '0', '0', '0'),
        ('1/4', '1/8', '1', '1', '0', '0', '0'),
        ('3/4', '1/8', '1', '0', '0', '0', '0'),
        ('1/2', '0', '1/6', '1', '0', '0', '0'),
        ('1', '1/8', '0', '0', '0', '0', '0'),
    ],
    [
        ('0', '1', '5/6', '1', '0', '1/3', '0'),
        ('3/10', '1', '1', '1/20', '0', '0', '0'),
        ('1/6', '1', '1/12', '1', '0', '0', '0'),
        ('1/2', '1', '1/12', '0', '0', '0', '0'),
        ('7/15', '1/30', '1', '1', '0', '0', '0'),
        ('4/5', '0', '1', '1/20', '0', '0', '0'),
        ('2/3', '0', '1/12', '1', '0', '0', '0'),
        ('1', '0', '0', '1/20', '0', '0', '0'),
    ],
    [
        ('0', '1', '2/3', '1', '0', '2/3', '0'),
        ('1/4', '1', '1', '1/8', '0', '0', '0'),
        ('1/6', '1', '1/6', '1', '0', '0', '0'),
        ('1/2', '1', '1/6', '0', '0', '0', '0'),
        ('5/12', '1/12', '1', '1', '0', '0', '0'),
        ('3/4', '0', '1', '1/8', '0', '0', '0'),
        ('2/3', '0', '1/6', '1', '0', '0', '0'),
        ('1', '0', '0', '1/8', '0', '0', '0'),
    ],
    [
        ('0', '1', '1/2', '1', '0', '1', '0'),
        ('5/12', '1', '1', '1/12', '0', '0', '0'),
        ('0', '1', '1/2', '1', '0', '1', '0'),
        ('1/2', '1', '2/3', '1/36', '0', '0', '0'),
        ('5/12', '1/12', '13/24', '1', '0', '0', '0'),
        ('11/12', '1/12', '13/24', '0', '0', '0', '0'),
        ('1/2', '1/36', '13/72', '1', '0', '0', '0'),
        ('1', '0', '0', '1/12', '0', '0', '0'),
    ],
    [
        ('0', '1', '1/2', '1', '0', '1', '0'),
        ('3/7', '1', '1', '1/14', '0', '0', '0'),
        ('0', '1', '1/2', '1', '0', '1', '0'),
        ('1/2', '1', '2/3', '1/42', '0', '0', '0'),
        ('3/7', '1/14', '15/28', '1', '0', '0', '0'),
        ('13/14', '1/14', '15/28', '0', '0', '0', '0'),
        ('1/2', '1/42', '5/28', '1', '0', '0', '0'),
        ('1', '0', '0', '1/14', '0', '0', '0'),
    ],
    [
        ('0', '1', '1/2', '1', '0', '1', '0'),
        ('5/12', '1', '13/24', '1/12', '0', '0', '0'),
        ('0', '1', '1/2', '1', '0', '1', '0'),
        ('1/2', '1', '13/72', '1/36', '0', '0', '0'),
        ('5/12', '1/12', '1', '1', '0', '0', '0'),
        ('11/12', '1/12', '1/2', '0', '0', '0', '0'),
        ('1/2', '1/36', '2/3', '1', '0', '0', '0'),
        ('1', '1/12', '0', '0', '0', '0', '0'),
    ],
]

WEIGHTS = [
    [
        ('3/10', '1/5', '1/2'),
        ('1/3', '1/5', '9/20', '1/60'),
        ('1/3', '3/20', '1/60', '1/2'),
        ('1/3', '1/5', '0', '9/20', '1/60'),
        ('17/60', '1/60', '1/5', '1/2'),
        ('1/3', '0', '1/5', '9/20', '1/60'),
        ('1/3', '0', '3/20', '1/60', '1/2'),
        ('17/60', '1/60', '1/5', '0', '1/2', '0'),
    ],
    [
        ('1/4', '1/4', '1/2'),
        ('1/3', '1/4', '3/8', '1/24'),
        ('1/3', '1/8', '1/24', '1/2'),
        ('1/3', '1/8', '1/24', '1/2', '0'),
        ('5/24', '1/24', '1/4', '1/2'),
        ('5/24', '1/24', '1/4', '1/2', '0'),
        ('1/3', '0', '1/8', '1/24', '1/2'),
        ('5/24', '1/24', '1/4', '0', '1/2', '0'),
    ],
    [
        ('1/2', '1/6', '1/3'),
        ('1/2', '1/5', '17/60', '1/60'),
        ('1/2', '3/20', '1/60', '1/3'),
        ('1/2', '3/20', '1/60', '1/3', '0'),
        ('9/20', '1/60', '1/5', '1/3'),
        ('1/2', '0', '1/5', '17/60', '1/60'),
        ('1/2', '0', '3/20', '1/60', '1/3'),
        ('1/2', '0', '1/5', '0', '17/60', '1/60'),
    ],
    [
        ('1/2', '1/6', '1/3'),
        ('1/2', '1/4', '5/24', '1/24'),
        ('1/2', '1/8', '1/24', '1/3'),
        ('1/2', '1/8', '1/24', '1/3', '0'),
        ('3/8', '1/24', '1/4', '1/3'),
        ('1/2', '0', '1/4', '5/24', '1/24'),
        ('1/2', '0', '1/8', '1/24', '1/3'),
        ('1/2', '0', '1/4', '0', '5/24', '1/24'),
    ],
    [
        ('1/2', '0', '1/2'),
        ('1/2', '1/12', '3/8', '1/24'),
        ('1/2', '0', '0', '1/2'),
        ('1/2', '0', '1/36', '11/24', '1/72'),
        ('3/8', '1/24', '1/12', '1/2'),
        ('3/8', '1/24', '1/12', '1/2', '0'),
        ('11/24', '1/72', '0', '1/36', '1/2'),
        ('1/2', '0', '1/12', '0', '3/8', '1/24'),
    ],
    [
        ('1/2', '0', '1/2'),
        ('1/2', '1/14', '11/28', '1/28'),
        ('1/2', '0', '0', '1/2'),
        ('1/2', '0', '1/42', '13/28', '1/84'),
        ('11/28', '1/28', '1/14', '1/2'),
        ('11/28', '1/28', '1/14', '1/2', '0'),
        ('13/28', '1/84', '0', '1/42', '1/2'),
        ('1/2', '0', '1/14', '0', '11/28', '1/28'),
    ],
    [
        ('1/2', '0', '1/2'),
        ('1/2', '1/12', '3/8', '1/24'),
        ('1/2', '0', '0', '1/2'),
        ('1/2', '0', '1/36', '11/24', '1/72'),
        ('3/8', '1/24', '1/12', '1/2'),
        ('3/8', '1/24', '1/12', '1/2', '0'),
        ('11/24', '1/72', '0', '1/36', '1/2'),
        ('3/8', '1/24', '1/12', '0', '1/2', '0'),
    ],
]

def verify():
    cases = list(product((0, 1), repeat=3))
    assert len(TRIPLES) == len(TARGETS) == len(WEIGHTS) == 7
    for ci, triples in enumerate(TRIPLES):
        assert len(triples) == 3
        assert len(TARGETS[ci]) == len(WEIGHTS[ci]) == 8
        for bi, off in enumerate(cases):
            rows = []
            for p, is_off in zip(triples, off):
                assert len(p) == 3 and all(isinstance(x,int) and x>=0 for x in p)
                assert sum(p) > 0
                if not is_off:
                    # Diagonal: R <= L A^-2.
                    rows.append((0, *p, *(-2*x for x in p)))
                else:
                    # Mean off-diagonal: R <= T A^-2.
                    rows.append((1, 0, 0, 0, *(-2*x for x in p)))
                    # Halasz off-diagonal: R <= LT A^-6.
                    rows.append((1, *p, *(-6*x for x in p)))
            weights = [F(x) for x in WEIGHTS[ci][bi]]
            assert len(rows) == len(weights)
            assert all(x>=0 for x in weights) and sum(weights)==1
            answer = [sum(w*r[j] for w,r in zip(weights,rows)) for j in range(7)]
            for j in (4,5,6): answer[j] += 2
            assert answer == [F(x) for x in TARGETS[ci][bi]], (ci,bi)
            assert 0 <= answer[0] <= 1  # Critical-height reduction is valid.
            assert all(x>=0 for x in answer[4:])
    print("Verified all 56 monomial identities using nonnegative rational Holder weights.")

if __name__ == '__main__':
    verify()
\end{filecontents*}
\begin{filecontents*}[overwrite]{region_stability_probe.py}
from fractions import Fraction as F
from itertools import product
from math import gcd
from functools import reduce
from collections import defaultdict

t0=F(1,24)
tII=F(1,22)

# A*M+B*H <= b0+b1*theta; reversing all signs permits lower bounds.
base=[]
def add(name,A,B,b0,b1,atIII=None,atII=None):
    b0,b1=F(b0),F(b1)
    if atIII is not None:
        assert b0+b1*t0==F(atIII),(name,b0+b1*t0,atIII)
    if atII is not None:
        assert b0+b1*tII==F(atII),(name,b0+b1*tII,atII)
    base.append((name,A,B,b0,b1))
def lo(name,A,B,b0,b1,atIII=None,atII=None):
    add(name,-A,-B,-F(b0),-F(b1),None if atIII is None else -F(atIII),None if atII is None else -F(atII))

# Li III Lemma 3.2, matching Li II cases 1,3,4,6,7,8,9.
add('1a',1,1,F(15,29),F(14,29),F(187,348),F(172,319))
lo('1b',0,1,F(2,11),-F(2,11),F(23,132),F(21,121))
add('1c',29,-1,9,20,F(59,6),F(109,11))
lo('1d',1,0,F(6,19),-F(6,19),F(23,76),F(63,209))
add('1e',-1,29,5,24,6,F(67,11))
lo('1f',12,11,6,-6,F(23,4),F(63,11))

add('2a',1,1,F(6,11),F(5,11),F(149,264),F(137,242))
lo('2b',0,1,F(1,5),-F(1,5),F(23,120),F(21,110))
add('2c',6,1,2,4,F(13,6),F(24,11))
lo('2d',1,0,F(2,7),-F(2,7),F(23,84),F(3,11))
add('2e',1,0,0,8,F(1,3),F(4,11))
add('2f',1,8,2,6,F(9,4),F(25,11))
lo('2g',6,5,3,-3,F(23,8),F(63,22))

# Derive affine coefficients for this face from the verified fixed-case values.
bIII,bII=F(53,76),F(146,209)
der=(bII-bIII)/(tII-t0)
add('3a',1,1,bIII-der*t0,der,bIII,bII)
lo('3b',0,1,F(2,11),-F(2,11),F(23,132),F(21,121))
add('3c',12,1,6,6,F(25,4),F(69,11))
lo('3d',1,0,F(14,29),-F(14,29),F(161,348),F(147,319))
add('3e',-1,19,3,16,F(11,3),F(41,11))
lo('3f',12,11,8,-8,F(23,3),F(84,11))

lo('4a',1,1,1,-8,F(2,3),F(7,11))
add('4b',1,1,F(5,7),F(2,7),F(61,84),F(8,11))
lo('4c',0,1,F(1,5),-F(1,5),F(23,120),F(21,110))
add('4d',6,1,3,3,F(25,8),F(69,22))
lo('4e',1,0,F(5,11),-F(5,11),F(115,264),F(105,242))
add('4f',-1,7,1,6,F(5,4),F(14,11))
lo('4g',6,5,4,-4,F(23,6),F(42,11))

lo('5a',1,1,1,-12,F(1,2))
add('5b',1,1,F(6,11),F(5,11),F(149,264),F(137,242))
add('5c',35,23,17,18,F(71,4),F(196,11))
lo('5d',2,1,F(10,11),-F(10,11),F(115,132),F(105,121))
add('5e',2,13,2,22,F(35,12),3)
lo('5f',70,59,36,-36,F(69,2),F(378,11))

add('6a',1,1,F(7,13),F(6,13),F(29,52),F(80,143))
add('6b',41,27,20,21,F(167,8),F(461,22))
lo('6c',2,1,F(12,13),-F(12,13),F(23,26),F(126,143))
add('6d',2,15,2,26,F(37,12),F(35,11))
lo('6e',82,69,42,-42,F(161,4),F(441,11))

add('7a',2,1,F(12,11),F(10,11),F(149,132),F(137,121))
add('7b',70,11,34,36,F(71,2),F(392,11))
lo('7c',1,0,F(5,11),-F(5,11),F(115,264),F(105,242))
add('7d',1,0,0,12,F(1,2))
add('7e',1,6,1,11,F(35,24),F(3,2))
lo('7f',35,12,18,-18,F(69,4),F(189,11))

# Watt Type I.
add('Wa',2,1,1,1,F(25,24))
add('Wb',4,6,3,1,F(73,24))
add('Wc',0,4,1,3,F(9,8))
add('Wd',1,0,F(1,2),F(1,2),F(25,48))
add('We',0,1,F(1,8),-F(1,8),F(23,192))

# Harman parametric regions; independently verified rational certificates
# supply their analytic proof. This file checks only boundary geometry.
lo('H0a',0,1,0,2,F(1,12))
add('H0b',-1,2,0,2,F(1,12))
add('H0c',1,0,0,4,F(1,6))
add('H1a',4,3,2,2,F(25,12))
lo('H1b',4,3,2,-2,F(23,12))
lo('H1c',0,1,0,2,F(1,12))
add('H1d',0,1,0,4,F(1,6))
add('H2a',4,1,2,2,F(25,12))
lo('H2b',4,1,2,-2,F(23,12))
lo('H2c',0,1,0,2,F(1,12))
add('H2d',0,1,0,4,F(1,6))

# Fixed geometric faces in the two-variable predicates.
a0=F(11,126)
lo('DomM',1,0,a0,0)
lo('DomH',0,1,a0,0)
add('DomMhalf',1,0,F(1,2),0)
add('DomOrder',-1,1,0,0)
add('DomResidual',1,2,1,0)

print('primitive_faces',len(base),'max_integer_coeff',max(max(abs(a),abs(b)) for _,a,b,_,_ in base),'max_rhs_theta_slope',max(abs(s) for _,_,_,_,s in base))
print('parametric_case3a',[(n,str(b0),str(b1)) for n,_,_,b0,b1 in base if n=='3a'])

# Four symbols = t1,t2,t3,beta with beta=1-t1-t2-t3.
# Assign each symbol to M, H, or neither. Including unused symbols is an
# intentional superset of all groupings actually used in the finite sieve.
symbols=[(F(0),(1,0,0)),(F(0),(0,1,0)),(F(0),(0,0,1)),(F(1),(-1,-1,-1))]
forms=[]
for assign in product(range(3),repeat=4):
    if 1 not in assign or 2 not in assign: continue
    mc=sum(symbols[i][0] for i,a in enumerate(assign) if a==1)
    hc=sum(symbols[i][0] for i,a in enumerate(assign) if a==2)
    mv=tuple(sum(symbols[i][1][j] for i,a in enumerate(assign) if a==1) for j in range(3))
    hv=tuple(sum(symbols[i][1][j] for i,a in enumerate(assign) if a==2) for j in range(3))
    for name,a,b,b0,b1 in base:
        av=tuple(a*x+b*y for x,y in zip(mv,hv))
        rhs0=b0-a*mc-b*hc
        forms.append((name,assign,av,rhs0,b1))

# The actual p3-domain for each parent pair, plus role-reversal boundaries.
for av,b0 in [((-1,0,0),-a0),((0,-1,0),-a0),((0,0,-1),-a0),
              ((1,0,0),F(1,2)),((-1,1,0),0),((0,-1,1),0),
              ((1,2,0),1),((1,1,2),1),((1,1,1),1-a0)]:
    forms.append(('geometry',(),av,F(b0),F(0)))

groups=defaultdict(list)
constant=[]
for name,assign,av,b0,b1 in forms:
    g=reduce(gcd,map(abs,av))
    if not g:
        if b0+b1*t0==0 and b1:
            constant.append((name,assign,str(b0),str(b1)))
        continue
    sg=1 if next(v for v in av if v)!=0 and next(v for v in av if v)>0 else -1
    divisor=g*sg
    canon=tuple(v//divisor for v in av)
    groups[(canon,(b0+b1*t0)/divisor)].append((name,assign,b0/divisor,b1/divisor))

collisions=[]
for (normal,at0),items in groups.items():
    coeffs=set((x[2],x[3]) for x in items)
    if len(coeffs)>1:
        collisions.append((normal,str(at0),[(n,a,str(b0),str(b1)) for n,a,b0,b1 in items]))
print('substituted_forms',len(forms),'max_integer_coeff_substituted',max(max(map(abs,x[2])) for x in forms))
print('constant_zero_faces',constant)
print('coincident_parallel_faces_nonidentical',len(collisions))
for item in collisions[:40]:print(item)

# These are exact assertions about the complete finite superset above,
# not numerical plausibility checks. The half-face collisions must still
# be shown harmless on each actual universally quantified parent range.
assert len(base)==64
assert len(forms)==3209
assert not constant
assert len(collisions)==7
for normal,at0,items in collisions:
    assert F(at0)==F(1,2)
    assert set(normal)<=set((0,1)) and any(normal)
    assert {x[0] for x in items} <= {'5a','7d','DomMhalf','DomOrder','geometry'}

# Nearest distinct parameter at which two parallel primitive faces meet.
by_normal=defaultdict(set)
for name,assign,av,b0,b1 in forms:
    g=reduce(gcd,map(abs,av))
    if not g: continue
    sg=1 if next(v for v in av if v)>0 else -1
    divisor=g*sg
    by_normal[tuple(v//divisor for v in av)].add((b0/divisor,b1/divisor))
nearest=[]
for normal,coeffset in by_normal.items():
    coeffs=list(coeffset)
    for i,(a,b) in enumerate(coeffs):
        for c,d in coeffs[i+1:]:
            if b==d:continue
            cross=(c-a)/(b-d)
            if F(1,25)<=cross<t0:
                nearest.append((t0-cross,cross,normal,(a,b),(c,d)))
nearest.sort()
print('nearest_distinct_parallel_parameter_contact',[(str(x[0]),str(x[1]),x[2]) for x in nearest[:5]])
assert nearest and nearest[0][0]==F(1,600) and nearest[0][1]==F(1,25)
print('PASS: all parameter-only degeneracies are absent; all seven endpoint collisions are the audited half-faces.')
\end{filecontents*}
\begin{filecontents*}[overwrite]{verify_negative_child_coverage.py}
"""Exact rational coverage of Li III's original universal negative-child intervals.
This checks all arrangement cells, not a numerical grid. Stdlib only.
"""
from fractions import Fraction as F
from pathlib import Path
import contextlib,io,runpy
with contextlib.redirect_stdout(io.StringIO()):
    B=runpy.run_path(str(Path(__file__).with_name('region_stability_probe.py')))['base']
theta=F(1,24); cutoff=F(11,126)

def region(keys,geometry=False):
    out=[(F(a),F(b),b0+b1*theta) for n,a,b,b0,b1 in B if n in keys]
    if geometry:out.extend([(F(-1),F(1),F(0)),(F(1),F(0),F(1,2))])
    return out
C=[region({n for n,*_ in B if n[0]==str(i)}) for i in range(1,8)]
H=[region({n for n,*_ in B if n.startswith(p)},True) for p in ('H1','H2')]
W=[region(set(('Wa','Wb','Wc'))),region(set(('Wd','We')))]

def at_m(reg,m):
    low=None;high=None
    for A,B,c in reg:
        if B>0:high=min(high,(c-A*m)/B) if high is not None else (c-A*m)/B
        elif B<0:low=max(low,(c-A*m)/B) if low is not None else (c-A*m)/B
        elif A*m>c:return None
    low=max(cutoff,low) if low is not None else cutoff
    high=min(F(1),high) if high is not None else F(1)
    return (low,high) if low<=high else None

def verify(name,ml,mh,qa,qb,regs):
    # Child top is qa*m+qb. All h-endpoints are affine functions of m.
    lines={(F(0),cutoff),(qa,qb)}
    critical={ml,mh}
    for reg in regs:
        for A,B,c in reg:
            if B:lines.add((-A/B,c/B))
            elif A:
                x=c/A
                if ml<=x<=mh:critical.add(x)
    lines=list(lines)
    for j,(a,b) in enumerate(lines):
        for c,d in lines[:j]:
            if a!=c:
                x=(d-b)/(a-c)
                if ml<=x<=mh:critical.add(x)
    critical=sorted(critical)
    tests=critical+[(x+y)/2 for x,y in zip(critical,critical[1:])]
    for m in tests:
        top=qa*m+qb
        if top<cutoff:continue
        spans=sorted(s for reg in regs if (s:=at_m(reg,m)) is not None)
        covered=cutoff
        for lo,hi in spans:
            if hi<covered:continue
            if lo>covered:break
            covered=max(covered,hi)
            if covered>=top:break
        assert covered>=top,(name,str(m),str(covered),str(top),spans)
    print(name,'verified',len(critical)-1,'open m-cells and',len(critical),'endpoints')

parents=[
 ('U0201','139/336','605/1456',-1,'29/52'),
 ('U0202','605/1456','17/40','-41/27','167/216'),
 ('U0203','17/40','51/112','-4/3','25/36'),
 ('U0204','461/1008','11/24','-1/6','35/144'),
 ('U0205','11/24','23/48',0,'1/6'),
 ('U0206','23/48','265/552',-4,'25/12'),
 ('U0207','265/552','4727/9816','-1/6','35/144'),
 ('U0208','4727/9816','151/312','-70/11','71/22'),
 ('U0209','151/312','503/1008',-4,'25/12')]
for name,lo,hi,a,b in parents:verify(name,F(lo),F(hi),F(a),F(b),C+H)
verify('U0303',F(115,264),F(11,24),F(-1,6),F(35,144),[C[6]]+W)
print('All original full negative-child covers verified over the rational arrangement.')
\end{filecontents*}
\begin{filecontents*}[overwrite]{verify_quantitative_sieve_stability.py}
"""Exact inventory and constants for a microscopic finite-sieve decrement.

This is a geometry certificate, not a new quadrature certificate.
The four Li III integral inequalities remain external inputs.
Use together with verify_negative_child_coverage.py and the analytic
certificates.  Standard library only; all arithmetic is rational.
"""
from fractions import Fraction as F
from itertools import product, combinations
from pathlib import Path
from math import comb
import contextlib, io, runpy, json

with contextlib.redirect_stdout(io.StringIO()):
    source = runpy.run_path(str(Path(__file__).with_name('region_stability_probe.py')))
base, raw = source['base'], source['forms']
t0, a = F(1,24), F(11,126)
rho, geometric_d, height_e = F(1,10**20), F(2,10**23), F(1,10**23)
final_Delta = geometric_d-height_e
assert final_Delta == F(1,10**23)

# 1. Complete negative-child primitive inventory: all disjoint nonempty
# M,H assignments of t1,t2,t3,beta=1-t1-t2-t3, plus the nine fixed
# child/parent geometry faces.  Actual negative covers are a subset.
assert len(raw)==3209
negative = set((tuple(F(x) for x in av), b0, b1)
               for name, assignment, av, b0, b1 in raw)
assert all(all(x.denominator==1 and abs(x)<=82 for x in av)
           and abs(b1)<=42 for av,b0,b1 in negative)

# A face is av.t <= b0+b1*theta.  If av[2] is nonzero its child
# endpoint is e0 + e1*theta + ep.(t1,t2).  Every change in the Boolean
# full-fibre coverage test is either a zero of a parent-only face or
# a comparison of two such endpoints.  This directly inventories all
# one-variable elimination events, without counting Boolean cells.
parent_only=[]
endpoints=set()
for av,b0,b1 in negative:
    if not av[2]:
        parent_only.append((tuple(-x for x in av[:2]),b0,b1))
    else:
        endpoints.add((tuple(-x/av[2] for x in av[:2]),b0/av[2],b1/av[2]))
endpoints=list(endpoints)
event_count=len(parent_only)+comb(len(endpoints),2)
minimum_normal=F(1,82**2)
maximum_slope=F(84)
half_contacts=0
constant_root_distances=[]

def check_event(normal,b0,b1):
    global half_contacts
    assert abs(b1)<=maximum_slope
    nz=[abs(x) for x in normal if x]
    if nz:
        assert max(nz)>=minimum_normal
        return
    if not b1:
        return  # a fixed sign or an identity, independent of theta
    root=-b0/b1
    if root==t0:
        half_contacts+=1
        return
    if root<t0:
        constant_root_distances.append(t0-root)
        assert t0-root>=F(1,600)

for normal,b0,b1 in parent_only:
    check_event(normal,b0,b1)
for (p,b0,b1),(q,c0,c1) in combinations(endpoints,2):
    check_event(tuple(x-y for x,y in zip(p,q)),b0-c0,b1-c1)
assert half_contacts>0

# The imported exact checker classifies EVERY coincident primitive face
# at theta0: all seven families are subset sums equal to 1/2 and involve
# only C5a, C7d, and fixed geometry.  They are inactive in the actual
# negative covers for the following independently checked reasons.
assert len(source['collisions'])==7
for normal,value,items in source['collisions']:
    assert F(value)==F(1,2)
    assert {item[0] for item in items} <= {'5a','7d','DomMhalf','DomOrder','geometry'}
# U0201--09: M=s, H=t3; C5a and C7d have these fixed margins.
assert F(139,336)+a-F(1,2)==F(1,1008)
assert F(1,2)-F(503,1008)==F(1,1008)
assert 12*geometric_d < F(1,2016)
# U0210--12 use C1--C3, so neither moving half-face occurs.
# U0301--02 and U04 use Watt, likewise neither occurs.
# U0303 uses C7 with s<=11/24; C7d has fixed margin at theta0.
assert F(1,2)-F(11,24)==F(1,24)
assert 12*geometric_d < F(1,48)

# 2. All terminal good predicates, including the reversed quadruple.
# An entry is (constant, coefficient vector) for one supplied exponent.
maps={
 'parent':[(0,(1,0)),(0,(0,1))],
 'triple':[(0,(1,0,0)),(0,(0,1,0)),(0,(0,0,1))],
 'quadruple':[(0,(1,0,0,0)),(0,(0,1,0,0)),
              (0,(0,0,1,0)),(0,(0,0,0,1))],
 'reversed_quadruple':[(1,(-1,-1,-1,0)),(0,(0,1,0,0)),
                       (0,(0,0,1,0)),(0,(0,0,0,1))],
}
terminal_counts={}
for map_name,symbols in maps.items():
    n=len(symbols[0][1]); count=0
    for assignment in product(range(3),repeat=len(symbols)):
        if 1 not in assignment or 2 not in assignment: continue
        mc=sum(F(symbols[i][0]) for i,x in enumerate(assignment) if x==1)
        hc=sum(F(symbols[i][0]) for i,x in enumerate(assignment) if x==2)
        mv=tuple(sum(symbols[i][1][j] for i,x in enumerate(assignment) if x==1) for j in range(n))
        hv=tuple(sum(symbols[i][1][j] for i,x in enumerate(assignment) if x==2) for j in range(n))
        for name,A,B,b0,b1 in base:
            av=tuple(A*x+B*y for x,y in zip(mv,hv))
            assert any(av)  # no parameter-only terminal predicate
            assert all(isinstance(x,int) and abs(x)<=82 for x in av)
            assert abs(b1)<=42
            count+=1
    terminal_counts[map_name]=count
assert terminal_counts=={'parent':128,'triple':768,'quadruple':3200,'reversed_quadruple':3200}
terminal_count=sum(terminal_counts.values())
assert terminal_count==7296

# 3. Fixed source parent boundaries.  There are exactly 27 displayed
# components: 7 U01, 3 V (including the empty printed V0), 12 U02,
# 3 U03, 2 U04.  Each has <=7 component inequalities and the five
# common Omega inequalities.  Exclusions of earlier components are
# Boolean operations and do not add faces.  The following counts are
# the displayed inequalities before the common five are appended.
component_counts=[6,7,6,7,6,5,6, 3,4,4,
                  4,4,4,4,4,4,4,4,4,6,7,6, 3,3,3, 5,3]
assert len(component_counts)==27 and max(component_counts)<=7
source_parent_faces=sum(component_counts)+5*27
assert source_parent_faces==261 and source_parent_faces<=27*12
# Normalize these fixed source faces by clearing only LHS denominators.
# Every nonzero LHS then has integer coefficients, hence normal >=1.
# Their theta slope is zero.  They therefore obey the common strip bound.

# Four final integration domains: Omega contributes five faces each;
# R1 adds six, R2 adds five, R3 adds two, I0 adds none, total 33.
# Every LHS is integral.  Role reversal already appears in the fourth
# map above and adds no new integration variable or nonlinear boundary.
fixed_domain_faces=4*5+6+5+2
assert fixed_domain_faces==33
total_exact_count=event_count+terminal_count+source_parent_faces+fixed_domain_faces
raw_count_bound=3209**2+3209+7296+27*12+40
assert total_exact_count<=raw_count_bound<11_000_000
generous_face_bound=11_000_000
assert raw_count_bound<generous_face_bound

# 4. Uniform comparator cost.  In dimensions 2--4, inside [0,1/2]^d,
# |g|<=2rho has volume <=4rho/minimum_normal: integrate along one
# coordinate whose coefficient has the stated magnitude; all other
# side lengths are <=1.  We pay the full union-bound cost separately
# for five integrals: discarded positive parent mass plus I1,I2,I0,I3.
assert a**-5 < 200000
number_of_comparator_integrals=5
density_bound=200000
strip_cost=(number_of_comparator_integrals*density_bound*
            generous_face_bound*4*rho/minimum_normal)
assert strip_cost == F(18491,6250000)
assert strip_cost < F(1,200)
assert maximum_slope*geometric_d < rho/4
assert geometric_d < F(1,12096)
assert geometric_d < F(1,600)
assert 0<height_e<geometric_d

# Closed Harman strips at theta_a=t0-d are interior strips at
# theta'=theta_a+e/2.  The actual interval exponent is theta_a+e.
theta_a=t0-geometric_d
theta_prime=theta_a+height_e/2
theta_star=theta_a+height_e
assert F(1,25)<theta_a<theta_prime<theta_star<t0
assert a-2*theta_prime>=F(1,252)
assert t0-theta_star==final_Delta
cstar=1/(1-theta_star)
cnew=F(24,23)-final_Delta
assert 1<cstar<cnew<F(24,23)
assert cnew-cstar == ((47*final_Delta-552*final_Delta**2)/
                     (23*(23+24*final_Delta)))
assert final_Delta<F(47,552)

report={
 'negative_raw_forms':len(raw),
 'negative_unique_forms':len(negative),
 'child_endpoints':len(endpoints),
 'parent_only_faces':len(parent_only),
 'negative_events':event_count,
 'terminal_counts':terminal_counts,
 'source_parent_faces':source_parent_faces,
 'fixed_domain_faces':fixed_domain_faces,
 'total_event_faces':total_exact_count,
 'raw_proven_count_bound':raw_count_bound,
 'generous_count_used':generous_face_bound,
 'minimum_nonzero_normal':str(minimum_normal),
 'maximum_theta_slope':str(maximum_slope),
 'inactive_half_contact_event_count':half_contacts,
 'nearest_other_constant_root_distance':str(min(constant_root_distances)),
 'erosion_rho':str(rho),
 'geometric_shift':str(geometric_d),
 'height_slack':str(height_e),
 'final_Delta':str(final_Delta),
 'mills_threshold_cstar':str(cstar),
 'explicit_mills_cnew':str(cnew),
 'mills_order_verified':'1 < cstar < cnew < 24/23',
 'mills_gap_identity_verified':'cnew-cstar=(47*Delta-552*Delta^2)/(23*(23+24*Delta))',
 'total_integral_strip_cost_bound':str(strip_cost),
 'strip_cost_less_than':'0.005',
 'status':'PASS; conditional on the four stated source integral inequalities and the separately certified analytic lemmas',
}
print(json.dumps(report,indent=2))
\end{filecontents*}
\begin{filecontents*}[overwrite]{quantitative_sieve_stability_report.json}
{
  "negative_raw_forms": 3209,
  "negative_unique_forms": 1503,
  "child_endpoints": 989,
  "parent_only_faces": 431,
  "negative_events": 488997,
  "terminal_counts": {
    "parent": 128,
    "triple": 768,
    "quadruple": 3200,
    "reversed_quadruple": 3200
  },
  "source_parent_faces": 261,
  "fixed_domain_faces": 33,
  "total_event_faces": 496587,
  "raw_proven_count_bound": 10308550,
  "generous_count_used": 11000000,
  "minimum_nonzero_normal": "1/6724",
  "maximum_theta_slope": "84",
  "inactive_half_contact_event_count": 12,
  "nearest_other_constant_root_distance": "1/600",
  "erosion_rho": "1/100000000000000000000",
  "geometric_shift": "1/50000000000000000000000",
  "height_slack": "1/100000000000000000000000",
  "final_Delta": "1/100000000000000000000000",
  "mills_threshold_cstar": "300000000000000000000000/287500000000000000000003",
  "explicit_mills_cnew": "2399999999999999999999977/2300000000000000000000000",
  "mills_order_verified": "1 < cstar < cnew < 24/23",
  "mills_gap_identity_verified": "cnew-cstar=(47*Delta-552*Delta^2)/(23*(23+24*Delta))",
  "total_integral_strip_cost_bound": "18491/6250000",
  "strip_cost_less_than": "0.005",
  "status": "PASS; conditional on the four stated source integral inequalities and the separately certified analytic lemmas"
}
\end{filecontents*}
\begin{filecontents*}[overwrite]{li_III_i0_convention_check.py}
"""Exploratory quadrature of Li III (24), with both V0 conventions.

Geometry is assembled from rational half-planes. Buchstab integration
and the final quadrature are numerical, not rigorous interval bounds.
"""
from fractions import Fraction as F
from pathlib import Path
from itertools import combinations
import json,math
import numpy as np
from scipy.integrate import solve_ivp, quad

a=F(11,126);theta=F(1,24)
domain=[(-1,0,-a),(1,0,F(1,2)),(0,-1,-a),(-1,1,0),(1,2,1)]

def clip(poly,line):
    A,B,C=map(F,line)
    if not poly:return []
    ans=[]
    for p,q in zip(poly,poly[1:]+poly[:1]):
        vp=A*p[0]+B*p[1]-C;vq=A*q[0]+B*q[1]-C
        if vp<=0:ans.append(p)
        if (vp<0 and vq>0) or (vp>0 and vq<0):
            t=vp/(vp-vq)
            ans.append((p[0]+t*(q[0]-p[0]),p[1]+t*(q[1]-p[1])))
    clean=[]
    for p in ans:
        if not clean or p!=clean[-1]:clean.append(p)
    if len(clean)>1 and clean[-1]==clean[0]:clean.pop()
    return clean

def polygon(lines):
    p=[(F(0),F(0)),(F(1),F(0)),(F(1),F(1)),(F(0),F(1))]
    for L in domain+lines:p=clip(p,L)
    return p

def regions(v0cap):
    # Direct transcription of Section 7, independent of the parametric
    # stability code used elsewhere in the project.
    out={
        'U0101':[(1,1,F(187,348)),(0,-1,-F(23,132)),(29,-1,F(59,6)),(-1,0,-F(23,76)),(-1,29,6),(-12,-11,-F(23,4))],
        'U0102':[(1,1,F(149,264)),(0,-1,-F(23,120)),(6,1,F(13,6)),(-1,0,-F(23,84)),(1,0,F(1,3)),(1,8,F(9,4)),(-6,-5,-F(23,8))],
        'U0103':[(1,1,F(53,76)),(0,-1,-F(23,132)),(12,1,F(25,4)),(-1,0,-F(161,348)),(-1,19,F(11,3)),(-12,-11,-F(23,3))],
        'U0104':[(-1,-1,-F(2,3)),(1,1,F(61,84)),(0,-1,-F(23,120)),(6,1,F(25,8)),(-1,0,-F(115,264)),(-1,7,F(5,4)),(-6,-5,-F(23,6))],
        'U0105':[(-1,-1,-F(1,2)),(1,1,F(149,264)),(35,23,F(71,4)),(-2,-1,-F(115,132)),(2,13,F(35,12)),(-70,-59,-F(69,2))],
        'U0106':[(1,1,F(29,52)),(41,27,F(167,8)),(-2,-1,-F(23,26)),(2,15,F(37,12)),(-82,-69,-F(161,4))],
        'U0107':[(2,1,F(149,132)),(70,11,F(71,2)),(-1,0,-F(115,264)),(1,0,F(1,2)),(1,6,F(35,24)),(-35,-12,-F(69,4))],
    }
    out['V0']=[(0,-1,-F(1,12)),(-1,2,F(1,12)),(1,0,F(v0cap))]
    out['V11']=[(4,3,F(25,12)),(-4,-3,-F(23,12)),(0,-1,-F(1,12)),(0,1,F(1,6))]
    out['V12']=[(4,1,F(25,12)),(-4,-1,-F(23,12)),(0,-1,-F(1,12)),(0,1,F(1,6))]
    cases=[
        (F(139,336),F(605,1456),F(1),F(29,52)),
        (F(605,1456),F(17,40),F(41,27),F(167,216)),
        (F(17,40),F(51,112),F(4,3),F(25,36)),
        (F(461,1008),F(11,24),F(1,6),F(35,144)),
        (F(11,24),F(23,48),F(0),F(1,6)),
        (F(23,48),F(265,552),F(4),F(25,12)),
        (F(265,552),F(4727,9816),F(1,6),F(35,144)),
        (F(4727,9816),F(151,312),F(70,11),F(71,22)),
        (F(151,312),F(503,1008),F(4),F(25,12)),
    ]
    for i,(lo,hi,p,q) in enumerate(cases,1):
        out[f'U02{i:02d}']=[(-1,-1,-lo),(1,1,hi),(0,-1,-a),(p,p+1,q)]
    out['U0210']=[(1,2,F(187,348)),(0,-1,a-F(23,132)),(29,-1,F(59,6)+a),(-1,0,-F(23,76)),(-1,58,6),(-12,-11,11*a-F(23,4))]
    out['U0211']=[(1,2,F(149,264)),(0,-1,a-F(23,120)),(6,2,F(13,6)),(-1,0,-F(23,84)),(1,0,F(1,3)),(1,16,F(9,4)),(-6,-5,5*a-F(23,8))]
    out['U0212']=[(1,2,F(53,76)),(0,-1,a-F(23,132)),(12,2,F(25,4)),(-1,0,-F(161,348)),(-1,38,F(11,3)),(-12,-11,11*a-F(23,3))]
    out['U0301']=[(2,3,F(25,24)),(4,10,F(73,24)),(0,4,F(9,8))]
    out['U0302']=[(1,4,F(25,24)),(6,8,F(73,24)),(4,0,F(9,8))]
    out['U0303']=[(-1,-1,-F(115,264)),(1,1,F(11,24)),(1,7,F(35,24))]
    out['U0401']=[(2,1,F(25,24)),(4,6,F(73,24)),(0,4,F(9,8)),(-2,-1,-F(23,24)),(-4,2,-F(23,24))]
    out['U0402']=[(1,0,F(25,48)),(0,1,F(23,192)),(-1,-1,-F(23,48))]
    return {name:[tuple(map(F,L)) for L in lines] for name,lines in out.items()}

class Buchstab:
    def __init__(self,rtol=2e-13,atol=2e-15):
        self.sol={}
        self.rtol=rtol;self.atol=atol
        for n in range(2,12):
            fn=lambda u,y:np.asarray([self.omega(u-1)])
            result=solve_ivp(fn,(n,n+1),[self.capital(n)],method='DOP853',rtol=rtol,atol=atol,dense_output=True,max_step=.05)
            assert result.success
            self.sol[n]=result.sol
    def capital(self,u):
        if u<=2:return 1.
        n=min(int(math.floor(u)),max(self.sol,default=1))
        if n not in self.sol:n-=1
        return float(self.sol[n](u)[0])
    def omega(self,u):return self.capital(u)/u
    def primitive(self,u):return self.capital(u+1)-1.

def slice_region(x,lines):
    lo=float(a);hi=min(x,(1-x)/2)
    for A,B,C in lines:
        A,B,C=float(A),float(B),float(C)
        if B>0:hi=min(hi,(C-A*x)/B)
        elif B<0:lo=max(lo,(C-A*x)/B)
        elif A*x>C+1e-15:return None
    return (lo,hi) if hi>lo+1e-15 else None

def bad_intervals(x,regs):
    intervals=[]
    for lines in regs.values():
        v=slice_region(x,lines)
        if v:intervals.append(v)
    intervals.sort();bad=[];left=float(a);top=min(x,(1-x)/2)
    for lo,hi in intervals:
        if lo>left:bad.append((left,lo))
        left=max(left,hi)
    if left<top:bad.append((left,top))
    return bad

def breakpoints(regs):
    points={a,F(1,3),F(1,2)};edges=[];polys={}
    for name,lines in regs.items():
        pp=polygon(lines);polys[name]=pp
        for p,q in zip(pp,pp[1:]+pp[:1]):
            points.add(p[0]);points.add(q[0])
            if p[0]==q[0]:continue
            slope=(q[1]-p[1])/(q[0]-p[0]);intercept=p[1]-slope*p[0]
            edges.append((slope,intercept,min(p[0],q[0]),max(p[0],q[0])))
    # Include domain and Buchstab piece boundaries as extra edges.
    edges.extend([(F(0),a,a,F(1,2)),(F(1),F(0),a,F(1,3)),(-F(1,2),F(1,2),F(1,3),F(1,2))])
    for j in range(2,12):edges.append((-F(1,j+1),F(1,j+1),a,F(1,2)))
    for (s,b,l,r),(ss,bb,ll,rr) in combinations(edges,2):
        if s==ss:continue
        x=(bb-b)/(s-ss)
        if max(l,ll)<=x<=min(r,rr) and a<=x<=F(1,2):points.add(x)
    return sorted(points),polys

def integrate(v0cap,rtol=2e-13,eps=2e-10):
    regs=regions(v0cap);pts,polys=breakpoints(regs);buch=Buchstab(rtol)
    def fun(x):
        ans=0.
        for lo,hi in bad_intervals(x,regs):
            ulo=(1-x)/lo-1;uhi=(1-x)/hi-1
            ans+=(buch.primitive(ulo)-buch.primitive(uhi))/(x*(1-x))
        return ans
    total=0.;err=0.
    for l,r in zip(pts[:-1],pts[1:]):
        ans,estimate=quad(fun,float(l),float(r),epsabs=eps/(len(pts)+1),epsrel=eps,limit=50)
        total+=ans;err+=estimate
    result={'V0_cap':str(F(v0cap)),'I0':total,'quad_error_estimate':err,'buchstab_rtol':rtol,'outer_epsilon':eps,'pieces':len(pts)-1,'nonempty_regions':[name for name,pp in polys.items() if len(pp)>=3],'V0_vertices':[[str(x),str(y)] for x,y in polys['V0']]}
    print(json.dumps(result),flush=True)
    return result

if __name__=='__main__':
    results=[]
    for cap in [F(1,16),F(1,6)]:results.append(integrate(cap))
    with Path(__file__).with_name('li_III_i0_convention_results.json').open('w') as f:json.dump(results,f,indent=2)
\end{filecontents*}
\begin{filecontents*}[overwrite]{li_III_i0_convention_results.json}
[
  {
    "V0_cap": "1/16",
    "I0": 0.966495918091412,
    "quad_error_estimate": 1.0730260214713147e-14,
    "buchstab_rtol": 2e-13,
    "outer_epsilon": 2e-10,
    "pieces": 189,
    "nonempty_regions": [
      "U0101",
      "U0102",
      "U0103",
      "U0104",
      "U0105",
      "U0106",
      "U0107",
      "V11",
      "V12",
      "U0201",
      "U0202",
      "U0203",
      "U0204",
      "U0205",
      "U0206",
      "U0207",
      "U0208",
      "U0209",
      "U0210",
      "U0211",
      "U0212",
      "U0301",
      "U0302",
      "U0303",
      "U0401",
      "U0402"
    ],
    "V0_vertices": []
  },
  {
    "V0_cap": "1/6",
    "I0": 0.966495918091412,
    "quad_error_estimate": 1.0730260214713147e-14,
    "buchstab_rtol": 2e-13,
    "outer_epsilon": 2e-10,
    "pieces": 194,
    "nonempty_regions": [
      "U0101",
      "U0102",
      "U0103",
      "U0104",
      "U0105",
      "U0106",
      "U0107",
      "V0",
      "V11",
      "V12",
      "U0201",
      "U0202",
      "U0203",
      "U0204",
      "U0205",
      "U0206",
      "U0207",
      "U0208",
      "U0209",
      "U0210",
      "U0211",
      "U0212",
      "U0301",
      "U0302",
      "U0303",
      "U0401",
      "U0402"
    ],
    "V0_vertices": [
      [
        "1/6",
        "1/8"
      ],
      [
        "23/252",
        "11/126"
      ],
      [
        "1/6",
        "11/126"
      ]
    ]
  }
]
\end{filecontents*}
% END EMBEDDED VERIFICATION MATERIALS

\begin{filecontents*}[overwrite]{mills_chain_5_4.gp}
\\ Greedy chain for c = 5/4.  J(x) = { m : m^4 >= x^5 and (m+1)^4 < (x+1)^5 }.
\\ Each term is proven prime by isprime() (PARI/GP primality proof).
default(parisize, 10^8);
{
p = 1009; L = List([p]);
for (n = 1, 21,
  lo = sqrtnint(p^5, 4); if (lo^4 < p^5, lo++);
  m = nextprime(lo);
  if (!((m+1)^4 < (p+1)^5), error("stuck at n=", n));
  if (!isprime(m), error("not proven prime"));
  listput(L, m); p = m);
}
{
for (i = 1, #L,
  q = L[i];
  if (i == 1, off = 0, a = L[i-1]; r = sqrtnint(a^5, 4); if (r^4 < a^5, r++); off = q - r);
  printf("%d  digits=%d  offset=%d\n", i, #digits(q), off));
}
{
ok = 1;
for (i = 1, #L-1, a = L[i]; b = L[i+1]; if (!(b^4 >= a^5 && (b+1)^4 < (a+1)^5), ok = 0));
print("nesting ok: ", ok);
}
write("chain54.txt", Vec(L));
\p 400
{
N = #L; A = exp(log(L[N] + 1/2) * (4/5)^N);
print("A approx = ", A);
bad = 0; for (i = 1, N, if (floor(A^((5/4)^i)) != L[i], bad++));
print("floor mismatches: ", bad);
}
quit;
\end{filecontents*}
\documentclass[11pt,a4paper]{article}
\usepackage[T1]{fontenc}
\usepackage{lmodern}
\usepackage[margin=25mm]{geometry}
\usepackage{amsmath,amssymb,amsthm,mathtools}
\allowdisplaybreaks[1]
\usepackage{booktabs,tabularx,array,enumitem,microtype,xcolor,listings}
\usepackage[colorlinks=true,linkcolor=blue!45!black,citecolor=blue!45!black,urlcolor=blue!45!black]{hyperref}
\usepackage[capitalise,nameinlink]{cleveref}
\usepackage{fancyhdr}
\pagestyle{fancy}\fancyhf{}
\fancyhead[L]{\small Harmonic branching and quantitative sieve bounds}
\fancyhead[R]{\small KEIO}\fancyfoot[C]{\thepage}
\setlength{\headheight}{14pt}
\newtheorem{theorem}{Theorem}[section]
\newtheorem{proposition}[theorem]{Proposition}
\newtheorem{lemma}[theorem]{Lemma}
\newtheorem{corollary}[theorem]{Corollary}
\theoremstyle{remark}\newtheorem{remark}[theorem]{Remark}
\newcommand{\eps}{\varepsilon}
\newcommand{\PP}{\mathcal P}
\newcommand{\Q}{\mathcal Q}
\newcommand{\W}{\mathcal W}
\newcommand{\meas}{\operatorname{meas}}
\newcommand{\R}{\mathbb R}
\newcommand{\N}{\mathbb N}

\theoremstyle{definition}
\newtheorem{hypothesis}[theorem]{Hypothesis}
\newtheorem{definition}[theorem]{Definition}
\newcommand{\Z}{\mathbb Z}
\newcommand{\Ch}{\operatorname{Ch}}
\newcommand{\LL}{\mathcal L}
\newcommand{\fl}[1]{\lfloor #1\rfloor}
\setlength{\emergencystretch}{1.5em}
\hypersetup{pdftitle={Harmonic Branching and Quantitative Sieve Bounds for Mills-Type Prime-Representing Functions},pdfauthor={KEIO}}
\lstset{basicstyle=\ttfamily\footnotesize,breaklines=true,columns=fullflexible,
keepspaces=true,showstringspaces=false,frame=single,rulecolor=\color{black!20}}
\title{\bfseries Harmonic Branching and Quantitative Sieve Bounds\\
for Mills-Type Prime-Representing Functions\\[7pt]
\large\normalfont An explicit conditional exponent below $24/23$}
\author{KEIO}
\date{Revised and consolidated edition\\9 November 2026}
\begin{document}
\maketitle
\begin{abstract}
For $c>1$, let $\W(c)$ be the set of real numbers $A>1$ for which
$\lfloor A^{c^n}\rfloor$ is prime for every $n\ge1$.
We prove a harmonic mass--compactness principle that converts a quantitative
almost-all prime lower bound into uncountably many representing constants.
An exceptional measure of size $X\exp[-D(\log\log X)^2]$, with sufficiently
large fixed $D$, controls all depths of the prime tree simultaneously.
The argument also treats bounded variable exponent sequences and proves
unboundedness, Lebesgue nullity, and nowhere density.
It recovers the zero-density ranges $c>6/5$ and $c>15/13$, and quantifies
the finite sieve inputs leading to $c>20/19$ and $c>43/41$ with their
respective source dependencies.
For the sharper extension, we identify an internal integer-moment
obstruction at the short-interval exponent $1/24$ and remove it by a
two-factor estimate valid at $1/25$, using the Chen--Gupta--Li large-value
theorem. Exact rational certificates verify the additional trilinear
regions and the finite boundary geometry.
Assuming four explicitly stated numerical Buchstab inequalities from
Li's third short-interval manuscript, we obtain
$\theta_*=1/24-10^{-23}$ and hence all
$c>24/(23+24\cdot10^{-23})$.
In particular, $c=24/23-10^{-23}$ is admissible.
The four numerical inequalities remain external hypotheses; the
exploratory recomputation of one integral is not an interval certificate.
The full prime-interval exponent $1/25$ is not claimed.
The single source embeds the exact verification programs and certificates.
\end{abstract}

\noindent\textbf{Keywords:} prime-representing functions; Mills exponent;
harmonic mass; primes in short intervals; Buchstab decomposition;
Dirichlet polynomials; exact rational certificates.\\
\noindent\textbf{MSC 2020:} 11N05, 11N35, 11M26.

\newpage
\tableofcontents
\newpage

\section{Introduction and principal results}\label{sec:introduction}

Mills' prime-representing construction asks for a real number whose
successive prescribed powers have prime integer parts~\cite{Mil47}.
Write $\PP$ for the primes and define
\[
 \W(c)=\{A>1:\lfloor A^{c^n}\rfloor\in\PP\text{ for every }n\ge1\}.
\]
Here $c$ is the exponent parameter, and $A$ is the representing constant.
An improvement in $c$ does not compute or alter the least usual Mills
constant associated with exponent $3$.
For a sequence $\boldsymbol c=(c_n)_{n\ge1}$, put
$C_n=c_1\cdots c_n$ and
\[
 \W(\boldsymbol c)=
 \{A>1:\lfloor A^{C_n}\rfloor\in\PP\text{ for every }n\ge1\}.
\]

The nested prime cells are
$J_c(m)=[m^c,(m+1)^c-1)$.
At height $X\asymp m^c$, their lengths are comparable to
$X^{1-1/c}$. A prime lower bound in almost all intervals of length
$X^\theta$ is therefore relevant when $\theta<1-1/c$.
The remaining difficulty is to keep prime-rich cells at every depth.
The harmonic construction measures the entire finite level by
$\sum_p1/p$. It compares its accumulated mass with the mass of all
defective integers at that scale, and then uses finite branching
compactness. The full argument, including the uniformity in depth,
appears in \cref{hb:sec:transfer}.

\begin{theorem}[explicit conditional Mills range]\label{thm:main}
Assume the four numerical inequalities of \cref{hyp:liIII}, with the
literal printed Li III regions and the conventions specified there.
Let
\[
 \Delta=10^{-23},\qquad
 \theta_* =\frac1{24}-\Delta,\qquad
 c_* =\frac{1}{1-\theta_*}=\frac{24}{23+24\Delta}.
\]
For every fixed $D>0$, there is $\kappa>0$ such that, for sufficiently
large $X$,
\begin{equation}\label{eq:main:prime}
 \meas\left\{x\in[X,2X]:
 \pi(x)-\pi(x-X^{\theta_*})
 <\kappa\frac{X^{\theta_*}}{\log X}\right\}
 \ll X\exp[-D(\log\log X)^2].
\end{equation}
For every sequence satisfying
$c_*<\inf_n c_n\le\sup_n c_n<\infty$,
the set $\W(\boldsymbol c)$ is uncountable, unbounded,
Lebesgue-null, and nowhere dense. In particular, this holds for every
constant $c>c_*$, including the exact rational value
\[
                   \boxed{c_{\mathrm{new}}=\frac{24}{23}-10^{-23}}.
\]
\end{theorem}

The lower boundary $c_*$ is strict. The displayed $c_{\mathrm{new}}$
is a convenient admissible point, since
\begin{equation}\label{eq:intro:exactgap}
 c_{\mathrm{new}}-c_*
 =\frac{\Delta(47-552\Delta)}{23(23+24\Delta)}>0.
\end{equation}
The endpoint $c=24/23$ is also included. No optimality is asserted for
$\Delta$, and the theorem does not include $c=c_*$.

\begin{theorem}[zero-density specializations]\label{thm:baseline}
The conclusions about $\W(\boldsymbol c)$ hold unconditionally when
$6/5<\inf_n c_n\le\sup_n c_n<\infty$, using the classical
Ingham--Huxley zero-density estimates. With the Guth--Maynard estimate,
the lower bound $6/5$ may be replaced by $15/13$.
Under the density hypothesis
$N(\sigma,T)\ll_\eps T^{2(1-\sigma)+\eps}$, it may be replaced by $1$.
\end{theorem}

\Cref{sec:density} proves these specializations without the four
numerical hypotheses used in \cref{thm:main}.
The intermediate sieve ranges are retained in \cref{sec:earlier-sieves}.
Their numerical inputs and those of the new theorem are separate.

\begin{table}[ht]
\centering\small
\caption{Inputs and the ranges established by the corresponding arguments.
The inequalities in the Mills column are strict.}
\label{tab:ranges}
\begin{tabularx}{\linewidth}{@{}>{\raggedright\arraybackslash}p{.23\linewidth}>{\raggedright\arraybackslash}p{.18\linewidth}>{\raggedright\arraybackslash}p{.17\linewidth}>{\raggedright\arraybackslash}X@{}}
\toprule
Input & Prime-interval scale & Mills range & Dependency\\
\midrule
Ingham--Huxley & $\theta>1/6$ & $c>6/5$ & Published density estimates\\
Guth--Maynard & $\theta>2/15$ & $c>15/13$ & Published density estimate\\
Jia finite sieve & $\theta>1/20$ & $c>20/19$ & Quantified source proof and its published numerical margin\\
Li I finite sieve & $\theta>2/43$ & $c>43/41$ & Its three stated numerical region bounds\\
Present extension & $\theta_*=1/24-10^{-23}$ & $c>c_*$ & Four Li III bounds; CGL large values; exact certificates\\
Density hypothesis & Every $\theta>0$ & Every $c>1$ & Conditional zero-density assertion\\
\bottomrule
\end{tabularx}
\end{table}

\paragraph{Proof dependencies.}
The harmonic transfer is proved independently of the sieve.
The analytic extension first repairs the two-factor obstruction, then
verifies the trilinear and Type I regions, and finally stabilizes a
finite signed Buchstab decomposition. The exact comparator-loss bound
is $0.00295856$, against an assumed original margin greater than $0.02$.
The positive margin supplies \eqref{eq:main:prime}; harmonic transfer
then proves the Mills assertion. The two-factor lemma itself works at
$1/25$, but no positive final sieve margin at $1/25$ is established.

The Chen--Gupta--Li result~\cite{CGL} and Li III~\cite{LiIII} are used at
the stated preprint/manuscript versions. The numerical hypotheses are
not replaced by the exact finite-geometry checks. This consolidated
paper contains the earlier harmonic and quantitative arguments, so
none of its proofs requires an unpublished predecessor by the author.
It makes no claim of priority over all other work.


% Integration note: the main preamble needs
% \newtheorem{hypothesis}[theorem]{Hypothesis}
% The seven raw regions \mathcal C_j(theta) are the locally reproduced
% regions in Appendix \ref{app:regions}; do not intersect them with Omega
% before using them in the definition of T_j below.

\subsection{The precise numerical sieve hypothesis}
\label{sec:numerical-input}

The numerical input is stated here as four definite inequalities.
Its regions are the literal regions in Li~\cite[Section~7,
equations (15), (23), (24), and (26)]{LiIII}; in particular the
printed empty region $V_0$ is not silently replaced by a larger one.
Boundary conventions on finitely many hyperplanes do not affect the
integrals below.

Put $a=11/126$, $\theta_0=1/24$, $s=t_1+t_2$, and
\[
 \Omega=\{(t_1,t_2):a\le t_1<1/2,\quad
             a\le t_2<\min(t_1,(1-t_1)/2)\}.
\]
Using the seven raw regions $\mathcal C_j(\theta_0)$ recorded in
\cref{app:regions}, define
\[
 U_{01}=\Omega\cap\bigcup_{j=1}^7\mathcal C_j(\theta_0).
\]
The three printed additional components are
\begin{align*}
 V_0&=\{(t_1,t_2):t_2>1/12,\quad
                       2t_2-1/12<t_1<1/16\},\\
 V_{11}&=\{(t_1,t_2):|4t_1+3t_2-2|<1/12,\quad
                                  1/12<t_2<1/6\},\\
 V_{12}&=\{(t_1,t_2):|4t_1+t_2-2|<1/12,\quad
                                  1/12<t_2<1/6\},\\
 V&=\Omega\cap(V_0\cup V_{11}\cup V_{12}).
\end{align*}
Here $\Omega\cap V_0=\varnothing$, since $a>1/16$.
The final displayed component is the second occurrence of the label
$V_{11}$ in the source and is called $V_{12}$ here, consistently with
the source's preceding union. This convention is fixed throughout
the perturbation.

For $1\le i\le9$, let
\[
 B_i=\{(t_1,t_2)\in\Omega:
            \ell_i\le s\le u_i,\quad t_2\le q_i s+r_i\},
\]
where the exact endpoints are
\begin{center}
\small
\begin{tabular}{@{}ccccc@{}}
\toprule
$i$ & $\ell_i$ & $u_i$ & $q_i$ & $r_i$\\
\midrule
1 & $139/336$ & $605/1456$ & $-1$ & $29/52$\\
2 & $605/1456$ & $17/40$ & $-41/27$ & $167/216$\\
3 & $17/40$ & $51/112$ & $-4/3$ & $25/36$\\
4 & $461/1008$ & $11/24$ & $-1/6$ & $35/144$\\
5 & $11/24$ & $23/48$ & $0$ & $1/6$\\
6 & $23/48$ & $265/552$ & $-4$ & $25/12$\\
7 & $265/552$ & $4727/9816$ & $-1/6$ & $35/144$\\
8 & $4727/9816$ & $151/312$ & $-70/11$ & $71/22$\\
9 & $151/312$ & $503/1008$ & $-4$ & $25/12$\\
\bottomrule
\end{tabular}
\end{center}
For $j=1,2,3$, put
\[
 T_j=\{(t_1,t_2)\in\Omega:
       (t_1,t_2+u)\in\mathcal C_j(\theta_0)
                     \text{ for every }u\in[a,t_2]\}.
\]
Because all defining inequalities are affine in $u$, testing its two
endpoints gives exactly the source components $U_{0210}$,
$U_{0211}$, and $U_{0212}$, respectively. Define
\[
 \widetilde U_{02}=\bigcup_{i=1}^9B_i\ \cup\ T_1\cup T_2\cup T_3.
\]

For the remaining components use the two fixed Type I regions
\begin{align*}
 W_1&=\{(m,h):2m+h\le25/24,\quad
                      4m+6h\le73/24,\quad4h\le9/8\},\\
 W_2&=\{(m,h):m\le25/48,\quad h\le23/192\}.
\end{align*}
Set
\begin{align*}
 \widetilde U_{03}=\Omega\cap\bigl(&
  \{(t_1,t_2):(s,t_2)\in W_1\}
  \cup\{(t_1,t_2):(2t_2,t_1)\in W_1\}\\
  &\cup\{(t_1,t_2):115/264\le s\le11/24,
                             \quad t_2\le35/144-s/6\}\bigr),\\
 \widetilde U_{04}=\Omega\cap\bigcup_{j=1}^2
  &\{(t_1,t_2):(t_1,t_2)\in W_j,
                            \ (1-s,t_2)\in W_j\}.
\end{align*}
The disjoint priority domains used in the signed decomposition are
\begin{align*}
 U_{02}&=\widetilde U_{02}\setminus(U_{01}\cup V),\\
 U_{03}&=\widetilde U_{03}\setminus(U_{01}\cup V\cup U_{02}),\\
 U_{04}&=\widetilde U_{04}\setminus(U_{01}\cup V\cup U_{02}\cup U_{03}),\\
 U_0&=U_{01}\cup V\cup U_{02}\cup U_{03}\cup U_{04}.
\end{align*}

For a supplied tuple $\boldsymbol t=(t_1,\ldots,t_n)$, write
$\mathcal{G}_n(\boldsymbol t)$ if there is an ordered partition
$\{1,\ldots,n\}=J\sqcup K$ into two nonempty blocks such that,
on putting $m=\sum_{j\in J}t_j$ and $h=\sum_{k\in K}t_k$, either
\[
 (m,h)\in V,
 \qquad\text{or}\qquad
 (m,h)\in U_{01}\quad\text{and}\quad
                       \min(|J|,|K|)=1.
\]
This fixes the literal two-block ``good'' convention immediately
preceding equation (12) in Li~\cite{LiIII}. The larger assignment
inventory in the verification appendix permits unused factors only
to overcount possible boundary faces; it does not enlarge this
numerical hypothesis's good sets.

Write $\beta=1-t_1-t_2-t_3$. The four integration domains are
\begin{align*}
 R_0={}&\Omega\setminus U_0,\\
 R_1={}&\bigl\{\boldsymbol t\in\mathbb R^4:
  (t_1,t_2)\in U_{03},\quad
  a\le t_3<\min(t_2,(1-s)/2),\quad \neg \mathcal{G}_3(t_1,t_2,t_3),\\
 &\hspace{28mm}a\le t_4<\min(t_3,(1-s-t_3)/2),\quad
                          \neg \mathcal{G}_4(t_1,t_2,t_3,t_4)\bigr\},\\
 R_2={}&\bigl\{\boldsymbol t\in\mathbb R^4:
  (t_1,t_2)\in U_{04},\quad
  a\le t_3<\min(t_2,(1-s)/2),\quad \neg \mathcal{G}_3(t_1,t_2,t_3),\\
 &\hspace{28mm}a\le t_4<t_1/2,\quad
                                \neg \mathcal{G}_4(\beta,t_2,t_3,t_4)\bigr\},\\
 R_3={}&\bigl\{\boldsymbol t\in\mathbb R^3:
  (t_1,t_2)\in R_0,\quad
  t_2<t_3<(1-s)/2,\quad \mathcal{G}_3(t_1,t_2,t_3)\bigr\}.
\end{align*}

Let $\omega$ be Buchstab's function, defined by
$\omega(u)=1/u$ for $1\le u\le2$ and
$(u\omega(u))'=\omega(u-1)$ for $u>2$. Use the source's upper
envelope
\[
 \omega_+(u)=
 \begin{cases}
  \omega(u),&1\le u<4,\\
  5617/10000,&u\ge4,
 \end{cases}
\]
which is the function denoted $\omega_1$ in Li~\cite[Section~7]{LiIII}.
Its upper-envelope property is the standard Buchstab bound used there;
see also Harman~\cite{Harman}.
All arguments of these functions below are at least one.
With $\mathrm d\boldsymbol t_n=\mathrm dt_n\cdots\mathrm dt_1$, put
\begin{align}
 I_0&=\int_{R_0}
       \frac{\omega((1-t_1-t_2)/t_2)}{t_1t_2^2}
                    \,\mathrm d\boldsymbol t_2,\label{eq:input-I0}\\
 I_1&=\int_{R_1}
       \frac{\omega_+((1-t_1-t_2-t_3-t_4)/t_4)}
                         {t_1t_2t_3t_4^2}
                    \,\mathrm d\boldsymbol t_4,\label{eq:input-I1}\\
 I_2&=\int_{R_2}
       \frac{\omega_+((t_1-t_4)/t_4)
             \omega_+((1-t_1-t_2-t_3)/t_3)}{t_2t_3^2t_4^2}
                    \,\mathrm d\boldsymbol t_4,\label{eq:input-I2}\\
 I_3&=\int_{R_3}
       \frac{\omega((1-t_1-t_2-t_3)/t_3)}{t_1t_2t_3^2}
                    \,\mathrm d\boldsymbol t_3.\label{eq:input-I3}
\end{align}

\begin{hypothesis}[the four Li III numerical inequalities]
\label{hyp:liIII}
For exactly the preceding domains and good predicate,
\[
             I_1<\frac1{50},\qquad I_2<\frac3{100},\qquad
             I_0<\frac{97}{100},\qquad I_3>\frac1{25}.
\]
\end{hypothesis}

These are the four numerical bounds recorded in Li's equations
(15), (23), (24), and (26), now isolated from the analytic assertions.
They imply $1-I_1-I_2-I_0+I_3>1/50$.
The present paper gives exact certificates for the finite algebraic
and geometric steps. It also reports an exploratory numerical
recomputation of $I_0$, but no rigorous enclosure of that integral.
$I_1$, $I_2$, and $I_3$ have not been independently recomputed here.
Thus none of the four inequalities in \cref{hyp:liIII} is replaced
by an independent quadrature certificate in this paper.

% Integration draft: standalone sections, without a preamble.
% All new labels have the prefix hb:.  The surrounding manuscript supplies
% the standard theorem environments and the macros listed by the editor.
% The reference sec:density is to the separately integrated density section.

\section{Harmonic branching and compactness}\label{hb:sec:mass}

The analytic estimates will control the total weight of defective
integers at a given scale.  They need not identify a good continuation
below each individual prime.  The following two results isolate the
global argument that makes such information sufficient.

\begin{definition}[Branching system]\label{hb:def:branching}
A branching system consists of a finite nonempty root set
$V_1\subset\Z_{\ge2}$, good sets $G_n\subset\Z_{\ge2}$, and finite
child sets $\Ch_n(x)\subset\Z_{\ge2}$.  At each fixed level, child
sets belonging to distinct parents are disjoint.  Put
\[
 V_{n+1}=\bigcup_{x\in V_n\cap G_n}\Ch_n(x),
 \qquad w(S)=\sum_{x\in S}\frac1x
\]
for finite sets $S$.  An infinite branch is a sequence $(x_n)$ with
$x_n\in V_n$ and $x_{n+1}\in\Ch_n(x_n)$ for every $n$.
\end{definition}

\begin{theorem}[Mass--compactness principle]\label{hb:thm:mass}
Suppose $\lambda_n>0$ and $\beta_n\ge0$ satisfy
\[
 w(\Ch_n(x))\ge\frac{\lambda_n}{x}
       \quad(x\in V_n\cap G_n),
 \qquad w(V_n\setminus G_n)\le\beta_n.
\]
Define
\[
 M_n=w(V_1)\prod_{k<n}\frac{\lambda_k}{2}.
\]
If $\beta_n\le M_n/2$ for every $n$, then $w(V_n)\ge M_n>0$
for every $n$, and the system has an infinite branch.
Every node on such a branch is good at its own level.
\end{theorem}

\begin{proof}
The assertion $w(V_1)\ge M_1$ is an equality.  If
$w(V_n)\ge M_n$, then
\[
 w(V_n\cap G_n)\ge w(V_n)-\beta_n
       \ge w(V_n)-M_n/2\ge w(V_n)/2.
\]
Disjointness of the child sets therefore gives
\[
 w(V_{n+1})\ge\lambda_n w(V_n\cap G_n)
       \ge\frac{\lambda_n}{2}w(V_n)\ge M_{n+1}.
\]
Thus every level is nonempty.  Each level is finite, each nonroot
node has a unique parent, and the root set is finite.  After adjoining
a common root, K\"onig's lemma gives an infinite path.  A node can have
a child in this forest only if it belongs to its good set.
\end{proof}

\begin{corollary}[Binary subdivision fixed by finite prefixes]
\label{hb:cor:binary}
Under the two local and defect bounds in \cref{hb:thm:mass}, put
\[
 s_n=\frac{w(V_1)}3\prod_{k<n}\frac{\lambda_k}{6}.
\]
Suppose, for every $n$,
\[
 \beta_n\le s_n/2,\qquad
 \max_{x\in V_n}\frac1x\le s_n,
 \tag{B}\label{hb:eq:binaryconditions}
\]
where the second condition is vacuous for an empty set.
Then the system has at least $2^{\aleph_0}$ distinct infinite branches.
\end{corollary}

\begin{proof}
A finite collection of weight at least $3s$, with each atom of weight
at most $s$, splits into two disjoint parts of weight at least $s$.
Indeed, in increasing integer order, stop at the first partial sum
at least $s$.  Its weight is at most $2s$, so the remaining weight is
at least $s$.

Fix these splits once for each finite binary word; they are not
chosen separately for different infinite addresses.
For a word $\sigma$ of length $n-1$, construct
$U_\sigma\subset V_n$ with $w(U_\sigma)\ge3s_n$, starting with
$U_\varnothing=V_1$.  Split $U_\sigma$ into
$S_{\sigma,0},S_{\sigma,1}$, each of weight at least $s_n$.
For $b\in\{0,1\}$ define
\[
 U_{\sigma b}
   =\bigcup_{x\in S_{\sigma,b}\cap G_n}\Ch_n(x).
\]
The defective weight in any selected subset is bounded by the
defective weight of the entire level.  Hence
\[
 w(U_{\sigma b})
 \ge\lambda_n\bigl(w(S_{\sigma,b})-\beta_n\bigr)
 \ge\lambda_ns_n/2=3s_{n+1}.
\]
The atom bound in \eqref{hb:eq:binaryconditions} permits the next
split.  Induction defines all finite-prefix families.

For an infinite address $\omega=(\omega_1,\omega_2,\ldots)$, the
selected level-$n$ set is
$S_{\omega|_{n-1},\omega_n}$.  These sets form a locally finite
forest with finitely many roots and a nonempty set at every depth.
K\"onig's lemma supplies an infinite branch in that forest.
If two addresses first differ at position $n$, their level-$n$
nodes belong to the two disjoint parts of the same $U_\sigma$.
Their branches are therefore distinct.
For definiteness, a branch for each address can be selected by taking,
at each step, the least node that has descendants at arbitrarily
large depths.  This also makes the simultaneous choice explicit.
\end{proof}

\begin{remark}\label{hb:rem:global}
Neither result assumes that each good parent has a good child.
The defect bound is global, and remains valid after any of the
finite-prefix selections.  Concentrating defective nodes beneath
one parent does not invalidate the comparison.
\end{remark}

\section{From almost-all short intervals to full prime cells}
\label{hb:sec:cells}

We formulate the prime input using centered intervals solely to
accommodate both orientations of the analytic estimates.
A lower bound in either $(u-X^\theta,u]$ or
$(u,u+X^\theta]$ implies the centered lower bound below by inclusion.
The same holds for a multiplicative window of length at most
$X^\theta$, provided its lower bound is a fixed positive multiple
of $X^\theta/\log X$.

\begin{hypothesis}[$\mathsf{AA}(\theta)$]\label{hb:hyp:AA}
Fix $0<\theta<1$.  For every fixed $D>0$, there are
$\kappa_D>0$ and $X_D$ such that, for all real $X\ge X_D$,
\[
 \meas E_{\theta,D}(X)
       \le X\exp\!\bigl[-D(\log\log X)^2\bigr],
 \tag{AA}\label{hb:eq:AA}
\]
where
\[
 E_{\theta,D}(X)=
 \left\{u\in[X,2X]:
  \pi(u+X^\theta)-\pi(u-X^\theta)
       <\kappa_D\frac{X^\theta}{\log X}\right\}.
\]
\end{hypothesis}

The lower-bound constant is allowed to depend on the fixed accuracy
$D$.  This causes no difficulty: the construction below chooses
$D$ once, before choosing the starting scale.
An implicit fixed multiplicative constant in \eqref{hb:eq:AA} is
harmless, since it is absorbed by first using a larger fixed accuracy.

\begin{lemma}[A poor cell uses a positive proportion of its length]
\label{hb:lem:average}
Let $h=X^\theta$ and let $E\subset[X,2X]$ be an exceptional set
outside which
\[
 \pi(u+h)-\pi(u-h)\ge\kappa h/\log X.
\]
If $I=[a,b)\subset[X,2X]$, $L=b-a\ge4h$, and
\[
 \#(\PP\cap I)<\frac{\kappa L}{16\log X},
\]
then $\meas(E\cap I)\ge L/4$.
\end{lemma}

\begin{proof}
For $u\in B_I=[a+h,b-h)$, the centered window is contained in $I$,
apart from endpoints of measure zero.  Each prime of $I$ is counted
for at most $2h$ units of $u$-measure.  Thus
\[
 \int_{B_I}\bigl(\pi(u+h)-\pi(u-h)\bigr)\,du
       \le2h\,\#(\PP\cap I)
       <\frac{\kappa hL}{8\log X}.
\]
The nonexceptional part of $B_I$ has measure at most $L/8$.
Since $\meas B_I=L-2h\ge L/2$, the exceptional part has
measure at least $3L/8$, which is more than required.
\end{proof}

For $c>1$ and an integer $m\ge2$ define
\[
 J_c(m)=[m^c,(m+1)^c-1),\qquad
 \PP_c(m)=\PP\cap J_c(m).
 \tag{J}\label{hb:eq:J}
\]
For a fixed $c$, these cells are pairwise disjoint.  The mean value
theorem gives
\[
 c\,m^{c-1}-1\le |J_c(m)|
       \le c\,(m+1)^{c-1}.
 \tag{L}\label{hb:eq:length}
\]

\begin{proposition}[Uniform bad-integer estimate]\label{hb:prop:bad}
Assume $\mathsf{AA}(\theta)$, and fix
\[
 \frac1{1-\theta}<c_-\le c_+<\infty.
\]
For every fixed $D>0$, there are $\delta_D\in(0,1]$ and $Y_D$
such that, for every real $Y\ge Y_D$ and every $c\in[c_-,c_+]$,
\[
 \#\left\{m\in[Y,2Y)\cap\Z:
   \#\PP_c(m)<
       \delta_D\frac{|J_c(m)|}{\log((m+1)^c)}
       \right\}
 \le Y\exp\!\bigl[-D(\log\log Y)^2\bigr].
 \tag{N}\label{hb:eq:bad}
\]
\end{proposition}

\begin{proof}
Use \cref{hb:hyp:AA} with a fixed accuracy $A>D$.
All cells under consideration lie in
$[Y^c,(2Y+1)^c]$, which is covered by $O_{c_+}(1)$ successive
dyadic intervals $[X,2X)$ beginning at $X=Y^c$.
At most $O_{c_+}(1)$ cells cross boundaries of this cover, because
the cells are disjoint.  Set those cells aside.

For every remaining cell $I=J_c(m)\subset[X,2X]$,
\[
 |I|\asymp_{c_-,c_+}Y^{c-1},\qquad X\asymp_{c_+}Y^c,
 \qquad
 \frac{|I|}{X^\theta}\gg
 Y^{\,c_-(1-\theta)-1}\longrightarrow\infty.
\]
All comparisons are uniform in $c\in[c_-,c_+]$.
Choose $\delta_D=\min(1,\kappa_A/16)$.
Since $\log((m+1)^c)\ge\log X$, a cell counted on the left of
\eqref{hb:eq:bad} satisfies the poor-cell hypothesis of
\cref{hb:lem:average}.  It accounts for at least $|I|/4$
of exceptional endpoint measure in its own dyadic range.
Disjointness therefore bounds the number of such cells by
\[
 O_{c_-,c_+}\!\left(
 Y\exp\!\bigl[-A(\log\log(Y^c))^2\bigr]\right).
\]
Now $\log\log(Y^c)=\log\log Y+\log c$, with $c$ in a fixed
compact interval above $1$.  Because $A>D$, the fixed constants,
the finite cover, and the $O_{c_+}(1)$ discarded boundary cells
are absorbed in the right side of \eqref{hb:eq:bad} for large $Y$.
\end{proof}

\begin{remark}\label{hb:rem:fullcell}
The averaging step charges a poor cell its full length
$\asymp Y^{c-1}$.  Charging only one interval of length $X^\theta$
would give a weaker count and would not prove \eqref{hb:eq:bad}.
There is no conversion of a fixed-logarithm theorem with a moving
logarithmic exponent in this argument.
\end{remark}

\section{The harmonic transfer theorem}\label{hb:sec:transfer}

For a sequence $\mathbf c=(c_k)_{k\ge1}$ with $c_k>1$, put
\[
 C_n=\prod_{k=1}^n c_k,\qquad
 \W(\mathbf c)=
 \{A>1:\lfloor A^{C_n}\rfloor\in\PP\text{ for every }n\ge1\}.
\]
The constant-sequence case is denoted by $\W(c)$.

\begin{lemma}[The strict upper endpoint]\label{hb:lem:nesting}
Suppose $p_n\ge2$ and
$p_{n+1}\in J_{c_{n+1}}(p_n)$ for every $n\ge1$.
Then
\[
 A=\lim_{n\to\infty}p_n^{1/C_n}
\]
exists, is greater than $1$, and satisfies
$\lfloor A^{C_n}\rfloor=p_n$ for every $n$.
Distinct sequences $(p_n)$ give distinct values of $A$.
\end{lemma}

\begin{proof}
Put $a_n=p_n^{1/C_n}$ and $b_n=(p_n+1)^{1/C_n}$.
The child inequalities give
\[
 a_{n+1}\ge a_n,\qquad b_{n+1}<b_n.
\]
The first sequence is increasing and bounded above by $b_1$,
so its limit $A$ exists.  For each fixed $n$,
\[
 a_n\le A\le b_{n+1}<b_n.
 \tag{E}\label{hb:eq:endpoints}
\]
Consequently $p_n\le A^{C_n}<p_n+1$, including the required
strict upper floor inequality.  Also $A\ge p_1^{1/c_1}>1$.
An $A$ uniquely determines every integer $\lfloor A^{C_n}\rfloor$,
which proves distinctness.
\end{proof}

\begin{proposition}[Nullity and category need no prime input]
\label{hb:prop:null}
More generally, for any sequence of positive exponents $C_n\to\infty$,
the set
\[
 \{A>1:\lfloor A^{C_n}\rfloor\in\PP\text{ for every }n\}
\]
has Lebesgue measure zero and is nowhere dense.  Its closure also
has measure zero.
\end{proposition}

\begin{proof}
Fix $1<a<b$.  The set satisfying the $n$th prime condition in
$[a,b]$ is a finite union of intersections with intervals
$[p^{1/C_n},(p+1)^{1/C_n})$, where
$a^{C_n}-1\le p\le b^{C_n}$.
For large $n$, each such interval has length at most
$b/(C_np)$ by the mean value theorem.  Mertens' theorem gives
\[
 \frac b{C_n}\sum_{a^{C_n}-1\le p\le b^{C_n}}\frac1p
 =\frac b{C_n}
   \left(\log\frac{\log b}{\log a}+o(1)\right)
 \longrightarrow0.
\]
Closing a finite union of these intervals adds only finitely many
endpoints.  The closure of the set satisfying every prime condition
in $[a,b]$ is contained in each of these closed unions, and is
therefore null.  To control the closure of the full set on a given
compact subinterval of $(1,\infty)$, apply this argument on a
slightly larger compact interval.  A countable cover, together
with the possible boundary point $1$, proves that the full closure
is null.  A closed null set has empty interior.
\end{proof}

\begin{theorem}[Transfer from one quantitative bad-node estimate]
\label{hb:thm:transfer}
Let $1<c_-\le c_+<\infty$.  Suppose there are constants
$\delta\in(0,1]$, $K>0$, $Y_0$ and
\[
                       D>\frac1{2\log c_-}
 \tag{T}\label{hb:eq:accuracythreshold}
\]
such that, for every $Y\ge Y_0$ and every $c\in[c_-,c_+]$,
\[
 \#\left\{m\in[Y,2Y)\cap\Z:
  \#\PP_c(m)<\delta\frac{|J_c(m)|}{\log((m+1)^c)}
       \right\}
 \le K Y\exp\!\bigl[-D(\log\log Y)^2\bigr].
 \tag{S}\label{hb:eq:nodehyp}
\]
Then, for every sequence $\mathbf c$ taking values in $[c_-,c_+]$,
$\W(\mathbf c)$ has cardinality $2^{\aleph_0}$, is unbounded,
has Lebesgue measure zero, and is nowhere dense.
\end{theorem}

\begin{proof}
Fix the constants in the theorem and then choose a large $P$.
No parameter below depends on the eventual depth of the tree.
Let
\[
 V_1=\PP\cap[P,2P],\qquad
 \Ch_n(m)=\PP_{c_{n+1}}(m).
\]
Call $m$ good at level $n$ when it satisfies the inequality
opposite to that defining the exceptional set in
\eqref{hb:eq:nodehyp}, with $c=c_{n+1}$.
Write
\[
 R_1=1,\qquad R_n=c_2\cdots c_n=C_n/c_1.
\]

\paragraph{Ranges and branching weights.}
Every $p\in V_n$ satisfies
\[
 P^{R_n}\le p,\qquad p+1\le(2P+1)^{R_n}.
 \tag{R}\label{hb:eq:levelrange}
\]
Indeed, the lower inequality propagates through
$q\ge p^{c_{n+1}}$, and the upper one propagates through
$q+1<(p+1)^{c_{n+1}}$.  These are direct inductions from $V_1$.
For a good parent $m$, \eqref{hb:eq:length} and the definition
of goodness give, uniformly for $c=c_{n+1}\in[c_-,c_+]$,
\[
 \begin{split}
 w(\Ch_n(m))
 &\ge
 \frac{\delta\,(c\,m^{c-1}-1)}
      {c\log(m+1)(m+1)^c}\\
 &\ge\frac{\delta}{4m\log(m+1)}
 \ge\frac{\lambda_n}{m},
 \qquad
 \lambda_n=\frac{\delta}{4R_n\log(4P)}.
 \end{split}
 \tag{W}\label{hb:eq:localweight}
\]
The second inequality holds for a single sufficiently large $P$,
uniformly on the fixed compact range of $c$.
The prime number theorem also gives
\[
                       w(V_1)\ge\frac1{4\log P}.
 \tag{P}\label{hb:eq:initialmass}
\]

\paragraph{Defect weights.}
The range in \eqref{hb:eq:levelrange} is covered by at most
$3R_n$ successive dyadic intervals whose lower endpoints
are at least $P^{R_n}$, once $P\ge2$.
Divide \eqref{hb:eq:nodehyp} by the lower endpoint of each
interval and sum.  With $\ell=\log P$ and $L=\log\ell$,
this gives the uniform upper bound
\[
 w(V_n\setminus G_n)
 \le 3K R_n\exp\!\bigl[-D(L+\log R_n)^2\bigr]
 =:\beta_n.
 \tag{D}\label{hb:eq:defectweight}
\]
All analytic hypotheses apply at all levels as soon as
$P\ge Y_0$; a new scale threshold is not chosen for each $n$.

\paragraph{One starting scale for every depth.}
Use the reserves $s_n$ from \cref{hb:cor:binary}.
Put $a=\log c_->0$, $r_n=\log R_n$, and $j=n-1$.
Since consecutive increments of $r_n$ are at least $a$,
\[
 j\le r_n/a,\qquad
 r_n-r_k\ge(n-k)a\quad(k<n).
\]
In particular,
\[
 \sum_{k<n}r_k
 \le j r_n-\frac{a j(j+1)}2
 \le\frac{r_n^2}{2a}.
 \tag{Q}\label{hb:eq:sumlogs}
\]
The last inequality follows by completing the square and dropping
the nonpositive term $-aj/2$.
Equations \eqref{hb:eq:localweight} and
\eqref{hb:eq:initialmass} therefore imply
\[
 -\log s_n
 \le L+O(1)+j(L+O(1))+\sum_{k<n}r_k
 \le\frac{r_n^2}{2a}+\frac{r_n L}{a}
             +O(r_n+L+1).
 \tag{M}\label{hb:eq:reservelog}
\]
All implied constants here are fixed and independent of $n$, $P$
and the particular sequence in $[c_-,c_+]$.
Combining \eqref{hb:eq:defectweight} and
\eqref{hb:eq:reservelog} yields
\[
 \begin{split}
 \log(\beta_n/s_n)
 &\le
 -D(L+r_n)^2+\frac{r_n^2}{2a}
           +\frac{r_n L}{a}+O(r_n+L+1)\\
 &=
 -\left(D-\frac1{2a}\right)(L+r_n)^2
 -\frac{L^2}{2a}+O(r_n+L+1).
 \end{split}
 \tag{U}\label{hb:eq:uniformcomparison}
\]
By \eqref{hb:eq:accuracythreshold}, the coefficient of the first
quadratic term is strictly negative.  Since $L+r_n\ge L$, choosing
$L$ sufficiently large makes this upper bound at most $-\log2$
simultaneously for every $r_n\ge0$.
Thus $\beta_n\le s_n/2$ at every level.

The atom bound is also uniform.  By \eqref{hb:eq:levelrange},
\[
 \max_{p\in V_n}\frac1p\le
          \exp(-\ell e^{r_n}).
\]
For $r\ge0$, $e^r\ge1+r+r^2/2$.  Because $\ell=e^L$,
a single sufficiently large $L$ makes the coefficients of
$\ell(1+r+r^2/2)$ dominate those of
\[
 \frac{r^2}{2a}+\frac{rL}{a}+O(r+L+1)
\]
for every $r\ge0$.  Consequently the last displayed atom bound
is at most $s_n$ for all $n$.
Both conditions in \eqref{hb:eq:binaryconditions} now hold
with the same $P$.

\paragraph{Branches, representing constants, and size.}
\Cref{hb:cor:binary} supplies $2^{\aleph_0}$ distinct prime
branches.  By \cref{hb:lem:nesting}, each gives a distinct member
of $\W(\mathbf c)$.  Since this is a subset of $\R$, its
cardinality is exactly $2^{\aleph_0}$.
Every such constant satisfies $A\ge P^{1/c_1}$.
The construction works for all sufficiently large $P$, so
$\W(\mathbf c)$ is unbounded.
Finally $C_n\ge c_-^n\to\infty$, and
\cref{hb:prop:null} supplies nullity and nowhere density.
\end{proof}

\begin{remark}[Comparison with the earlier sufficient threshold]
\label{hb:rem:oldthreshold}
The earlier comparison bounded $R_n$ above and below separately
and imposed
\[
 D>\frac{\log c_+}{2(\log c_-)^2}.
\]
That condition remains sufficient: $\log c_+\ge\log c_-$ implies
that it is at least as strong as \eqref{hb:eq:accuracythreshold}.
It also forces the mixed coefficient $2D\log c_--1$ to be positive.
The sharper estimate \eqref{hb:eq:sumlogs} compares the accumulated
loss with the actual current scale and avoids the artificial
$c_+$ loss.  This refines the required exceptional accuracy, without
changing the prime-interval exponent.
\end{remark}

\begin{corollary}[Transfer of a quantitative almost-all prime input]
\label{hb:cor:AAtransfer}
Assume $\mathsf{AA}(\theta)$.
If $\mathbf c$ satisfies
\[
 \frac1{1-\theta}<\inf_k c_k\le\sup_k c_k<\infty,
\]
then $\W(\mathbf c)$ has cardinality $2^{\aleph_0}$, is unbounded,
is Lebesgue-null, and is nowhere dense.
\end{corollary}

\begin{proof}
Set $c_-=\inf c_k$ and $c_+=\sup c_k$.
Choose one fixed $D>1/(2\log c_-)$.
\Cref{hb:prop:bad} supplies \eqref{hb:eq:nodehyp} for this $D$,
with a fixed $\delta_D>0$.  Apply \cref{hb:thm:transfer}.
\end{proof}

\begin{remark}[Quantifiers and endpoints]\label{hb:rem:quantifiers}
The order is: fix the exponent range, choose the exceptional
accuracy $D$, obtain the fixed prime lower-bound constant, and
then choose one starting scale $P$ for the whole infinite tree.
A bound $O_B(X\log^{-B}X)$ separately for each fixed $B$ does not
by itself justify this order, or permit replacing $B$ by a function
of $X$ or the depth.
The strict condition
$\theta<1-1/c_-$ is used in the full-cell averaging step.
Thus an input at exactly $\theta$ gives $c>1/(1-\theta)$.
An improvement to $\theta<1/24$ includes $c=24/23$ because
the needed inequality is then strict.
The stronger Vinogradov--Korobov exceptions in
\cref{sec:density} also meet the quantitative node hypothesis.
\end{remark}

\section{Primes in almost all short intervals with Vinogradov--Korobov savings}\label{sec:density}

Let $N(\sigma,T)$ denote the number of zeros $\rho=\beta+i\gamma$ of $\zeta(s)$ with $\beta\ge\sigma$ and
$|\gamma|\le T$. Write $\mathcal L(X)=(\log X)^{1/3}(\log\log X)^{-1/3}$.

\begin{hypothesis}[$\mathrm Z(\theta_1)$]\label{hyp:Z}
Let $0<\theta_1<1$. There are $\sigma_1\in[1/2,1)$, $\delta_0>0$ and $B\ge0$ such that, with
$a_*=2/(1-\theta_1)-\delta_0$,
\begin{enumerate}[label=\textup{(\alph*)}]
\item for every $\eps>0$: $N(\sigma,T)\ll_\eps T^{a_*(1-\sigma)+\eps}$ uniformly for $1/2\le\sigma\le\sigma_1$, $T\ge2$;
\item $N(\sigma,T)\ll T^{a_*(1-\sigma)}(\log T)^B$ uniformly for $\sigma_1\le\sigma\le1$, $T\ge2$.
\end{enumerate}
\end{hypothesis}

We use the following classical facts.
\begin{itemize}[leftmargin=1.5em]
\item[(F1)] Explicit formula \cite[Ch.~17]{Dav00}: for $x\ge2$ and $T\ge2$,
$\psi(x)=x-\sum_{|\gamma|<T}x^\rho/\rho+O\bigl(x(\log xT)^2/T+\log x\bigr)$.
\item[(F2)] Zero counting \cite[Ch.~15]{Dav00}: $N(t+1)-N(t)\ll\log(|t|+2)$ for the number $N(t)$ of zeros with $0<\gamma\le t$.
\item[(F3)] Vinogradov--Korobov \cite[Thm.~6.1]{Ivi85}, \cite{Ford}: there is $c_0>0$ such that $\zeta(\sigma+it)\neq0$ for
$\sigma\ge1-c_0(\log|t|)^{-2/3}(\log\log|t|)^{-1/3}$, $|t|\ge3$.
\item[(F4)] Ingham \cite{Ing40}, \cite[Ch.~11]{Ivi85}: $N(\sigma,T)\ll T^{3(1-\sigma)/(2-\sigma)}(\log T)^5$ for $1/2\le\sigma\le1$.
\item[(F5)] Huxley \cite{Hux72}, \cite[Ch.~11]{Ivi85}: $N(\sigma,T)\ll T^{3(1-\sigma)/(3\sigma-1)}(\log T)^{44}$ for $3/4\le\sigma\le1$.
\item[(F6)] Guth--Maynard \cite{GM}: for each fixed $\sigma\in[7/10,8/10]$
and every $\eps>0$, $N(\sigma,T)\ll_{\sigma,\eps} T^{15(1-\sigma)/(3+5\sigma)+\eps}$.
\end{itemize}

\begin{lemma}[verification of $\mathrm Z(\theta_1)$]\label{lem:Z}
\begin{enumerate}[label=\textup{(\roman*)}]
\item $\mathrm Z(\theta_1)$ holds for every $\theta_1>1/6$, using \textup{(F4)} and \textup{(F5)} only.
\item $\mathrm Z(\theta_1)$ holds for every $\theta_1>2/15$, using \textup{(F4)}, \textup{(F5)} and \textup{(F6)}.
\item Under the density hypothesis, $\mathrm Z(\theta_1)$ holds for every $\theta_1>0$.
\end{enumerate}
\end{lemma}

\begin{proof}
(i) On $[1/2,3/4]$ the exponent $3/(2-\sigma)$ of (F4) is at most $12/5$, and on $[3/4,1]$ the exponent
$3/(3\sigma-1)$ of (F5) is at most $12/5$. Hence $N(\sigma,T)\ll T^{12(1-\sigma)/5}(\log T)^{44}$ on $[1/2,1]$.
If $\theta_1>1/6$ then $12/5<2/(1-\theta_1)$; take $\sigma_1=1/2$ and $\delta_0=2/(1-\theta_1)-12/5$ in (b).

(ii) Let $\theta_1>2/15$, so that $30/13<2/(1-\theta_1)$. On $[1/2,7/10]$, (F4) gives exponent
$3/(2-\sigma)\le30/13$ with a logarithmic factor; on $[7/10,4/5]$, (F6) gives $15/(3+5\sigma)\le30/13$ with a
factor $T^\eps$; on $[4/5,1]$, (F5) gives $3/(3\sigma-1)\le 15/7<30/13$ with a logarithmic factor. Take
$\sigma_1=4/5$ and $\delta_0=\tfrac12\bigl(2/(1-\theta_1)-30/13\bigr)$, so that $a_*=30/13+\delta_0$. Then (b) holds on
$[4/5,1]$ directly from (F5) with $B=44$, because $15/7<a_*$; on $[1/2,7/10]$ the logarithmic factor in (F4) is
$\ll_\eps T^\eps$, which gives (a) there. For (a) on $[7/10,4/5]$, where (F6) is a pointwise statement, use that $N(\sigma,T)$ is nonincreasing in
$\sigma$: for a mesh $7/10=s_0<s_1<\dots<s_r=4/5$ with $s_{j+1}-s_j\le\delta_0/20$ and $\sigma\in[s_j,s_{j+1}]$,
$N(\sigma,T)\le N(s_j,T)\ll T^{\frac{30}{13}(1-s_j)+\eps}\le T^{\frac{30}{13}(1-\sigma)+\frac{30}{13}\cdot\frac{\delta_0}{20}+\eps}$, and the extra
exponent is absorbed into $\delta_0(1-\sigma)$ because $1-\sigma\ge1/5$ and $\frac{30}{13}\cdot\frac1{20}<\frac15$.

(iii) Let $\theta_1>0$. The density hypothesis gives (a) with exponent $2<2/(1-\theta_1)$ on $[1/2,1]$. For (b)
choose $\sigma_1\in[3/4,1)$ with $3/(3\sigma_1-1)<2/(1-\theta_1)$, which is possible because $3/(3\sigma-1)\to3/2$ as
$\sigma\to1$, and use (F5) on $[\sigma_1,1]$. Take $\delta_0$ smaller than both gaps. If the density hypothesis is
only available pointwise in $\sigma$, the mesh argument of (ii) gives the required uniformity.
\end{proof}

\begin{proposition}[mean square with VK savings; classical, cf.\ \cite{SV77,Zac98,BP00,GT25}]\label{prop:ms}
Assume $\mathrm Z(\theta_1)$ and let $\theta_1\le\theta_2<1$. There are $\kappa>0$ and $X_0$ such that for
$X\ge X_0$ and $X^{\theta_1}\le h\le X^{\theta_2}$,
\[
 \int_X^{2X}\bigl|\psi(u+h)-\psi(u)-h\bigr|^2\,du\le h^2X\exp\bigl(-\kappa\,\mathcal L(X)\bigr).
\]
\end{proposition}

\begin{proof}
Let $T=X^{1-\theta_1}\exp(2\sqrt{\log X})$, replacing it if necessary by a height in $[T,T+1]$ away from zero ordinates; then $2\le T\le X$ for large $X$. By (F1), for $u\in[X,2X]$,
\begin{gather*}
 \psi(u+h)-\psi(u)-h=-S(u)+O\Bigl(\frac{X\log^2X}{T}+\log X\Bigr),\\
 S(u)=\sum_{|\gamma|<T}\frac{(u+h)^\rho-u^\rho}{\rho}=\int_u^{u+h}F(v)\,dv,
\end{gather*}
with $F(v)=\sum_{|\gamma|<T}v^{\rho-1}$. The error term is at most $h\exp(-\sqrt{\log X})$ for large $X$, since
$X\log^2X/T=X^{\theta_1}\log^2X\,e^{-2\sqrt{\log X}}$ and $h\ge X^{\theta_1}$.

By Cauchy--Schwarz, $|S(u)|^2\le h\int_u^{u+h}|F(v)|^2dv$, and integrating over $u\in[X,2X]$ (each $v$ lies
in at most $h$ units of $u$-measure),
\[
 \int_X^{2X}|S(u)|^2du\le h^2\int_X^{3X}|F(v)|^2dv=h^2\sum_{\rho,\rho'}\int_X^{3X}v^{\rho+\overline{\rho'}-2}dv .
\]
For $\rho=\beta+i\gamma$ and $\rho'=\beta'+i\gamma'$ the inner integral is at most $2X^{\beta+\beta'-1}$ in
absolute value, and also at most $6X^{\beta+\beta'-1}/|\gamma-\gamma'|$ by integrating the oscillating factor.
Hence it is $\le12X^{\beta+\beta'-1}/(1+|\gamma-\gamma'|)$. With $X^{\beta+\beta'}\le(X^{2\beta}+X^{2\beta'})/2$
and symmetry,
\[
 \int_X^{3X}|F|^2\le 12\sum_{|\gamma|<T}X^{2\beta-1}\sum_{|\gamma'|<T}\frac1{1+|\gamma-\gamma'|}
 \ll(\log T)^2\sum_{|\gamma|<T}X^{2\beta-1},
\]
by (F2). Consequently
\[
 \int_X^{2X}|\psi(u+h)-\psi(u)-h|^2du\ll h^2X\Bigl((\log X)^2\Sigma+e^{-2\sqrt{\log X}}\Bigr),
 \qquad \Sigma=\sum_{|\gamma|<T}X^{2\beta-2}.
\]
Zeros with $\beta<1/2$ contribute at most $N(T)X^{-1}\ll T\log T/X\le X^{-\theta_1/2}$ to $\Sigma$.
For $\beta\ge1/2$, partial summation gives
\[
 \sum_{\beta\ge1/2}X^{2\beta-2}=X^{-1}N(\tfrac12,T)+2\log X\int_{1/2}^1X^{2\sigma-2}N(\sigma,T)\,d\sigma .
\]
The endpoint term $X^{-1}N(1/2,T)\ll T\log T/X$ has the same fixed power saving as the zeros with $\beta<1/2$.
By (F3), after decreasing $c_0$ to cover the finitely many bounded-height zeros, $N(\sigma,T)=0$ for $\sigma>1-\eta$ with $\eta=c_0(\log T)^{-2/3}(\log\log T)^{-1/3}$.
Write $a_*(1-\theta_1)=2-\delta_1$ with $\delta_1=\delta_0(1-\theta_1)>0$. Then
$T^{a_*(1-\sigma)}X^{2\sigma-2}=\exp\bigl(-(1-\sigma)[\delta_1\log X-2a_*\sqrt{\log X}]\bigr)
\le\exp(-(1-\sigma)\tfrac{\delta_1}{2}\log X)$ for large $X$.
On $[\sigma_1,1-\eta]$ hypothesis (b) bounds the integrand by $(\log X)^B\exp(-\tfrac{\delta_1}{2}\eta\log X)$.
On $[1/2,\sigma_1]$ hypothesis (a), with $\eps=(1-\sigma_1)\delta_1/4$, bounds it by $X^{-(1-\sigma_1)\delta_1/4}$.
Since $\log T\asymp\log X$, we have $\eta\log X\gg\mathcal L(X)$, and therefore
$\Sigma\ll(\log X)^{B+1}\exp(-c'\mathcal L(X))$ for some $c'>0$. The proposition follows with any
$\kappa<c'$.
\end{proof}


\begin{proof}[Proof of \cref{thm:baseline}]
For the applicable density threshold $\theta_0\in\{1/6,2/15,0\}$,
choose fixed exponents
$\theta_0<\theta_1<\vartheta<1-1/c_-$ and then
$\vartheta<\theta_2<1$. The preceding lemma verifies
$\mathrm Z(\theta_1)$, with the density hypothesis assumed only in
the third case. Apply the mean-square proposition with
$h=X^\vartheta$. Chebyshev's inequality gives a prime-weighted lower
bound outside measure $O(Xe^{-a\LL(X)})$ for some $a>0$.

Prime powers of exponent at least two contribute uniformly
\[
 O\bigl(hX^{-1/2}(\log X)^2+(\log X)^2\bigr)=o(h)
\]
to the difference of $\psi$ on an interval of length $h$ at height
comparable to $X$: count the possible $k$th roots in each interval
and sum over $2\le k\ll\log X$.
Thus the number of primes is at least a fixed positive multiple of
$h/\log X$. Translation and a bounded number of comparable dyadic
ranges give the same assertion for backward windows
$[x-X^\vartheta,x]$.
Since $e^{-a\LL(X)}$ is smaller than
$\exp[-D(\log\log X)^2]$ for every fixed $D$, the quantitative
short-interval hypothesis in \cref{hb:sec:transfer} applies.
Its bounded-variable-exponent conclusion proves all three assertions.
\end{proof}

\section{Why fixed logarithmic estimates cannot simply be substituted}

For each separately fixed $B$, the stated almost-all sieve theorem supplies exceptional
measure $O_B(X(\log X)^{-B})$.
The harmonic construction needs a uniform estimate of the form
\[
             X\exp(-D(\log\log X)^2).
\]
The first assertion alone does not imply the second.
For example,
\[
             r(X)=\exp(-(\log\log X)^{3/2})
\]
is $O_B((\log X)^{-B})$ for every fixed $B$, but exceeds
$\exp(-D(\log\log X)^2)$ eventually for every fixed $D>0$.

This distinction has a concrete tree interpretation.
At level $n$, $\log Y_n$ is comparable to $c^n$, and repeatedly retaining
a prime proportion of size $1/\log Y_n$ loses
\[
 \prod_{j<n}\frac1{\log Y_j}
       =\exp\left(-\frac{\log c}{2}n^2+O(n)\right).
\]
A fixed logarithmic exceptional rate gives only
$\exp(-Bn\log c+O(1))$, which eventually exceeds this retained mass.
The analytic proof must therefore be made quantitative before it is
combined with the tree argument.

For $D>0$ define, for sufficiently large $X$,
\[
 \Q_D(X)=\exp(D(\log\log X)^2),\qquad
 \LL(X)=\frac{(\log X)^{1/3}}{(\log\log X)^{1/3}}.
\]
For any fixed $D,K,d,\nu>0$,
\begin{equation}\label{eq:hierarchy}
 (\log X)^d=\Q_D(X)^{o(1)},\quad
 \Q_D(X)^K=X^{o(1)},\quad
 \Q_D(X)^K e^{-\nu\LL(X)}=o(\Q_D(X)^{-d}).
\end{equation}
The order of choices below is essential: fix the interval exponent,
then fix all sieve and moment parameters, then choose the finite
accuracy powers, and finally let $X$ increase.

\section{A quantitative prime-polynomial estimate}

\begin{lemma}[uniform prime-polynomial bound]\label{lem:vk}
Fix $\delta>0$, $M_0>0$, and a bounded real interval of values of $\sigma$.
There is $a>0$ such that, whenever $X$ is sufficiently large,
$X^\delta\le u\le v\le2u\ll X$, and $|t|\le X^{M_0}$,
\[
 \sum_{u<n\le v}\frac{\Lambda(n)}{n^{\sigma+it}}
   =\int_u^v y^{-\sigma-it}\,dy
     +O_{\delta,M_0,\sigma}\left(u^{1-\sigma}e^{-a\LL(X)}\right).
\]
The constants are uniform when $\sigma$ stays in the indicated compact
interval. Consequently
\[
 \left|\sum_{u<n\le v}\frac{\Lambda(n)}{n^{\sigma+it}}\right|
 \ll u^{1-\sigma}
       \left(\frac1{1+|t|}+e^{-a\LL(X)}\right).
\]
The analogous upper bound for primes with coefficient $1$ holds, with at
most a fixed logarithmic factor.
\end{lemma}

\begin{proof}
Let $s=\sigma+it$. Apply truncated Perron's formula to
$-\zeta'(s+w)/\zeta(s+w)$ with difference kernel
\[
                    \frac{v^w-u^w}{w}.
\]
Take the initial line with $\Re(s+w)>1$, and choose a truncation height
$X^M$ with fixed $M$ sufficiently large relative to $M_0$ and the allowed
range of $\sigma$.
The endpoint convention for Perron's formula changes the normalized
sum by $O(u^{-\sigma}\log X)$, which is absorbed by the asserted error.

The Vinogradov--Korobov zero-free region~\cite{Ford,Ivi85} permits a shift to
$\Re(s+w)=1-\eta$, where
\[
 \eta=b_0(\log X)^{-2/3}(\log\log X)^{-1/3}
\]
and $b_0>0$ is sufficiently small. Decreasing $b_0$ handles bounded
heights as well. In this narrower zero-free strip,
$-\zeta'/\zeta(z)-(z-1)^{-1}$ has a fixed power of $\log X$ as an upper
bound. The only pole crossed is at $w=1-s$, and its residue is
\[
              \frac{v^{1-s}-u^{1-s}}{1-s}=\int_u^v y^{-s}\,dy.
\]
The difference kernel is regular at $w=0$, so it creates no additional
residue there. The shifted vertical integral is
\[
 \ll u^{1-\sigma-\eta}(\log X)^{O(1)}
 \ll u^{1-\sigma}e^{-a\LL(X)},
\]
because $\log u\ge\delta\log X$. The horizontal integrals and Perron
truncation error are smaller on increasing $M$.
This proves the expansion.

For the main integral, writing $y=ue^z$ shows that its size is
$O(u^{1-\sigma}/(1+|t|))$, uniformly for $0\le\log(v/u)\le\log2$.
Removing prime powers costs at most a fixed logarithmic factor times
$u^{1/2-\sigma}+u^{-\sigma}$, which is negligible here.
Partial summation removes the remaining $\log p$ weight.
\end{proof}

\begin{corollary}\label{cor:Qinput}
For every fixed $D,K>0$, the prime polynomials in \cref{lem:vk} satisfy
\[
 \left|\sum_{u<n\le v}\Lambda(n)n^{-\sigma-it}\right|
 \ll u^{1-\sigma}\Q_D(X)^{-K}
\]
throughout $\Q_D(X)^{K+1}\le |t|\le X^{M_0}$, for sufficiently large $X$.
The low-frequency main integral is retained rather than discarded.
\end{corollary}

\begin{lemma}[fixed products of large primes]\label{lem:products}
The same conclusion, with a possible fixed logarithmic loss, holds for
Dirichlet polynomials formed from a fixed number of prime variables
$p_j\ge X^\delta$, with fixed order, range, or monomial constraints,
provided the selected prime variable has no arbitrary oscillating
coefficient.
It continues to hold after convolution with bounded sieve coefficients
in other variables, with fixed divisor and logarithmic losses.
\end{lemma}

\begin{proof}
Fix all but one prime variable and let $q$ be their product.
The allowed values of the remaining prime lie in a bounded union of
intervals inside $[M/q,2M/q]$, subject to a lower bound $X^\delta$.
Use \cref{lem:vk} on that variable. Its error, multiplied by $q^{-\sigma}$,
is $O(M^{1-\sigma}q^{-1}e^{-a\LL(X)})$.
The sum of $1/q$ over a fixed number of primes in $[X^\delta,O(X)]$ is
bounded in terms of $\delta$ and the number of variables, by Mertens'
estimate. Fixed logarithmic weights add only fixed logarithmic factors.

One can fix the other sieve variables first as well. Bounded sieve
coefficients have harmonic absolute sums bounded by fixed logarithmic
powers; grouping a fixed number of variables gives a fixed divisor
bound. These factors are absorbed using \eqref{eq:hierarchy}.
This statement does not assert cancellation for arbitrary coefficients
on the prime variable itself.
\end{proof}

\section{Earlier finite-sieve specializations}\label{sec:earlier-sieves}

We use Li's finite decomposition at the short-interval exponent $2/43$.
Its combinatorial regions and final numerical integrals are unchanged.
The issue is the analytic error budget.

\begin{proposition}[accuracy parameter in the finite sieve]\label{prop:finiteQ}
Fix a sufficiently small $\eps>0$ in the argument of \cite{LiI}, and put
\[
            \eta=\frac12X^{-41/43+\eps}.
\]
For every fixed $D>0$, there exist $X_0$ and $\kappa>0$ such that
\begin{equation}\label{eq:sieveQ}
 \meas\left\{x\in[X,2X]:
   \pi(x+\eta x)-\pi(x)<\kappa\frac{\eta x}{\log X}\right\}
 \le X\Q_D(X)^{-1}
 \qquad (X\ge X_0).
\end{equation}
Here the finite region inequalities (17), (18), and (20) of \cite{LiI}
are used as external numerical inputs.
\end{proposition}

\begin{proof}
Fix $\eps$ first; in particular $\delta=\eps^{1/3}$ is fixed.
Whenever necessary, take $\eps$ smaller so that the deterministic
sieve approximation errors fit strictly inside Li's positive margin.
No accuracy parameter below changes $\delta$, a moment order, or the
number of Buchstab iterations.

Write $Q=\Q_D(X)$.
We rerun the finite analytic estimates with an auxiliary accuracy
$Q^{-R}$, where $R$ is a sufficiently large fixed constant.
This is not an application of Li's theorem with a moving $B$.
The following checks show directly that its proof allows the new
accuracy.

\paragraph{1. Pointwise input.}
Replace the low-frequency cutoff by $T_0=Q^{R_0+2}$ with a sufficiently
large fixed $R_0$. Prime factors satisfy every fixed power of $Q^{-1}$
above that cutoff by \cref{cor:Qinput,lem:products}.
In particular, the saving $Q^{-R_0}$ remains available after the
fixed logarithmic losses.
Below it, keep the main integral in \cref{lem:vk}.
The contribution of its remainder is bounded by a fixed power of
$T_0$ and $\log X$ times $e^{-a\LL(X)}$, and is therefore smaller than
any prescribed fixed power of $Q^{-1}$.

\paragraph{2. Mean and large values.}
The estimates in Li's Lemmas 3.1--3.6 use fixed moments and fixed
interpolations of classical mean-value and large-value inequalities.
In Lemma 3.1 the integer moment satisfies $k\ll1/\delta$, and the
logarithmic loss has exponent $O(\delta^{-2})$; both are fixed.
Every resulting case has one of two forms:
\[
     (\log X)^d Q^{-\alpha R_0}
             \quad\hbox{or}\quad
     (\log X)^d X^{-\nu},
\]
where $d$ is fixed, $\alpha>0$, and $\nu>0$ depends on the fixed
exponent margin. There are finitely many such cases.
Thus the minimum relevant positive $\alpha$ can be chosen before
$R_0$. By increasing $R_0$ and then $X$, every one is
$O(Q^{-R})$ for any specified fixed $R$.
The extra finite polynomial products in Lemmas 3.3 and 3.4 do not
alter this conclusion. Lemmas 3.5 and 3.6 are regroupings and
permutations of Lemma 3.2 and preserve the same accuracy accounting.

For completeness, one displayed interpolation in Lemma 3.2,
case (1)(a), can be justified without its copied intermediate line.
If the simultaneous large-value count $F$ is bounded by
\[
 V^{-6}M^3,\quad W^{-10}H^5,\quad U^{-4}K^2,
 \qquad MHK=X,
\]
interpolate with weights $3/10$, $1/5$, and $1/2$. Then
\begin{align*}
 U^2V^2W^2F
 &\ll U^2V^2W^2
     (V^{-6}M^3)^{3/10}
     (W^{-10}H^5)^{1/5}(U^{-4}K^2)^{1/2}\\
 &=V^{1/5}M^{9/10}HK.
\end{align*}
If $V\le M^{1/2}Q^{-R_0}$, this is at most
$XQ^{-R_0/5}$. The positive power of the pointwise saving is all that
is needed. This is the interpolation already present in Jia's
Lemma 5, case I(a), printed page 28.

\paragraph{3. The Watt-type estimates.}
Li's Section 4 uses a fixed exponent margin in its mean-value
estimates. The resulting error is a fixed power saving in $X$,
or one of the small-factor cases already covered in step 2.
Both dominate every required fixed power of $Q^{-1}$ by
\eqref{eq:hierarchy}. The constants in these mean-value estimates
depend on the fixed margins and coefficient bounds, not on a
logarithmic accuracy exponent.

\paragraph{4. Perron and the exceptional measure.}
Here are the details behind the required error bookkeeping.
For a relevant Dirichlet polynomial $F$, put
\[
                W_\eta(s)=\frac{(1+\eta)^s-1}{s}.
\]
On a dyadic range $|t|\asymp T$,
$|W_\eta(b+it)|\ll\min(\eta,T^{-1})$, where $b=1+1/\log X$.
Opening the square, integrating $x^{s+\overline{s'}}$, and applying
$2|ab|\le |a|^2+|b|^2$, gives
\begin{align}
 &\int_X^{2X}\left|
     \int_{T_0\le |t|\le X}F(b+it)W_\eta(b+it)x^{b+it}\,dt
                       \right|^2dx \nonumber\\
 &\hspace{7mm}\ll X^3(\log X)^3
       \max_{T_0\le T\le X}
       \min(\eta,T^{-1})^2
          \int_{T\le |t|\le2T}|F(b+it)|^2dt. \label{eq:perronL2}
\end{align}
This is the estimate used in Jia's Lemmas 16--19~\cite{Jia}, whose
Perron steps underlie Li's Section 5.
If the right-hand maximum is $\ll\eta^2Q^{-R}$, Chebyshev's inequality,
with threshold $\eta XQ^{-r}$, yields exceptional measure
\[
                       \ll X(\log X)^3Q^{-R+2r}.
\]
Choose $R>2r+10$ after all fixed decompositions. This leaves more
than enough accuracy for \eqref{eq:sieveQ}.
The low-frequency terms are evaluated as main terms, with errors
$\ll\eta XQ^{-r}$ by step 1.
Terms such as $\eta^2XT_0^j$ and Perron truncation errors are harmless
since $T_0^j=X^{o(1)}$ and $\eta X$ is a fixed positive power of $X$.
Endpoint conventions contribute smaller errors as well.
The same argument with the appropriate zeta factor handles the
Type I terms.

\paragraph{5. Separation of the finite constraints.}
The number of prime variables and the number of defining inequalities
are fixed. Standard dyadic partitions introduce only a fixed power
of $\log X$. Alternatively, a relative mesh $Q^{-r_1}$ separates the
variables into at most $Q^{O_\eps(r_1)}$ boxes.
Boundary strips in fixed monomial and ordering constraints have
total tuple weight $\ll XQ^{-r_1}(\log X)^L+X^{1-\delta_0}(\log X)^L$
for fixed $L$ and $\delta_0>0$. The divisor short-interval estimate
proved in \cref{app:reproduce} justifies this bound after all but one
variable are fixed.
Fubini and Markov transfer it to the short sums.
Choose $r_1$ for the needed boundary accuracy first, and then choose
$R$ large enough to cover the number of boxes.
This avoids a circular choice of accuracy powers.
For the lower bound one may also shrink the admissible polyhedral
regions by a sufficiently small fixed amount: their boundary loss is
absorbed in the fixed positive numerical margin.

\paragraph{6. The sieve and its positive margin.}
In Li's Lemma 5.6 the sieve level $X^{10^{-1000}}$ and the sifting
threshold $X^\delta$ are fixed powers of $X$.
The bounded sieve coefficients and the grouping of finitely many
variables preserve fixed divisor bounds. The deterministic errors of
size $O(\eps)$ in the normalized main term need not become $Q^{-1}$:
choose $\eps$ small once. Only the fluctuating error and its exceptional
set require the stronger accuracy.

In applying the small-factor estimates, a coefficient supported on
rough integers is not automatically declared to have cancellation.
The relevant groups contain prescribed prime factors, to which
\cref{lem:products} applies after the remaining variables are fixed.
If the group denoted $N$ in Li's Lemma 4.4 consists only of the tiny
sieve variable, use the Type I estimate in Lemma 5.1 instead:
the relevant range has $QR<X^{45/86}$ and $N<X^{41/344}$.
Thus the proof does not need a pointwise estimate for arbitrary
sieve weights.

Li's recorded numerical region estimates give the positive margin
\[
                  1-0.03-0.992+0.027=0.005.
\]
Choose the fixed deterministic errors and region-boundary losses less
than half this margin. Steps 1--5 make all remaining errors smaller
than a further fixed portion outside measure $XQ^{-1}$, for sufficiently
large $X$. The same fixed sieve decomposition therefore gives
\eqref{eq:sieveQ}, with a positive constant $\kappa$.
The numerical inequalities themselves are precisely the external
component specified in the proposition.
\end{proof}

\begin{remark}[what is and is not imported]
The standard statement $O_B(X/\log^B X)$ is not the imported
quantitative conclusion. We use the finite algebraic estimates and
the fixed positive sieve margin, and verify their accuracy dependence.
The general free-parameter variance estimate in
Matomäki--Radziwiłł--Shao--Tao--Teräväinen~\cite[Lemma 3.3]{MRSTT}
provides a separate modern formulation of the Perron/variance step.
Their other results are not used to claim a new short-interval exponent.
\end{remark}

\subsection{The published-source alternative}

\begin{proposition}[the published-source alternative]\label{prop:jiaQ}
For sufficiently small fixed $\eps>0$, the counterpart of
\cref{prop:finiteQ} holds with
$\eta=\tfrac12X^{-19/20+\eps}$, using Jia's published proof
in place of Li's refined decomposition.
\end{proposition}

\begin{proof}
Jia fixes $\eps$, $\eps_1=\eps^2$ and $\delta=\eps^{1/3}$ before
the accuracy parameter. His mean-value arguments in Sections 2--5
have the same structure as steps 1--3 in \cref{prop:finiteQ}:
fixed moments, fixed powers of logarithms, and either a positive
power of the pointwise saving or a fixed power saving in $X$.
His Lemmas 16--19 are the Perron calculations quantified directly
in \eqref{eq:perronL2}. The subsequent fixed Buchstab regions,
linear sieve levels and deterministic $O(\eps)$ or $O(\delta)$
errors do not depend on the requested logarithmic accuracy.
Thus the identical order of parameter choices gives an exceptional
measure at most $X\Q_D(X)^{-1}$ for every fixed $D$.
Jia's published positive numerical sieve margin supplies the
constant $\kappa>0$. We use that published margin as an external
input just as in a normal application of his theorem.
\end{proof}


\begin{corollary}[intermediate Mills ranges]\label{cor:earlier-ranges}
Using Jia's finite decomposition and published numerical margin,
the four conclusions about $\W(\boldsymbol c)$ hold for bounded
sequences with $\inf_n c_n>20/19$.
Assuming the three numerical region inequalities (17), (18), and
(20) of Li I, the corresponding lower bound is $43/41$.
\end{corollary}
\begin{proof}
Choose the fixed $\eps>0$ in the appropriate proposition so that
its interval exponent is less than $1-1/c_-$.
The full-cell averaging argument in \cref{hb:sec:transfer} applies
also to the forward multiplicative windows here. Indeed, if
$I=[a,b)$ has length $L\gg\eta X$, integrate over
$[a,b/(1+\eta))$. Each prime is counted for at most $2\eta X$
units of endpoint measure. A cell with fewer than
$\kappa L/(16\log X)$ primes therefore contains at least $L/4$
measure of exceptional endpoints. Disjoint cells give the same
bad-node bound, with arbitrary fixed accuracy, and the harmonic
transfer proves the assertions. The numerical inputs of this
corollary are independent of \cref{hyp:liIII}.
\end{proof}

\section{Why the positive sieve margin alone is insufficient}

Li III has the recorded lower-bound budget
\[
             1-0.02-0.03-0.97+0.04=0.02.
\]
It is tempting to argue that continuity immediately permits a slightly
smaller interval exponent. This misses a uniform analytic input.

Let $MH=X$, $H=X^h$, $M=X^{1-h}$, and put $b=1+1/\log X$.
On a one-separated large-value class write
\[
 V\le\sqrt M\,|M(b+it)|<2V,\qquad
 W\le\sqrt H\,|H(b+it)|<2W,\qquad R=\#\{\text{class}\}.
\]
The classical mean-value and Halasz bounds~\cite{Harman} give, up to fixed powers
of $\log X$,
\begin{equation}\label{eq:classic}
 R\ll\min\left\{
 \frac{M+T}{V^2},\frac M{V^2}+\frac{MT}{V^6},
 \frac{H^k+T}{W^{2k}},
 \frac{H^k}{W^{2k}}+\frac{H^kT}{W^{6k}}\right\}.
\end{equation}
The four-case argument in Li II~\cite[Lemma 3.1]{LiII} bounds $V^2W^2R$ by,
respectively,
\[
 X\cdot(\text{pointwise saving}),\quad
 M^{1/(2k)}T,\quad MH^{1/3}T^{1/(3k)},\quad
 M^{1/(2k)}HT^{1-1/k}.
\]
For intervals of length $X^\theta$, the last three are admissible when
\begin{equation}\label{eq:oldcover}
 h>1-2k\theta,\qquad h>\frac{1-\theta}{2k},\qquad
 h<\frac{1+2(k-1)\theta}{2k-1}.
\end{equation}
At $\theta=1/24$, the $k=11$ and $k=12$ ranges meet at $h=1/12$.
At $\theta=1/24-d$ a gap opens:
\[
 \frac1{12}-\frac{22d}{23}<h<\frac1{12}+22d.
\]
Reducing the outer cutoff $11/126$ by an infinitesimal amount does
not remove this internal gap.

There is also a useful exponent witness.
At $h=1/12$, $T=X^{23/24+d}$, set
\[
 V=X^{11/48},\quad W=X^{1/48-d/24},\quad R=X^{1/2+d}.
\]
Every individual bound in \eqref{eq:classic}, for every positive
integer $k$, permits these exponents, although
$V^2W^2R=X^{1+11d/12}$.
Adding the trivial mean value bound for $MH$ caps this at $X$:
replace $R$ by $X^{1/2+d/12}$. Equality then remains possible at
amplitudes far below their pointwise maxima, so no logarithmic
saving follows from that additional trivial bound.
This is an obstruction to these estimates, not an example of actual
Dirichlet polynomials or a counterexample about primes.

\section{An explicit modern two-factor bridge}

Theorem 1.1 of Chen--Gupta--Li~\cite{CGL}, with $q=q_1=1$, gives
for $N\ge T^{2/3}$ and ordinary polynomial amplitude at least $U$,
\[
 R\ll T^{o(1)}
 \left(\frac{N^2}{U^2}+
 \frac{T^{1/2}N^3}{U^4}+
 \frac{TN^2}{U^4}+
 \frac{TN^{12/5}}{U^4}\right).
\]
If $N\ge T^{5/6}$, the final two terms are dominated by the second.
For a fixed divisor-bounded polynomial of length $M$, normalizing
$U\asymp\sqrt M V$ and absorbing the coefficient bound into an
arbitrarily small fixed power gives
\begin{equation}\label{eq:modern}
             R\ll_\xi X^\xi
                \left(\frac M{V^2}+\frac{MT^{1/2}}{V^4}\right)
 \qquad(\xi>0,\ M\ge T^{5/6}).
\end{equation}
The coefficient normalization is also detailed in
\cite[Lemma 2.3]{MPPTW}. We use the underlying estimate, not an
exponent extrapolation from its applications.

\begin{lemma}[two-factor estimate at $1/25$]\label{lem:pair}
Fix $\delta,D>0$, a desired accuracy $B_0>0$, and a divisor bound
for the coefficients of $M(s)$.
Let $MH=X$, $X^\delta\le H\le X^{11/126}$, and
\[
 H(s)=\sum_{h\asymp H}\Lambda(h)h^{-s},\quad
 Q=\exp[D(\log\log X)^2],\quad \eta=X^{-24/25}.
\]
There is a fixed $A>0$ such that, uniformly for $Q^A\le T\le X$,
\[
 \min(\eta,T^{-1})^2\int_T^{2T}|M(b+it)H(b+it)|^2\,dt
                 \ll \eta^2Q^{-B_0}.
\]
\end{lemma}

\begin{proof}
The quantitative prime-polynomial estimate in \cref{lem:vk} gives, for every fixed $A_1$, a bound
\[
 |H(b+it)|\ll Q^{-A_1}
 \qquad(Q^{A_1+2}\le |t|\le 2X).
\]
Indeed the uniform error is
$\exp[-a(\log X)^{1/3}/(\log\log X)^{1/3}]$, smaller than every fixed
power of $Q^{-1}$, and the retained main integral is $O(1/|t|)$.
The difference kernel $(v^w-u^w)/w$ removes an artificial pole at
$w=0$ in the Perron argument.
Hence
$W\ll\sqrt H Q^{-A_1}$ and $V\ll\sqrt M(\log X)^C$.
Standard separation and dyadic amplitude classes cost fixed
logarithmic powers. It suffices to bound
$\min(\eta,T^{-1})^2V^2W^2R/X$.

First suppose $\delta\le h\le81/1000$ and take $k=\lfloor1/h\rfloor$.
Then $k\ge12$. The normalized exponents in the three non-diagonal
classical cases are
\[
 e_b=\frac{1-h}{2k}-\frac1{25},\quad
 e_c=-\frac{2h}{3}+\frac{24}{75k},\quad
 e_d=\frac{(2k-1)h-1-(2k-2)/25}{2k}.
\]
They satisfy
\[
 e_b\le-\frac1{650},\qquad e_c\le-\frac{8\delta}{25},
\]
and
\[
 e_d\le
 \begin{cases}-17/24000,&k=12,\\-6/4225,&k\ge13.\end{cases}
\]
The diagonal case retains the pointwise $Q^{-A_1}$ saving.

It remains to cover $81/1000\le h\le11/126$.
Here $M\ge X^{115/126}\ge T^{5/6}$ for $T\le X$.
Use $k=11$ and classify terms as follows.
If $T\le M$ or $V\ge T^{1/4}$, the classical bounds give
$R\ll MV^{-2}(\log X)^C$.
Otherwise \eqref{eq:modern} and the classical mean value bound give
\[
 R\ll X^\xi MT^{1/2}V^{-4},\qquad
 R\ll TV^{-2}(\log X)^C.
\]
Similarly, if $T\le H^k$ or $W\ge T^{1/(4k)}$, then
$R\ll H^kW^{-2k}(\log X)^C$.
Otherwise both
$R\ll TW^{-2k}(\log X)^C$ and
$R\ll H^kTW^{-6k}(\log X)^C$ hold.

If both factors are in their diagonal regimes, interpolation with
weights $1-1/(2k)$ and $1/(2k)$ gives
\[
 V^2W^2R
 \ll V^{1/k}WM^{1-1/(2k)}H^{1/2}(\log X)^C
 \ll XQ^{-A_1}(\log X)^C.
\]
In particular, this case does not incur an $X^\xi$ loss.
The other three cases give, with logarithmic factors suppressed,
\[
 X^{\xi/k}M^{1/k}T^{1-1/(2k)},\qquad
 MH^{1/3}T^{1/(3k)},\qquad
 X^{\xi/k}M^{1/k}HT^{1-3/(2k)}.
\]
For the first and third expressions, use the interpolation weights
$1-2/k,1/k,1/k$.
Since
$\min(\eta,T^{-1})^2T^a\le\eta^{2-a}$ for $0\le a\le2$,
their normalized exponents, before the arbitrarily small $\xi$ loss,
are
\[
 e'_b=\frac{1-2h-(2k-1)/25}{2k},\quad
 e'_c=-\frac{2h}{3}+\frac{24}{75k},\quad
 e'_d=\frac{(2k-2)h-1-(2k-3)/25}{2k}.
\]
Uniformly on the stated patch these obey
\begin{equation}\label{eq:pairgaps}
 e'_b\le-\frac1{11000},\qquad
 e'_c\le-\frac{137}{5500},\qquad
 e'_d\le-\frac1{1575}.
\end{equation}
Take $\xi<1/44000$. All non-diagonal cases retain a fixed power
saving. Because $Q=X^{o(1)}$, these powers dominate $Q^{-B_0}$.
Finally choose $A_1$ large enough for the diagonal case and its
fixed logarithmic losses.
\end{proof}

\begin{remark}
The new input repairs a whole interval of factor lengths, not just
one exceptional amplitude tuple.
The classical $k=12$ upper range and the modern $k=11$ lower range
overlap once
\[
 \frac{1+22\theta}{23}>
       \frac{1-21\theta}{2},
 \quad\text{equivalently}\quad \theta>\frac{21}{527}.
\]
This indicates further room in this particular lemma.
For the full sieve we only need the explicitly proved $1/25$ version.
\end{remark}

\section{Exact trilinear certificates}

Let $M=X^m$, $H=X^h$, $K=X^{1-m-h}$ and
$T=X^{1-\theta}$. The additional regions needed from the
Harman-type input are
\begin{align*}
 \mathcal H_1(\theta):\;&
 2\theta<h\le m<\tfrac12,\quad
 2h-2\theta<m<4\theta;\\
 \mathcal H_2(\theta):\;&
 2\theta<h<4\theta,\quad
 |4m+3h-2|<2\theta,\quad h\le m<\tfrac12;\\
 \mathcal H_3(\theta):\;&
 2\theta<h<4\theta,\quad
 |4m+h-2|<2\theta,\quad h\le m<\tfrac12.
\end{align*}
In $\mathcal H_1$ the factor $M$ must supply the prime-polynomial
pointwise saving; in the other two regions it is $H$.
We use compact interiors of these regions, so equalities along
geometric boundaries may be omitted without changing the integrals.

Write the normalized amplitude deficits as
\[
 |M(1/2+it)|\asymp X^{m/2-a},\quad
 |H(1/2+it)|\asymp X^{h/2-b},\quad
 |K(1/2+it)|\asymp X^{(1-m-h)/2-c},
\]
and let the large-value cardinality have exponent $f$.
For a monomial $M^rH^sK^u$, put
\[
 L=rm+sh+u(1-m-h),\qquad d=ra+sb+uc.
\]
Classical estimates give
\begin{equation}\label{eq:monomials}
 f\le\max(2d,2d+1-\theta-L),\qquad
 f\le\max(2d,6d+1-\theta-2L),
\end{equation}
up to fixed logarithmic factors.
These bounds include mixed products; individual powers alone need
not give the desired conclusion.

\begin{lemma}[rationally certified regions]\label{lem:harman}
Uniformly for $\theta\in[1/25,1/24]$, the inequalities
\eqref{eq:monomials} for the finitely many monomials in the attached
certificates imply
\begin{equation}\label{eq:certgain}
                 f-2(a+b+c)\le-\frac z{100}.
\end{equation}
Here $z$ is the minimum of the saving factor's amplitude deficit
and the positive slacks in the defining region inequalities.
Consequently, on every compact interior of any $\mathcal H_i$,
the corresponding trilinear mean square has arbitrary fixed
$Q^{-1}$ accuracy whenever the indicated prime factor does.
\end{lemma}

\begin{proof}
Each maximum in \eqref{eq:monomials} is a disjunction of two affine
inequalities. The proof certificates use rational coefficients in
the eight variables
\[
                  (m,h,a,b,c,f,z,\theta).
\]
They cover all disjunctions by a finite binary tree.
At each leaf, nonnegative rational multiples of the allowed
inequalities sum exactly to
$f-2a-2b-2c+z/100\le0$.
The two certificates have the following sizes:
\[
\begin{array}{c|r|r|c}
 \text{region}&\text{leaves}&\text{maximum depth}
      &\text{maximum sum of dual weights}\\ \hline
 \mathcal H_1&70&17&1521/400\\
 \mathcal H_2&36&9&573/200
\end{array}
\]
The verifier uses only Python's integer arithmetic, \texttt{Fraction},
and JSON parsing.
It checks the provenance of every inequality, the exact extremal
lengths on the region vertices, complete tree coverage, nonnegativity
of every dual coefficient, exact equality of coefficient vectors,
and nonpositivity of the resulting constant.
It does not import the numerical optimizer that found the trees.

For $\mathcal H_3$, replace $M$ by $K$:
\[
 |4(1-m-h)+3h-2|=|4m+h-2|.
\]
The required order $h\le1-m-h<1/2$ holds in this range
($\theta\le1/15$ suffices).

To interpret the certificate analytically, first discard classes
where a normalized amplitude is less than $1$:
the ordinary second moment for the product of the other factors
gives a fixed power saving.
On the other classes, all deficits lie between $0$ and half the
corresponding length exponent, apart from harmless logarithmic
normalizations.
On a fixed compact interior, let the region slack be at least $\rho$.
If the saving factor has deficit less than $\rho$,
\eqref{eq:certgain} retains a fixed positive power of its pointwise
$Q^{-A}$ saving.
If its deficit is at least $\rho$, the certificate gives the fixed
power saving $X^{-\rho/100}$.
All logarithmic errors in the input inequalities are multiplied
by bounded sums of rational dual weights.
Taking $A$ large covers these errors and the finitely many amplitude
classes. This proves the $Q$-accuracy conclusion.
\end{proof}

\section{A short-to-long comparison with a free accuracy}\label{sec:comparison}

The extra trilinear input is compared to a longer interval.
We give the comparison explicitly, in particular retaining the
cancellation of the low-frequency kernels.
For a coefficient sequence of fixed divisor bound, let
\[
 F(s)=\sum_{n\asymp X}a_n n^{-s},\qquad
 S_h(x)=\sum_{x-hx<n\le x}a_n,
\]
and put
\[
 \eta_1=e^{-3(\log X)^{1/3}},\qquad
 T_0=e^{(\log X)^{1/3}},\qquad T=X^\eps/\eta.
\]
Suppose that
\[
 \int_{T_0}^{T}|F(1/2+it)|^2\,dt\ll XQ^{-R}.
\]
Then, for every fixed $r$, increasing $R$ as necessary gives
\begin{equation}\label{eq:comparison}
 S_\eta(x)=\frac{\eta}{\eta_1}S_{\eta_1}(x)+O(\eta XQ^{-r})
\end{equation}
outside measure $O(XQ^{-r})$.

To see this, use the Mellin kernels
$k_h(s)=(1-(1-h)^s)/s$.
For $|t|\le T_0$,
\[
 \left|k_\eta(s)-\frac{\eta}{\eta_1}k_{\eta_1}(s)\right|
 \ll\eta\eta_1(1+|t|).
\]
The low-frequency contribution is therefore at most
\[
 \eta X\eta_1T_0^2(\log X)^C
      =\eta X e^{-(\log X)^{1/3}}(\log X)^C,
\]
smaller than every fixed power of $Q^{-1}$ relative to $\eta X$.
Mellin--Plancherel controls the middle frequencies by
$O(\eta^2X^3Q^{-R})$ in mean square over $x\in[X,2X]$.
For the tail, the classical bound
\[
 \int_U^{2U}|F(1/2+it)|^2dt\ll(U+X)(\log X)^C
\]
and kernel decay $O(U^{-1})$ give a mean square
\[
 \ll X^2(1/T+X/T^2)(\log X)^C,
\]
which has a fixed power saving relative to $\eta^2X^3$.
Chebyshev's inequality proves \eqref{eq:comparison}.
Thus no multiplication of a weak pointwise error by the exponentially
large cutoff $T_0$ is required.
The correct scaling is $\eta/\eta_1$; the occurrence of $\eta/\delta$
in the transcribed comparison lemma in Li II/III is not used.

\section{Stability of the signed finite sieve}\label{sec:stability}

Keep the initial cutoff
\[
                         a_0=\frac{11}{126}
\]
fixed. This is possible because \cref{lem:pair} retains that cutoff
throughout an interval of short-interval exponents below $1/24$.
The original decomposition begins with
\[
 S(A,(2X)^{1/2})=S_1-S_2+S_3,\qquad
 S_3=\sum_{(t_1,t_2)\in\Omega}S(A_{p_1p_2},p_2),
\]
where all summands of $S_3$ are nonnegative. The limiting exponent
domain, after the usual fixed-factor endpoint conventions, is
\[
 \Omega=\{a_0\le t_2\le t_1\le1/2,\quad t_1+2t_2\le1\}.
\]

\begin{lemma}[finite signed sieve stability]\label{lem:stability}
Suppose a finite Buchstab decomposition has a strictly positive
normalized lower-bound margin at $\theta_0$.
Assume its initial signed terms remain uniformly evaluable in a
neighborhood below $\theta_0$, and its other evaluations use finitely
many continuously moving affine exponent regions.
For negative children, assume that complete closed interval coverings
remain available, and that a reserved positive interval-exponent
slack makes the estimates valid on those closed coverings.
For freely selected good terms, compact interior estimates suffice.
Assume also that the universally quantified child ranges have no
parameter-only zero constraint or coincident moving boundary that
creates a positive-measure discontinuity at $\theta_0$.
If all discarded comparator densities are integrable, the same
scheme has positive margin for some $\theta<\theta_0$.
\end{lemma}

\begin{proof}
Remove thin boundary strips from the positive parent regions before
applying any further signed decomposition.
The removed parent sums are nonnegative, so discarding them preserves
the desired lower bound.
Within each retained parent region, repeat the exact Buchstab
identities, protecting all negative children by their complete
coverings and the reserved exponent slack. Do not erode individual
members of a closed interval covering and create gaps at shared endpoints.
For the terminal positive terms, use the eroded good regions and
discard the others.

The corresponding comparator sums satisfy the same identities.
Their total additional normalized deficiency is bounded by integrals
over the removed strips and the symmetric differences of the
terminal regions.
These tend to zero as the strip thickness tends to zero and the
parameter tends to $\theta_0$.
Integrability gives dominated convergence.
The hypothesis concerning child ranges ensures that removing
parent boundary strips is sufficient; an internal disappearing
coverage band of positive parent measure would violate it.
Choose the strip thickness to consume less than one quarter of the
original margin, then move $\theta$ by a positive amount small enough
to retain every strict analytic inequality and consume less than
another quarter.
The remaining margin is positive.
\end{proof}

We apply this lemma to the actual finite decomposition, rather than
postulating continuity of the published theorem.
Here are the checks specific to Li III.

\paragraph{Initial signed terms.}
The initial $S_1$ and $S_2$ use the Type I sieve input, with
$t_1\le1/2+o(1)$, and the small-prime terminal evaluation.
The Type I inequalities have strict room at these lengths;
for instance $2t_1\le1<1+\theta$.
The terminal prime factors below $X^{a_0}$ are covered uniformly by
\cref{lem:pair}. This is precisely where the old integer-moment
argument would fail. The auxiliary small sieve exponent (written
$10^{-1000}$ in the source) may be chosen sufficiently small after
the retained region margins have been fixed. Its value and all
sieve depths then remain fixed as $X$ grows.

\paragraph{The remaining analytic regions.}
Li III, Lemma 3.2, uses seven fixed moment patterns. We verified all
$7\cdot8=56$ diagonal/off-diagonal branches by exact nonnegative
rational H\"older identities; the patterns and complete moving
regions are recorded in \cref{app:regions}.
Two faces must move with the auxiliary exponent:
\[
 \mathcal C_5:\quad m+h\ge1-12\theta,
 \qquad
 \mathcal C_7:\quad m\le12\theta.
\]
Keeping the printed halves fixed would give an incorrect perturbation.
All non-pointwise monomials contain a strictly positive power of $T$.
Thus closed length faces at an auxiliary exponent $\theta_a$ acquire
a fixed power saving when the actual interval exponent is
$\theta_a+e$, $e>0$. Diagonal cases retain a positive fixed power of
the prime-polynomial saving.
The Watt/Type I faces are
\[
 2m+h\le1+\theta,\quad 4m+6h\le3+\theta,\quad
 4h\le1+3\theta,
\]
or $m\le(1+\theta)/2$, $h\le(1-\theta)/8$.
They likewise allow a fixed positive height slack; their harmless
$T^{\eps}$ loss is chosen smaller than that slack.
The additional trilinear regions are supplied independently by
\cref{lem:harman}.

\paragraph{The actual negative-child ranges.}
For $U_{0201},\ldots,U_{0209}$, put $s=t_1+t_2$ and $h=t_3$.
The parent conditions give
\[
 \frac{139}{336}\le s\le\frac{503}{1008},\qquad h\ge a_0.
\]
Consequently
\begin{equation}\label{eq:halfmargin}
 s+h-\frac12\ge\frac1{1008},\qquad
 \frac12-s\ge\frac1{1008}.
\end{equation}
Both moving half-faces therefore remain valid if
$\theta_a=1/24-d$ and $0<d<1/12096$.
$U_{0210},\ldots,U_{0212}$ use only $\mathcal C_1,\mathcal C_2,
\mathcal C_3$, with $m=t_1$, $h=t_2+t_3$; the two corrected faces
do not occur there.
For $U_{0303}$, $m=s\le11/24$ gives a $1/24$ margin from the
upper half-face, and $s+h\ge115/264+11/126>1/2$.

The full $U_{02}$ coverage uses the union of the classical regions
and the Harman strip regions; it is not a classical-only covering.
To use the closed strip endpoints rigorously, reserve $0<e<d$ and put
\[
 \theta'=\theta_a+e/2<1/24,\qquad
 \theta_{\rm actual}=\theta_a+e<1/24.
\]
The region $\mathcal H_1$ is impossible here because $m=s>4\theta'$.
For $\mathcal H_2$ or $\mathcal H_3$ closed at $\theta_a$, the same
point is in a compact interior at $\theta'$: its upper $h$ margin
is at least $2e$, its two strip margins are at least $e$, and
\[
 h-2\theta'\ge a_0-\frac1{12}=\frac1{252}.
\]
For $\mathcal H_2$, $m-h\ge83/336$ and $1/2-m\ge1/1008$.
For the $\mathcal H_3$ swap $m'=1-s-h$, one has
$m'-h\ge169/1008$ and $1/2-m'\ge1/1008$.
The truncation height $X^{e/4}/\eta\le X^{1-\theta'}$ for large $X$ fits
\cref{lem:harman} and leaves a fixed power saving in the transfer
tail. Thus shared negative-child endpoints are retained without gaps.

The other negative children use the following uniform Type I
groupings. In $U_{0301}$ use $(m,h)=(t_1+t_2,t_3)$, and in
$U_{0302}$ use $(m,h)=(t_2+t_3,t_1)$; the worst endpoint is
$t_3=t_2$ in both cases. In $U_{0401},U_{0402}$ the reversed
grouping is $(m,h)=(1-t_1-t_2,t_2)$, independent of $t_3$.
The defining parent conditions are precisely sufficient for these
full ranges. The combination of $\mathcal C_7$ and Type I controls
$U_{0303}$.

\paragraph{Finite boundary geometry.}
The primitive classical forms have integer left-hand coefficients
of absolute value at most $82$ and right-hand $\theta$-slopes at
most $42$. Substituting the finitely many groupings of
$t_1,t_2,t_3,1-t_1-t_2-t_3$ gives affine child endpoints.
Their order can change only on finitely many affine parent faces.
An exact rational check of the $3209$ substituted forms finds no
parameter-only zero face; coincident moving faces at $1/24$ are
the half-face families just handled by \eqref{eq:halfmargin} and
the $U_{0303}$ margins.
Every other nonidentity endpoint comparison has nonzero parent
dependence and can be protected by removing a parent boundary strip.
The next distinct parallel-face contact below $1/24$ is at $1/25$.
These checks supplement the complete analytic case proofs and the
actual negative-child analysis; a collision scan alone would not
prove sieve stability.

\paragraph{Signed pruning and numerical margin.}
Whole positive parents may be pruned before their signed refinements.
The relevant exact identities include
\begin{align*}
 S_{303}&=S_{3031}-S_{3032}-S_{3033}+S_{3034}+S_{3035},\\
 S_{304}&=S_{3041}-S_{3042}-S_{30431}+S_{30432}+S_{30433}.
\end{align*}
Shrinking the ``good'' predicate enlarges the negative terms
$S_{3033}$ and $S_{30431}$, so their full-range Type I evaluations
are essential. Only the positive terminal failures $S_{3035}$ and
$S_{30433}$ are discarded. For the reverse-Buchstab recovery, retain
the old bad-parent domain and compact subsets of its new good triples;
this reduces a positive saving by a controlled comparator integral.
No short-interval upper bound on a removed boundary strip is asserted.
The cost is charged solely to the same identity on the long comparator.

After reversal, the $\beta$ coefficient is the exact indicator of
$p_3$-rough integers, rather than an arbitrary divisor-bounded sequence.
Its unique ordered factorization has a bounded number of primes,
all at least $X^{a_0}$. Condition on all but one of those primes and
apply the uniform prime-polynomial bound on its remaining interval.
The outer reciprocal-prime sums are bounded. This supplies the
required quantitative prime-factored bound even when a group consists
only of $\beta$; see \cref{app:beta}.

All prime exponent variables in the four integrals are at least
$a_0$. The respective densities are bounded by
$a_0^{-3},a_0^{-5},a_0^{-5},a_0^{-4}$, hence by
\[
 a_0^{-5}=\frac{31757969376}{161051}<200000.
\]
Their boundary losses tend to zero by dominated convergence.
Using the four recorded numerical inequalities
\[
  \text{loss}_1<0.02,\quad \text{loss}_2<0.03,\quad
  \text{loss}_3<0.97,\quad \text{recovery}>0.04
\]
leaves positive room to perform this pruning.

The fully enumerated budget in \cref{app:explicit} permits
\[
 \rho=10^{-20},\qquad d=2\cdot10^{-23},\qquad e=10^{-23}.
\]
Here $\rho$ is the parent/terminal erosion width, the analytic
length regions use $\theta_a=1/24-d$, and the actual interval
exponent is $\theta_a+e$.
The total comparator cost is at most $0.00295856<0.005$.
Also $84d<\rho/4$ and $d<\min(1/600,1/12096)$, so every retained
endpoint-order sign and all fixed half-face gaps persist.
There is a remaining normalized margin greater than $0.017$ before
the fixed deterministic sieve errors, which are chosen smaller
than that margin. Thus the complete lower-bound sieve works at
\[
                 \theta_*=\frac1{24}-10^{-23}.
\]
All parameters, region erosions, sieve levels, and moment orders
are fixed before the requested exceptional accuracy is chosen.
The arguments of \cref{lem:pair,lem:harman} and
\eqref{eq:comparison} supply arbitrary fixed powers of
$Q^{-1}$.
The remaining Perron and Chebyshev steps lose only fixed powers of
$Q$ or $\log X$.
Thus the exceptional measure is
$X\exp[-D(\log\log X)^2]$ for any prescribed fixed $D$.
For the fixed-length formulation in \cref{thm:main}, take
$\eta=\tfrac12 X^{\theta_*-1}$. Then
$X^{\theta_*}/2\le\eta x\le X^{\theta_*}$ for $x\in[X,2X]$.
The prime count in the shorter multiplicative window gives the
claimed fixed-window lower bound after decreasing $\kappa$.
All moment and truncation arguments allow these fixed factors.
This proves the analytic assertion of \cref{thm:main}.

\begin{remark}[source convention and independent numerical check]
Our numerical hypotheses are stated for the literal printed regions,
including the empty $V_0$ with $t_1<1/16$.
Correcting this cap to $1/6$, as suggested by Lemma 5.2, is also safe
under those same hypotheses. Indeed the added region has
$t_1\le1/6$, $t_2\le1/8$ and lies inside the raw $U_{0301}$ region:
\[
 2t_1+3t_2\le17/24<25/24,\quad
 4t_1+10t_2\le23/12<73/24,\quad 4t_2\le1/2<9/8.
\]
Thus the union $U_0$ and its loss integral are unchanged.
Moving these positive parents into $V$ and enlarging the good predicate
can only reduce the first two losses and increase the recovery.
This is a monotonicity argument, not an assumption about which
convention the author used in an unpublished calculation.

An independent exploratory quadrature gives
\[
 I_0\approx0.966495918091412
\]
for the third loss under either convention. Gauss orders $12,24,48$
agree at the displayed precision; the source's conservative Buchstab
upper envelope gives approximately $0.967062500408101$.
These computations are not interval-certified enclosures and do not
replace the explicit external numerical hypotheses. The other three
integrals have not been independently recomputed in this note.
\end{remark}


\begin{proof}[Completion of \cref{thm:main}]
The signed sieve argument proves \eqref{eq:main:prime} under
\cref{hyp:liIII}. For $c_->c_*$ one has
$\theta_*<1-1/c_-$. The transfer theorem of
\cref{hb:sec:transfer} therefore applies to every sequence in the
stated range. The strict nesting argument gives exact floors, and
the measure/category proposition supplies the last two conclusions.
Finally, \eqref{eq:intro:exactgap} proves admissibility of
$c_{\mathrm{new}}$; $24/23>c_*$ proves the other stated endpoint.
\end{proof}

\section{What remains to reach \texorpdfstring{$25/24$}{25/24}}

The new two-factor estimate and the independently certified
trilinear regions both work at $\theta=1/25$.
This does not imply that the full signed sieve has positive margin
there.
To obtain the full interval exponent $1/25$ and then the range
$c>25/24$, one must quantitatively optimize the remaining admissible
regions and certify their final Buchstab deficiency below $1$.

The present conclusion is strictly stronger than importing the
statement at $1/24+\eps$: it supplies an actual positive movement
below $1/24$ and includes $c=24/23-10^{-23}$.
The explicit decrement is a conservative consequence of the checked
finite-geometry budget, rather than a numerical optimization of the sieve.
The embedded verification materials contain the exact rational
certificates and an independent verifier.
The single \LaTeX\ source embeds those materials so that the
mathematical computation can be reproduced without a floating-point
optimizer.

\appendix
\section{The seven classical moving regions}\label{app:regions}

Put $\tau=1-\theta$, $k=1-m-h$. For a separated value class,
$V,W,U$ denote the normalized amplitudes of $M,H,K$.
For a product polynomial $P=M^rH^sK^u$, let its length be
$L=M^rH^sK^u$ and its amplitude be $A=V^rW^sU^u$.
The classical bounds split exhaustively into a diagonal row
$R\ll LA^{-2}$ or two simultaneous off-diagonal rows
\[
                   R\ll TA^{-2},\qquad R\ll LTA^{-6},
\]
with fixed logarithmic factors understood. The following three
product choices give eight branches each.
\begin{center}
\begin{tabular}{@{}clc@{}}
\toprule
Region & Product polynomials & Pointwise saving factor\\
\midrule
$\mathcal C_1$ & $M^3,H^5,K^2$ & $|M(b+it)|^{1/5}$\\
$\mathcal C_2$ & $M^3,H^4,K^2$ & $|M(b+it)|^{1/2}$\\
$\mathcal C_3$ & $M^2,H^5,K^3$ & $|H(b+it)|^{1/3}$\\
$\mathcal C_4$ & $M^2,H^4,K^3$ & $|H(b+it)|^{2/3}$\\
$\mathcal C_5$ & $M^2H,H^6,K^2$ & $|H(b+it)|$\\
$\mathcal C_6$ & $M^2H,H^7,K^2$ & $|H(b+it)|$\\
$\mathcal C_7$ & $M^2,H^6,K^2H$ & $|H(b+it)|$\\
\bottomrule
\end{tabular}
\end{center}

The embedded \texttt{verify\_threefactor\_fraction.py} contains all
56 target monomials and their rational H\"older weights.
It reconstructs the valid rows, checks nonnegative weights summing
to one, verifies every exponent exactly, and checks that every
$T$-power is in $[0,1]$. The pointwise cases are $X$ times the
positive powers in the table; all other branches have a positive
$T$-power and give precisely the length conditions below.

For example, branch OOD of $\mathcal C_5$ uses weights
$11/24,1/72,1/36,1/2$ on the first product's mean row, its
Hal\'asz row, the second product's Hal\'asz row, and the third
product's diagonal row. Their sum is one and they give
\[
 V^2W^2U^2R\ll T^{1/2}M^{1/36}H^{13/72}K(\log X)^C.
\]
Putting $T=X^{1-\theta}$ yields $70m+59h\ge36\tau$.

{\interlinepenalty=10000
With $m,h,k>0$, the following closed inequalities suffice at the
auxiliary exponent. The actual exponent has the reserved positive slack.\par}
\begin{align*}
\mathcal C_1:\quad
 &m+h\le(15+14\theta)/29,\qquad h\ge2\tau/11,\\
 &29m-h\le9+20\theta,\qquad m\ge6\tau/19,\\
 &-m+29h\le5+24\theta,\qquad12m+11h\ge6\tau.\\[4pt]
\mathcal C_2:\quad
 &m+h\le(6+5\theta)/11,\qquad h\ge\tau/5,\\
 &6m+h\le2+4\theta,\qquad2\tau/7\le m\le8\theta,\\
 &m+8h\le2+6\theta,\qquad6m+5h\ge3\tau.\\[4pt]
\mathcal C_3:\quad
 &m+h\le(13+6\theta)/19,\qquad h\ge2\tau/11,\\
 &12m+h\le6+6\theta,\qquad m\ge14\tau/29,\\
 &-m+19h\le3+16\theta,\qquad12m+11h\ge8\tau.\\[4pt]
\mathcal C_4:\quad
 &1-8\theta\le m+h\le(5+2\theta)/7,\qquad h\ge\tau/5,\\
 &6m+h\le3+3\theta,\qquad m\ge5\tau/11,\\
 &-m+7h\le1+6\theta,\qquad6m+5h\ge4\tau.\\[4pt]
\mathcal C_5:\quad
 &1-12\theta\le m+h\le(6+5\theta)/11,\\
 &35m+23h\le17+18\theta,\qquad2m+h\ge10\tau/11,\\
 &2m+13h\le2+22\theta,\qquad70m+59h\ge36\tau.\\[4pt]
\mathcal C_6:\quad
 &m+h\le(7+6\theta)/13,\qquad41m+27h\le20+21\theta,\\
 &2m+h\ge12\tau/13,\qquad2m+15h\le2+26\theta,\\
 &82m+69h\ge42\tau.\\[4pt]
\mathcal C_7:\quad
 &2m+h\le(12+10\theta)/11,\qquad70m+11h\le34+36\theta,\\
 &5\tau/11\le m\le12\theta,\qquad m+6h\le1+11\theta,\\
 &35m+12h\ge18\tau.
\end{align*}
Three all-off-diagonal conditions are redundant on these regions.
For $\mathcal C_1$, the slack in $m\le20\theta$ is at least
$(6168\theta-107)/319\ge3493/7975$.
For $\mathcal C_3$, the condition $m+h\ge1-20\theta$ has the same
lower bound on its slack. For $\mathcal C_6$, the condition
$m+h\ge1-14\theta$ has slack at least
$(176\theta-7)/13\ge1/325$.
These assertions hold throughout $[1/25,1/24]$.

Every non-pointwise branch is a monomial $T^\alpha M^rH^sK^u$
with $0<\alpha\le1$. Reducing the critical height by $X^e$
gives a power saving $X^{-\alpha e}$ even on a closed length face.
For arbitrary $T$, the factor
$\min(\eta,T^{-1})^2T^\alpha$ is largest relative to $\eta^2$
at $T=1/\eta$. This proves the full weighted mean-square assertion,
not merely its value at one height.

\section{Direct verification of the moving Type I faces}\label{app:watt}

Let $MHL=X$ and $L(s)$ be the unweighted zeta-type polynomial.
Watt's estimate~\cite{Watt}, in the form stated in
\cite[Lemma 4.1]{LiIII}, is
\[
 \int_T^{2T}|L(1/2+it)|^4|N(1/2+it)|^2dt
 \ll\bigl(T+N^2T^{1/2}+NL^2T^{-2}\bigr)X^\xi
\]
for arbitrarily small fixed $\xi>0$.
For the first Type I region, Cauchy's inequality with the second
moment of $M^2H$ and Watt applied to $N=H$ gives
\[
 \int_T^{2T}|MHL|^2dt
 \ll(M^2H+T)^{1/2}
        (T+H^2T^{1/2}+HL^2T^{-2})^{1/2}X^\xi.
\]
The first four terms are
\[
 MH^{1/2}T^{1/2},\quad MH^{3/2}T^{1/4},\quad T,\quad HT^{3/4}.
\]
At $T=X^{1-\theta}$ these are at most $X$ under exactly the
three Type I inequalities displayed in \cref{sec:stability},
with the condition $\theta>0$ for the third term.
The last Watt term contributes at most $2X/\sqrt T$.
The length faces give $\log_X L\ge65/192$; since this input is
used only for $T\ge\sqrt L$, the last term has a uniform power saving.

For the second Type I region, use the second moment of $M^2$ and
Watt with $N=H^2$. The first four terms become
\[
 MT^{1/2},\quad MH^2T^{1/4},\quad T,\quad H^2T^{3/4}.
\]
They are at most $X$ if
$m\le(1+\theta)/2$ and $h\le(1-\theta)/8$.
The last Watt term is again at most $2X/\sqrt T$, and now
$\log_XL\ge23/64$.
The potentially saturated terms have $T$-powers at least $1/4$.
On closed faces at $\theta_a$, an actual interval exponent
$\theta_a+e$ therefore gains at least $X^{-e/4}$.
Choosing $\xi<e/8$ preserves a fixed power saving. This directly
derives the moving Type I faces; it does not interpolate between
two published endpoint statements.

\section{The rough factor after reversal}\label{app:beta}

Let $B\ge X^{a_0}$, $z\ge X^{a_0}$, and
\[
 F_{B,z}(s)=\sum_{\substack{B<n\le2B\\p\mid n\Rightarrow p\ge z}}n^{-s},
 \qquad B\le O(X).
\]
For $1/2\le\sigma\le1+1/\log X$ and $|t|\le2X$, the uniform
prime-polynomial estimate gives
\[
 |F_{B,z}(\sigma+it)|\ll_{a_0}
 \frac{B^{1-\sigma}}{\log X}
 \left((1+|t|)^{-1}+
 \exp\!\left[-c\frac{(\log X)^{1/3}}{(\log\log X)^{1/3}}\right]\right).
\]
To prove this, write each rough integer uniquely as an ordered
product $n=p_1\cdots p_r$, $z\le p_1\le\cdots\le p_r$.
The number $r$ is bounded in terms of $a_0$. Fixing
$p_2,\ldots,p_r$ and writing their product as $q$, the variable
$p_1$ lies in $(B/q,2B/q]\cap[z,p_2]$ (with the evident $r=1$
convention). Apply the uniform prime estimate on that interval.
Its contribution is bounded by $B^{1-\sigma}/q$ times the common
error factor and $1/\log X$. The sum of $1/q$ is bounded by a
fixed power of $\sum_{X^{a_0}\le p\le O(X)}1/p=O_{a_0}(1)$.
Repeated primes are included once by the ordered factorization.
Finite order and monomial restrictions leave bounded unions of
intervals after conditioning and are handled in the same way.

Thus for $|t|\ge Q^A$ this rough indicator has an arbitrary fixed
$Q^{-1}$ saving after increasing $A$.
The assertion relies on its exact rough-indicator coefficients.
It is not asserted for an arbitrary rough-supported coefficient sequence.
This distinction is needed in the reversed $S_{304}$ terms.

\section{An explicit decrement from the finite geometry}\label{app:explicit}

We now make the continuity argument quantitative. Throughout this
construction the printed $V_0$ remains empty. Its unused
$\mathcal H_1$ faces may be included in the inventory as an
overestimate; no new positive-volume good region is silently activated.

For the universally quantified negative children, assign each of
$t_1,t_2,t_3,1-t_1-t_2-t_3$ to $M,H$, or neither, with both groups
nonempty. There are $3^4-2\cdot2^4+1=50$ assignments.
The $64$ primitive faces and nine fixed geometry faces give
$64\cdot50+9=3209$ forms.
Eliminating the single child variable $t_3$ only requires comparisons
between its affine endpoints and parent-only face tests. Thus at most
$3209^2+3209$ events suffice.

The ordinary good predicates use four exponent tuples:
\[
 (t_1,t_2),\quad(t_1,t_2,t_3),\quad(t_1,t_2,t_3,t_4),\quad
 (1-t_1-t_2-t_3,t_2,t_3,t_4).
\]
Their respective nonempty group assignment counts are $2,12,50,50$,
so they contribute $64(2+12+50+50)=7296$ faces.
The $27$ original parent components have at most $12$ faces each
including the common domain conditions. There are fewer than $40$
additional fixed integration-domain faces. Boolean unions and
exclusions create no additional affine boundary.
Consequently the complete event count obeys
\[
 N\le3209^2+3209+7296+324+40
      =10\,308\,550<11\,000\,000.
\]
The exact inventory, after partial deduplication, has $496\,587$
events; the proof uses the larger bound $N=11\,000\,000$.

Write a primitive face as $A\cdot t\le b_0+b_1\theta$.
Its integer coefficients have magnitude at most $82$, and
$|b_1|\le42$. A child endpoint is
$(b_0+b_1\theta-A_1t_1-A_2t_2)/A_3$.
After comparing two endpoints, every nonconstant parent event has
a normal coordinate of magnitude at least
\[
                  \nu=\frac1{82^2}=\frac1{6724},
\]
because its denominator divides $A_3A'_3$.
Its $\theta$-slope is at most $84$.
All terminal and fixed source faces obey these same, or stronger,
bounds. The rational verifier checks these inequalities for every
listed endpoint comparison.

Parameter-only identities keep their truth values. The nearest
distinct lower parameter contact is $1/25$, at distance $1/600$.
Contacts exactly at $1/24$ are only the classified half-face families,
which are inactive on the actual child domains by the $1/1008$ and
$1/24$ margins. Thus no constant event opens a hidden full-fibre gap.

Take
\[
 \rho=10^{-20},\quad d=2\cdot10^{-23},\quad e=10^{-23},\quad
 \theta_a=1/24-d.
\]
Remove parent points in the nonconstant strips $|g(\theta_0)|\le\rho$.
Every face moves by at most
\[
 84d=1.68\cdot10^{-21}<\rho/4,
\]
so all retained endpoint-order signs remain fixed.
For ordinary terminal predicates, any change is supported in the
old strips $|g(\theta_0)|\le2\rho$.
In dimensions $2,3,4$, inside $[0,1/2]^r$, each such strip has volume
at most $4\rho/\nu$: integrate along a coordinate whose coefficient
has magnitude at least $\nu$ and bound the other side lengths by one.

Charge this union of strips separately to five comparator integrals:
the discarded positive-parent mass, the three losses, and the positive
recovery. Their densities are each at most $200000$.
The original parent classifications are disjoint, and repeated good
predicates are charged within each affected integral; there is no
additional Boolean-cell multiplicity. The total cost is at most
\begin{align*}
 5\cdot200000\cdot11\,000\,000\cdot4\rho\cdot6724
       &=0.00295856\\
       &<0.005<0.02.
\end{align*}
All of this is a comparator cost. Only whole positive parents and
positive terminal failures are discarded in the short sequence.

The small sieve parameters are selected independently of $e$.
For example a level exponent $10^{-1000}$, a small-prime exponent
$\delta_{\rm sieve}=10^{-2000}$, and an internal error parameter
$\eps_{\rm sieve}=\delta_{\rm sieve}^3=10^{-6000}$ preserve the usual
source ordering while remaining much smaller than $\rho$ and $e$.
The corresponding moment and sieve orders are enormous but fixed.
The actual interval exponent is still $\theta_a+e$; using a much
smaller internal error parameter does not reduce that interval slack.
The fixed fundamental-lemma accuracy is chosen within the remaining
positive budget.

We have therefore obtained
\[
                \Delta=d-e=10^{-23},\qquad
                \theta_*=1/24-10^{-23}.
\]
Finally set $c_{\rm new}=24/23-\Delta$.
Its strict position above the transfer boundary is exact:
\[
 c_{\rm new}-\frac{24}{23+24\Delta}
   =\frac{\Delta(47-552\Delta)}{23(23+24\Delta)}>0.
\]
The last inequality holds since $\Delta<47/552$.
This proves the explicit exponent claimed in \cref{thm:main} under
its stated numerical hypotheses. The calculation bounds a permissible
decrement; it does not identify the best one.

\section{Quantitative bookkeeping and reproducibility}\label{app:reproduce}

The arithmetic and geometric parameters are selected first.
For $Q=\exp[D(\log\log X)^2]$, one then uses arbitrarily large
\emph{fixed} powers $Q^{-R}$ in the mean-square inputs.
For every fixed $r,C,\sigma>0$,
\[
 (\log X)^C=Q^{o(1)},\qquad Q^r=X^{o(1)},\qquad
 X^{-\sigma}Q^r(\log X)^C\longrightarrow0.
\]
The Perron/mean-square transfer behind the sieve asymptotics has
the form
\[
 \int_X^{2X}|\mathrm{error}(x)|^2dx
 \ll X^3(\log X)^C
       \max_T\left[\min(\eta,T^{-1})^2
                       \int_T^{2T}|F(b+it)|^2dt\right].
\]
An input $\eta^2Q^{-R}$ and error threshold $\eta XQ^{-r}$ give
exceptional measure $O(X(\log X)^CQ^{-R+2r})$.
The low-frequency prime-polynomial main term is kept explicitly;
at cutoff $Q^K$, its uniform VK error is padded to $Q^{-3K}$.
The longer-interval comparison is proved separately in \cref{sec:comparison}.

Every factorization depth, divisor order, moment, and variable count
is fixed. Thus dyadic or finer rectangular decompositions cost only
fixed powers of $\log X$ or $Q$, which can be reserved in advance.
Coupled product and ordering cuts can be bracketed by smooth cuts
at relative width $Q^{-K}$, with fixed $K$ chosen for the desired
accuracy. Conditioning on all but one variable bounds the global
boundary count by $O(XQ^{-K}(\log X)^L+X^{1-\delta_0}(\log X)^L)$
for fixed $L$ and some fixed $\delta_0>0$.
Here is an explicit justification for divisor-bounded coefficients.
For fixed $r\ge1$, $y\asymp N$, and $0<v\le N$,
\[
 \sum_{y<n\le y+v}\tau_r(n)
 \ll_r v(\log 2N)^{r-1}
       +N^{1-1/r}(\log 2N)^{r-1}.
\]
Choose a largest factor in each ordered $r$-factorization.
The product $q$ of the other factors is $O_r(N^{1-1/r})$;
summing $v/q+1$ with weight $\tau_{r-1}(q)$ gives the bound
(the case $r=1$ is immediate).
If the selected variable has length $N\ge X^\delta$, the second
term, after the other variables are counted, is at most
$X^{1-\delta/r}(\log X)^L$.
Thus one may use $\delta_0=\delta/r$; it need not equal the
small-prime exponent itself.
Its short-window first moment and Markov's inequality give any
required $XQ^{-R}$ exceptional budget after increasing $K$.
The resulting Mellin separation has only fixed powers of $Q$ in
its norms and frequency truncations. Raise the low-frequency cutoff
above those frequency shifts; the uniform prime-polynomial estimate
then retains its saving. For the Harman comparison, the cutoff
$\exp[(\log X)^{1/3}]$ already dominates every fixed power of $Q$.
This fine analytic separation is distinct from the fixed geometric
parent strips, whose cost was charged only to comparator integrals.
Prime powers in a factor of length at least $X^\delta$ have a
power-saving first moment; Markov's inequality consequently also
puts their exception inside any fixed $XQ^{-R}$ budget.
Relative errors $O(1/\log X)$ in removing prime logarithmic weights,
and the fixed fundamental-lemma error, only consume a chosen small
portion of the positive main-term margin; they need not be
$Q$-small. No theorem with a constant depending on a fixed $B$
is used by replacing $B$ with a growing function of $X$.

The single source file contains standard \LaTeX\ \texttt{filecontents}
blocks. Compiling it writes the following verification files into
the compilation directory:
\begin{itemize}[itemsep=2pt]
\item \texttt{verify\_harman\_certificates.py} and its two JSON
certificates for $\mathcal H_1$ and $\mathcal H_2$;
\item \texttt{verify\_threefactor\_fraction.py};
\item \texttt{region\_stability\_probe.py};
\item \texttt{verify\_negative\_child\_coverage.py};
\item \texttt{verify\_quantitative\_sieve\_stability.py} and its
exact rational report;
\item the exploratory \texttt{li\_III\_i0\_convention\_check.py}
and its numerical results.
\end{itemize}
Run the exact checks with ordinary Python 3:
\begin{verbatim}
python3 verify_harman_certificates.py
python3 verify_threefactor_fraction.py
python3 region_stability_probe.py
python3 verify_negative_child_coverage.py
python3 verify_quantitative_sieve_stability.py
\end{verbatim}
They require only the standard library and use exact rational
arithmetic. Expected counts are $70+36=106$ Harman leaves,
$56$ classical H\"older identities, $3209$ substituted affine
forms, and $144$ open $U_{02}$ arrangement cells plus $2$
$U_{0303}$ cells, with their endpoints, in the original universally
quantified child coverings.
That coverage check verifies complete affine arrangements, not a sampling grid.
The final check verifies the full $496\,587$-event inventory,
the conservative count and normal/slope bounds, the explicit loss
budget, and the exact inequalities for $c_{\rm new}$.
The source also writes \texttt{mills\_chain\_5\_4.gp}, described
in \cref{app:chain}. The exploratory integral script additionally uses NumPy and
SciPy and is explicitly not a rigorous quadrature certificate.
None of these scripts is executed by \LaTeX\ itself.

The exact checks validate the stated finite inequalities and geometry.
They do not independently validate the external large-value theorem,
all four numerical sieve integrals, or an optimized value of $\Delta$.


\section{A finite illustration at exponent \texorpdfstring{$5/4$}{5/4}}\label{app:chain}

The original harmonic-branching manuscript recorded a finite greedy
chain starting at $p_1=1009$, using PARI/GP's \texttt{isprime} for
primality proofs~\cite{PARI}. Its first terms are
\[
1009,\ 5689,\ 49409,\ 736657,\ 21581503,\ 1470964273,\
288073066607,\ 211046580457399.
\]
The reported $22$ terms have digit counts
\[
4,4,5,6,8,10,12,15,18,23,28,35,44,55,69,86,107,134,167,209,261,326.
\]
This earlier computation is retained as a reproducible illustration;
it is not a numerical construction at $c_{\mathrm{new}}$ and is not
used in any infinite-existence proof.
For $c=5/4$ every nesting check is an exact integer comparison:
\[
 m\in J_{5/4}(p)
 \quad\Longleftrightarrow\quad
 m^4\ge p^5\quad\hbox{and}\quad(m+1)^4<(p+1)^5.
\]
The source embeds the original GP script as
\texttt{mills\_chain\_5\_4.gp}. The command
\begin{verbatim}
gp -q mills_chain_5_4.gp
\end{verbatim}
reconstructs the $22$ terms,
proves each selected term prime, checks every exact nesting inequality,
and writes the full chain. The optional decimal floor comparison at
the end is a numerical consistency check; the two integer inequalities
above are the nesting certificate.
The GP primality computation is distinct from the standard-library
rational verifiers in \cref{app:reproduce}.
No claim that this particular greedy chain continues indefinitely
follows from a finite calculation or from the global mass theorem.

\section*{Acknowledgment and revision record}
\addcontentsline{toc}{section}{Acknowledgment and revision record}
The initial harmonic-branching manuscript acknowledged research
development and manuscript preparation with Claude (Anthropic).
This revision consolidates that manuscript, the quantitative sieve
extension, and the explicit below-$1/24$ analysis into a single proof
sequence. It preserves the external-input status of the numerical
sieve margins and includes the exact verification materials used for
the new finite inequalities and geometry.


\begin{thebibliography}{99}\small\raggedright
\bibitem{Mil47}
W.~H. Mills, A prime-representing function,
\emph{Bull. Amer. Math. Soc.} \textbf{53} (1947), 604.
\url{https://doi.org/10.1090/S0002-9904-1947-08849-2}.

\bibitem{CGL}
B. Chen, V. Gupta, and Y.~C. Li,
\emph{Large Value Estimates for Dirichlet Polynomials with Characters
and Zero Density of Dirichlet $L$-Functions},
arXiv:2507.08296v2, 27 July 2026.
\url{https://arxiv.org/abs/2507.08296v2}.

\bibitem{LiIII}
R. Li, \emph{Primes in almost all short intervals III},
author-site manuscript, 13 pages, accessed 9 November 2026.
\url{https://runbolicarey.com/assets/downloads/Primes_in_almost_all_short_intervals_III.pdf}.

\bibitem{Harman}
G. Harman, \emph{Prime-Detecting Sieves},
London Mathematical Society Monographs, vol.~33,
Princeton University Press, Princeton, NJ, 2007.
\url{https://press.princeton.edu/books/hardcover/9780691124377/prime-detecting-sieves-lms-33}.

\bibitem{Dav00} H.~Davenport, \emph{Multiplicative Number Theory}, 3rd ed., revised by H.~L. Montgomery, Graduate Texts in Math. 74, Springer, 2000.

\bibitem{Ivi85} A.~Ivi\'c, \emph{The Riemann Zeta-Function}, Wiley, 1985.

\bibitem{Ford}
K. Ford, Vinogradov's integral and bounds for the Riemann zeta function,
\emph{Proc. London Math. Soc.} (3) \textbf{85} (2002), no.~3, 565--633.
\url{https://doi.org/10.1112/S0024611502013655}.

\bibitem{Ing40} A.~E. Ingham, On the estimation of $N(\sigma,T)$, \emph{Quart. J. Math. Oxford} 11 (1940), 201--202.

\bibitem{Hux72} M.~N. Huxley, On the difference between consecutive primes, \emph{Invent. Math.} 15 (1972), 164--170.

\bibitem{GM}
L. Guth and J. Maynard, New large value estimates for Dirichlet polynomials,
\emph{Ann. of Math.} (2) \textbf{203} (2026), no.~2, 623--675.
\url{https://doi.org/10.4007/annals.2026.203.2.6}.

\bibitem{SV77} B.~Saffari, R.~C. Vaughan, On the fractional parts of $x/n$ and related sequences II, \emph{Ann. Inst. Fourier} 27 (1977), 1--30.

\bibitem{Zac98} A.~Zaccagnini, Primes in almost all short intervals, \emph{Acta Arith.} 84 (1998), 225--244.

\bibitem{BP00} D.~Bazzanella, A.~Perelli, The exceptional set for the number of primes in short intervals, \emph{J. Number Theory} 80 (2000), 109--124.

\bibitem{GT25} A.~Gafni, T.~Tao, On the number of exceptional intervals to the prime number theorem in short intervals, preprint, arXiv:2505.24017 (2025).

\bibitem{LiI}
R. Li, \emph{Primes in almost all short intervals},
arXiv:2407.05651v6, 25 September 2025.
\url{https://arxiv.org/abs/2407.05651v6}.

\bibitem{Jia}
C. Jia, Almost all short intervals containing prime numbers,
\emph{Acta Arith.} \textbf{76} (1996), no.~1, 21--84.
\url{https://doi.org/10.4064/aa-76-1-21-84}.

\bibitem{MRSTT}
K. Matomäki, M. Radziwiłł, X. Shao, T. Tao, J. Teräväinen,
Higher uniformity of arithmetic functions in short intervals II.
Almost all intervals, \emph{Inventiones Mathematicae} 244 (2026),
967--1091.
\url{https://doi.org/10.1007/s00222-026-01408-6}.

\bibitem{LiII}
R. Li, \emph{Primes in almost all short intervals II},
Cambridge Open Engage working paper, version~1, 17 October 2025.
\url{https://doi.org/10.33774/coe-2025-jrnjl}.

\bibitem{MPPTW}
K. Matom\"aki, M. Pandey, J. Pliego, J. Ter\"av\"ainen, and M. Wang,
\emph{Higher order uniformity of the primes and cancellation of the
M\"obius function in shorter intervals},
arXiv:2610.09567v1, 7 October 2026.
\url{https://arxiv.org/abs/2610.09567v1}.

% Add direct citations in the prose:
% - Cite Watt at the displayed fourth-moment theorem and Appendix app:watt.
% - Cite Ford (and the retained standard zero-free-region reference if any)
%   when the Vinogradov--Korobov estimate is first invoked.
% - Cite Harman for the classical mean-value/Halasz machinery and the
%   standard Buchstab envelope, independently of the new exact certificates.
% - Cite CGL Theorem 1.1 as the primary modern input, and MPPTW Lemma 2.3
%   specifically for coefficient normalization, not as its origin.

\bibitem{Watt}
N. Watt, Kloosterman sums and a mean value for Dirichlet polynomials,
\emph{J. Number Theory} \textbf{53} (1995), no.~1, 179--210.
\url{https://doi.org/10.1006/jnth.1995.1086}.

\bibitem{PARI} The PARI~Group, \emph{PARI/GP version 2.15.4}, Univ. Bordeaux, 2023, \url{https://pari.math.u-bordeaux.fr/}.
\end{thebibliography}
\end{document}
